\documentclass[11pt,oneside]{article}
\usepackage[T1]{fontenc}
\usepackage[utf8]{inputenc}
\usepackage{lmodern}
\usepackage[a4paper,margin=27mm,headheight=15pt]{geometry}
\usepackage{amsmath,amssymb,amsthm,mathtools,bm}
\usepackage{microtype,booktabs,longtable,tabularx,array,enumitem}
\newcolumntype{L}[1]{>{\raggedright\arraybackslash}p{#1}}
\newcolumntype{Y}{>{\raggedright\arraybackslash}X}
\usepackage[dvipsnames]{xcolor}
\usepackage{fancyhdr}
\usepackage{tikz}
\definecolor{ink}{HTML}{172D46}
\definecolor{accent}{HTML}{285F78}
\definecolor{muted}{HTML}{53616E}
\definecolor{TierE}{HTML}{1B5E20}
\definecolor{TierC}{HTML}{0D47A1}
\definecolor{TierA}{HTML}{8D4004}
\definecolor{TierI}{HTML}{4A4A4A}
\usepackage{hyperref}
\usepackage{xurl}
\usepackage{bookmark}
\usepackage[nameinlink,noabbrev]{cleveref}
\usepackage{aliascnt}
\hypersetup{colorlinks=true,linkcolor=accent,citecolor=accent,urlcolor=accent,
 pdftitle={Surreal Scales in Quantum Theory and Gauge Models},
 pdfsubject={Merged research report: rare quantum events, postselection cost,
  exact elimination, soft modes, holonomy discrimination, and quaternionic curvature},
 pdfkeywords={surreal numbers, Hahn series, effective Hamiltonians, postselection,
  quantum instruments, Schur complements, surquaternions, Wilson loops, Yang--Mills}}
\pagestyle{fancy}
\fancyhf{}
\fancyhead[L]{\small\color{muted}Surreal scales in quantum theory and gauge models}
\fancyhead[R]{\small\color{muted}Merged research report}
\fancyfoot[C]{\thepage}
\renewcommand{\headrulewidth}{0.3pt}
\setlength{\parindent}{1.2em}
\setlength{\parskip}{0.28em}
\setlength{\emergencystretch}{3em}
\widowpenalty=10000 \clubpenalty=10000
\allowdisplaybreaks[2]
\setcounter{tocdepth}{2}
\numberwithin{equation}{section}
% One shared counter; alias counters so that cleveref prints the correct
% environment name (proposition, lemma, example, ...) in cross-references.
\theoremstyle{plain}
\newtheorem{theorem}{Theorem}[section]
\newaliascnt{proposition}{theorem}
\newtheorem{proposition}[proposition]{Proposition}
\aliascntresetthe{proposition}
\newaliascnt{lemma}{theorem}
\newtheorem{lemma}[lemma]{Lemma}
\aliascntresetthe{lemma}
\newaliascnt{corollary}{theorem}
\newtheorem{corollary}[corollary]{Corollary}
\aliascntresetthe{corollary}
\theoremstyle{definition}
\newaliascnt{definition}{theorem}
\newtheorem{definition}[definition]{Definition}
\aliascntresetthe{definition}
\newaliascnt{example}{theorem}
\newtheorem{example}[example]{Example}
\aliascntresetthe{example}
\newaliascnt{assumption}{theorem}
\newtheorem{assumption}[assumption]{Assumption}
\aliascntresetthe{assumption}
\newaliascnt{convention}{theorem}
\newtheorem{convention}[convention]{Convention}
\aliascntresetthe{convention}
\theoremstyle{remark}
\newaliascnt{remark}{theorem}
\newtheorem{remark}[remark]{Remark}
\aliascntresetthe{remark}
\crefname{theorem}{Theorem}{Theorems}
\crefname{proposition}{Proposition}{Propositions}
\crefname{lemma}{Lemma}{Lemmas}
\crefname{corollary}{Corollary}{Corollaries}
\crefname{definition}{Definition}{Definitions}
\crefname{example}{Example}{Examples}
\crefname{assumption}{Assumption}{Assumptions}
\crefname{convention}{Convention}{Conventions}
\crefname{remark}{Remark}{Remarks}
\crefname{section}{Section}{Sections}
\crefname{subsection}{Section}{Sections}
\crefname{appendix}{Appendix}{Appendices}
\crefname{table}{Table}{Tables}
\crefname{figure}{Figure}{Figures}
\crefname{equation}{}{}
\creflabelformat{equation}{(#2#1#3)}
\newcommand{\R}{\mathbb R}
\newcommand{\C}{\mathbb C}
\newcommand{\Q}{\mathbb Q}
\newcommand{\Z}{\mathbb Z}
\newcommand{\N}{\mathbb N}
\newcommand{\No}{\mathbf{No}}
\newcommand{\HH}{\mathbb H}
\newcommand{\HF}{\mathbb H_F}
\newcommand{\OO}{\mathcal O}
\newcommand{\mm}{\mathfrak m}
\newcommand{\TD}{\mathcal D}
\newcommand{\qi}{\mathbf i}
\newcommand{\qj}{\mathbf j}
\newcommand{\qk}{\mathbf k}
\newcommand{\Gap}{\mathsf G}
\newcommand{\Heff}{H_{\mathrm{eff}}}
\newcommand{\Corr}{\mathcal C}
\newcommand{\Lop}{\mathcal L}
\newcommand{\Green}{\mathcal G}
\newcommand{\QH}{\mathcal Q_H}
\newcommand{\ob}{\operatorname{ob}}
\newcommand{\eff}{\mathrm{eff}}
\newcommand{\fin}{\mathrm{fin}}
\newcommand{\expfin}{\exp_{\mathrm{fin}}}
\newcommand{\SYM}{S_{\mathrm{YM}}}
\newcommand{\Oval}[1]{O_{v}(#1)}
\DeclareMathOperator{\st}{st}
\DeclareMathOperator{\tr}{tr}
\DeclareMathOperator{\Tr}{Tr}
\DeclareMathOperator{\Rea}{Re}
\DeclareMathOperator{\Ima}{Im}
\DeclareMathOperator{\diag}{diag}
\DeclareMathOperator{\rank}{rank}
\DeclareMathOperator{\Span}{span}
\DeclareMathOperator{\lc}{lc}
\DeclareMathOperator{\ini}{in}
\DeclareMathOperator{\id}{id}
\DeclareMathOperator{\supp}{supp}
\DeclareMathOperator{\SU}{SU}
\DeclareMathOperator{\Sp}{Sp}
\DeclareMathOperator{\Herm}{Herm}
\DeclareMathOperator{\ad}{ad}
\newcommand{\op}{\mathrm{op}}
\newcommand{\HS}{\mathrm{HS}}
\newcommand{\abs}[1]{\left|#1\right|}
\newcommand{\norm}[1]{\left\|#1\right\|}
\newcommand{\ket}[1]{\lvert#1\rangle}
\newcommand{\bra}[1]{\langle#1\rvert}
\newcommand{\proj}[1]{\ket{#1}\!\bra{#1}}
\newcommand{\repo}{\texttt{VladimirReshetnikov/Surreal}}
% ---- claim tiers, as in docs/physics/surreal-scalars-and-spacetime --------
\newcommand{\TE}{\textcolor{TierE}{\textsf{\bfseries[E]}}}
\newcommand{\TC}{\textcolor{TierC}{\textsf{\bfseries[C]}}}
\newcommand{\TA}{\textcolor{TierA}{\textsf{\bfseries[A]}}}
\newcommand{\TI}{\textcolor{TierI}{\textsf{\bfseries[I]}}}
\newcommand{\status}[1]{\par\smallskip\noindent{\small\textbf{Status.} #1}\par\smallskip}
\newcommand{\src}[1]{{\small\textsf{[#1]}}}
\newcommand{\merge}{{\small\textsf{[merge]}}}

\begin{document}
\hypersetup{pageanchor=false}
\begin{titlepage}
\color{ink}
\vspace*{10mm}
{\small\sffamily A MERGED RESEARCH REPORT\par}
\vspace{9mm}
{\Huge\bfseries Surreal Scales in\par Quantum Theory\par and Gauge Models\par}
\vspace{8mm}
{\Large Rare-event shadows, postselection cost, exact elimination,\par
soft modes, holonomy discrimination, and quaternionic curvature\par}
\vspace{11mm}
\noindent\rule{\textwidth}{0.7pt}
\vspace{4mm}
{\large Prepared for Vladimir Reshetnikov\par}
{\normalsize 23 September 2026; merged from four manuscripts written independently on that day\par}
\vspace{9mm}
\noindent\begin{minipage}{0.94\textwidth}
\textbf{Central question.} Can surreal, surcomplex, or surquaternion scalars do something useful in theoretical physics without simply renaming infinities?

\medskip
\textbf{Answer developed here.} They provide an exact language for competing scales, including scales invisible to every power of a chosen perturbation parameter. In finite quantum and gauge models this language supports precise theorems about rare conditioning and its cost, virtual transitions and soft modes, the resolution of a small non-Abelian holonomy by finite quantum protocols, and the curvature of gauge fields. It also exposes operational and compositional obstructions that scalar extension cannot remove.

\medskip
\textbf{Scope.} Set-sized Hahn workspaces (or real closed fields) embedded in the surreals; fixed finite-dimensional quantum systems and finite protocols; finite lattice loops; fixed ordinary bundles on closed ordinary four-manifolds. No continuum quantum field theory, new experimental prediction, or resolution of a gravitational singularity is asserted.
\end{minipage}
\vfill
{\small\color{muted}Unrefereed, AI-assisted mathematical research report. Proofs are supplied in the text. The four accompanying exact symbolic programs verify selected finite identities and examples, not the general theorems. No Lean verification is claimed.\par}
\end{titlepage}
\hypersetup{pageanchor=true}

\begin{abstract}
We investigate a conservative role for surreal, surcomplex, and surquaternion numbers in theoretical physics: not a replacement for dynamics or measurement theory, but an exact coefficient language for separated scales. We work in a set-sized real-closed Hahn field $F=\R((t^\Gamma))$, its complexification $K=F[i]$, and the quaternion algebra $\HF$, and, where only finite algebra is needed, in any real closed field $F\supseteq\R$ with its natural valuation. This report merges four independently written manuscripts; \cref{qgs:sec:provenance} states what each contributed.

For finite quantum instruments, a leading-Gram formula computes the ordinary conditional state of any nonzero finite Kraus branch, even when its probability is infinitesimal; identical complete shadows can nevertheless conceal any ordinary channel, and all power-order data in one parameter can fail to determine a conditioned state. The sharp success-weighted inequality $\max(p,q)\,\TD(\rho',\sigma')\le\TD(\rho,\sigma)$ prices every postselected amplification; a one-sided filter concentrates infinitesimal entanglement at probability $d\lambda_d$. For Hermitian Hahn matrices we prove an exact invariant-graph elimination of several unlimited sectors with independently separated scales, an exact support-controlled block diagonalization at an ordinary gap with its inverse-splitting-scale dynamics, a finite-depth spectral hierarchy, a multiscale Gibbs reduction, and a three-level example in which virtual transitions defeat a bias below every power of the coupling. For positive block matrices with a soft eliminated block we classify every possible shadow of the Schur complement: it is $S_0-\Xi$ with $0\le\Xi\le S_0$ and $\rank\Xi\le\dim\ker D_0$, and every such defect occurs; a gap-relative window separates this static reduction from the Feshbach pencil.

For unit surquaternions we prove the exact commutator identity $1-\Rea c=2\abs{u_1\times u_2}^2$. For the known gates $I$ and $\chi(c)$ the optimal $m$-query unambiguous identification probability, with finite ancillas and adaptive control, is $w_m=1-T_m(1-w)$, $w=1-\Rea c$; the attained postselection frontier has valuation $\max\{0,\kappa-\beta/2\}$, where $\kappa$ is the commutator valuation and $\beta$ the valuation of the success probability, in every divisible value group. A balanced controlled-loop interferometer heralds an exact $\SU(2)$ unitary at its dark port. Initial quaternionic curvature determines the valuation and leading coefficient of a positive Wilson action. On a fixed $\SU(2)$ bundle, positive-support Hahn deformations preserve the charge exactly, the excess Yang--Mills action has twice the valuation of the wrong-duality curvature, and under stated Hodge--Green hypotheses a support-controlled Hahn Kuranishi map exists, whose nonzero leading quadratic obstruction cannot be removed by smaller scales. Quaternionic instruments have complex simulations, a naive quaternionic composite is impossible, and the Bell and canonical-commutation bounds survive.

Established ingredients are distinguished from proposed formulations throughout, with the claim tiers of the collection's physics report. All applications are mathematical model applications, not claims of empirically new physics.
\end{abstract}

\noindent\textbf{Keywords:} surreal numbers; Hahn series; postselection; quantum instruments; effective Hamiltonians; Schur complements; unambiguous discrimination; surquaternions; Wilson loops; Yang--Mills deformations.

\tableofcontents
\newpage

\section{Purpose, contributions, and limits}\label{qgs:sec:intro}

\subsection{The useful question is not whether infinity becomes a number}
A formula can acquire a value in a larger field without acquiring a physical interpretation. An unlimited curvature scalar is still not a regular spacetime point; an infinitesimal Born weight is still not an ordinary real frequency; an algebraic matrix exponential is not automatically available at an arbitrary infinite imaginary argument. The collection's existing physics report makes these distinctions particularly clearly \cite{repoPhysics}. This report does not repeat its black-hole calculations or present its negative conclusions as new results.

Instead, we ask more specific questions. Can one extract mathematically invariant information from a quantum process whose leading ordinary shadow vanishes, and what does that extraction cost? Can several independently small mixing scales, or soft eliminated modes, be eliminated without replacing them by a single perturbation parameter? How much of an infinitesimal non-Abelian holonomy can a finite quantum protocol resolve? Can quaternionic noncommutativity retain useful scale information after the ordinary gauge field has become flat, and can a hierarchy of scales repair a gauge-theoretic obstruction? These questions admit precise answers in finite algebra, or on a fixed ordinary manifold with explicitly imported elliptic input, where no infinite-dimensional functional analysis or noncanonical global oscillatory exponential is needed.

The operative architecture is
\begin{equation}\label{qgs:eq:architecture}
\text{ordinary finite model}
\ \longrightarrow\ \text{valued coefficient model}
\ \longrightarrow\ \text{specified ordinary readout}.
\end{equation}
The first arrow introduces exact scale separation. The second is a modelling choice, usually standard part, sometimes standard part \emph{after} normalization or energy rescaling. The order of these operations matters. Equivalently, one compares
\begin{equation}\label{qgs:ops:eq:diagram}
 \text{perform the operation, then reduce}
 \qquad\text{versus}\qquad
 \text{reduce, then perform the operation}.
\end{equation}
For bounded polynomial operations the two procedures agree. They can disagree for conditioning, for inversion of a small stiffness or gap, and for evolution over an unlimited time. Most of the useful results below identify exactly when two orders of operation agree, what additional leading data recover the answer when they do not, which discrepancies are possible, and what physical resource is consumed.

\subsection{Main results and their status}\label{qgs:sec:ledger}
\Cref{qgs:tab:contributions} is a claim ledger. ``Proposed'' means that the stated combination and formulation were developed in one of the four source manuscripts and were not identified in their targeted repository and literature reviews. It is not a certification of priority. Each underlying ingredient is attributed where used, and the proofs remain subject to independent checking. The letters in brackets name the source manuscripts, described in \cref{qgs:sec:provenance}.

\begin{longtable}{@{}L{0.24\textwidth}L{0.41\textwidth}L{0.24\textwidth}@{}}
\caption{What is proved, and what is not claimed to be new.}\label{qgs:tab:contributions}\\
\toprule
Result & Content & Claim status\\
\midrule
\endfirsthead
\toprule
Result & Content & Claim status\\
\midrule
\endhead
\bottomrule
\endfoot
Rare-branch initial state, \cref{qgs:thm:rare} \src{03, 07, 08}
& Computes a conditional ordinary state from leading Kraus--purification products, including infinitesimal probabilities; universal rare embedding (\cref{qgs:ops:prop:rare-embedding}).
& Proposed arbitrary-rank formulation; the commutative case is \texttt{fsp:thm:skeleton}.\\[3pt]
Postselection cost, \cref{qgs:thm:cost} \src{04, 07, 08; weaker 03}
& $\max(p,q)\TD(\rho',\sigma')\le\TD(\rho,\sigma)$; equality can occur.
& Standard data-processing method, re-proved over $F$.\\[3pt]
Invisible power data, \cref{qgs:thm:alljets} \src{07}
& Identical data modulo every $\epsilon^N$ can give orthogonal conditional shadows.
& Proposed explicit obstruction.\\[3pt]
Entanglement filter, \cref{qgs:ops:thm:entanglement} \src{08}
& One-sided concentration probability $d\lambda_d$.
& Special case of Vidal's theorem.\\[3pt]
Multiscale elimination, \cref{qgs:thm:elimination} \src{07}
& Exact invariant graph, finite effective block, unrestricted ratios of small scales, positive unlimited complement.
& Proposed; Schur and Schrieffer--Wolff mechanisms are established.\\[3pt]
Support-controlled block diagonalization and dynamics, \cref{qgs:obs:thm:block,qgs:obs:thm:dynamics,qgs:obs:thm:hierarchy,qgs:obs:thm:gibbs} \src{03}
& Exact unitary block at an ordinary gap with support certificates; inverse-splitting-scale dynamics; finite-depth hierarchy; multiscale Gibbs reduction.
& Proposed support formulation of classical mechanisms.\\[3pt]
Virtual selection, \cref{qgs:prop:virtual} \src{07}
& A virtual $\epsilon^2$ splitting defeats a bias below every power of $\epsilon$.
& Worked counterexample, not a new mechanism.\\[3pt]
Soft-mode Schur defects, \cref{qgs:ops:thm:defect,qgs:ops:thm:classification,qgs:ops:thm:window} \src{08}
& Complete fixed-shadow classification with rank bound and realization; displacement obstruction; static energy window.
& Proposed candidate contribution.\\[3pt]
Quaternionic simulation and composition, \cref{qgs:sec:quaternions} \src{03, 04, 07}
& Complex simulation of finite instruments; obstruction to the naive $mn$-level composite (two proofs).
& Established structures, explicit scalar-extension consequences.\\[3pt]
Exact commutator identity, \cref{qgs:ops:thm:commutator} \src{04, 08; leading term 07}
& $1-\Rea c=2\abs{u_1\times u_2}^2$ for unit quaternions.
& Elementary; not a new quaternion theorem.\\[3pt]
Holonomy discrimination, \cref{qgs:hol:thm:query,qgs:hol:thm:frontier,qgs:hol:thm:valuation} \src{04}
& Optimal $m$-query overlap $T_m(s)$ with ancillas and adaptivity; $w_m=1-T_m(1-w)$; attained frontier and its valuation law.
& Proposed synthesis of Chefles--Barnett and Ac\'in.\\[3pt]
Dark-port unitary, \cref{qgs:ops:thm:dark} \src{08}
& Exact heralded $\SU(2)$ unitary with probability $\abs{u_1\times u_2}^2$.
& Proposed application of finite $\SU(2)$ algebra.\\[3pt]
Wilson initial curvature, \cref{qgs:thm:wilson} \src{04, 07}
& Positive face weights turn leading loop curvature into an exact minimum-valuation action formula.
& Proposed arbitrary-rank formulation of classical gauge algebra.\\[3pt]
Continuum gauge deformations, \cref{qgs:obs:thm:chern,qgs:obs:thm:action,qgs:obs:thm:kuranishi,qgs:obs:thm:obstruction,qgs:obs:thm:floor,qgs:obs:thm:torus-obstruction} \src{03}
& Exact charge rigidity; action-gap valuation; Hahn Kuranishi map; persistent obstruction and sharp action floor.
& Proposed arbitrary-support formulation; Kuranishi theory is classical.\\
\end{longtable}

\subsection{What would count as a physical application?}
\TA\ An application in this report means an exact mathematical statement about a specified Hamiltonian, instrument, quadratic form, lattice model, or connection on a fixed bundle. A stronger claim of new physics would also require a derivation of that model, a readout rule, an accessible experimental regime, and a difference from existing real or complex predictions. None follows merely from writing a coefficient as a surreal number.

All scale parameters below are dimensionless. A physical Hamiltonian is obtained by multiplying the displayed matrix by a fixed reference energy $E_\ast$; time may then be measured in $\hbar/E_\ast$, and in quantum dynamics formulas we set $\hbar=1$. Similarly, displayed Wilson and Yang--Mills actions, stiffnesses, and ratios such as $\abs z/\delta$ are dimensionless, dimensional quantities having been divided by stated reference units before a parameter is declared infinitesimal. An expression such as $t^{-\gamma}$ is not silently interpreted as the cardinality of an experiment. This convention prevents comparisons between dimensionally incompatible quantities disguised as valuation comparisons.

The most defensible benefits are exact asymptotic bookkeeping, certificates of scale dominance, detection of invalid truncations, and proofs that a proposed readout cannot reveal a hidden correction under stated resource constraints. These are useful mathematical tools even when every finite-scale prediction agrees with ordinary quantum mechanics.

\subsection{Claim tiers}\label{qgs:sec:tiers}
This report is filed beside the collection's physics report and adopts its claim-status firewall (\texttt{phys:conv:tiers} in \cite{repoPhysics}), because the characteristic failure of surreal-physics writing is a correct calculation presented at the wrong epistemic grade.

\begin{convention}[Claim tiers]\label{qgs:conv:tiers}
Numbered results and substantive paragraphs carry one of the following tags.
\begin{description}[style=nextline,leftmargin=0pt,labelsep=0pt]
\item[\TE\ Exact symbolic identity.] A finite algebraic identity or finite example checked exactly, never to a tolerance, by at least one of the four accompanying programs (\cref{qgs:app:verification}), or derived from such an identity by finite algebra. This is the entire content of what the programs establish.
\item[\TC\ Mathematical theorem.] A statement with an ordinary mathematical proof in the text, under hypotheses stated in full. The hypotheses are established for no physical model; a theorem about a finite matrix over a Hahn field is not a theorem about an experiment. Nothing in this tier is machine-checked.
\item[\TA\ Assessment.] An argument rather than a result: a modelling choice, a proposed application, a judgement of what a construction does or does not accomplish, or a negative claim scoped as a missing implication.
\end{description}
\TI\ marks a result imported from the literature or from another report, used with its own hypotheses and not reproved here; it is a provenance marker, not a fourth tier.
\end{convention}

\subsection{Four sources, one report}\label{qgs:sec:provenance}
This report was assembled from four manuscripts, written independently on 23 September 2026 and delivered together. They are numbered by their position in the delivery; the source manuscripts themselves are not shipped with the report, but their code and data are, under the file prefixes listed.

\begin{longtable}{@{}L{0.05\textwidth}L{0.27\textwidth}L{0.09\textwidth}L{0.49\textwidth}@{}}
\caption{The four source manuscripts.}\label{qgs:tab:sources}\\
\toprule
No. & Manuscript and file prefix & Pin & Contribution to this report\\
\midrule
\endfirsthead
\toprule
No. & Manuscript and file prefix & Pin & Contribution to this report\\
\midrule
\endhead
\bottomrule
\endfoot
07 & \emph{Surreal Scales in Quantum Theory and Gauge Models} (base); \nolinkurl{07-scales-in-physics-} & \texttt{c0a36ee} & The structure and title; \cref{qgs:sec:workspace,qgs:sec:quantum} in large part; the leading-Gram theorem in the form printed; the all-power-data obstruction; the sharp cost theorem in the form printed; the multiscale elimination theorem, its real-parameter counterpart and dressed dynamics; the scalar trichotomy and virtual selection; the evaluation-matrix proof of the composite obstruction; the Wilson action theorem and the two-scale plaquette; the Bell and canonical-commutation obstructions; the benchmarks.\\[3pt]
03 & \emph{Infinitesimal Scales, Observable Limits, and Gauge Topology}; \nolinkurl{03-observable-limits-} & \texttt{934810a} & The positive-support evaluation lemma and the two-reductions discussion; the adaptive operational reduction; the weaker postselection bound (now \cref{qgs:obs:cor:postselection}); a proof of the leading-Gram theorem and the conditional qubit; the support-controlled block diagonalization, finite-angle dynamics, three-level mediated coupling, spectral hierarchy and Gibbs reduction; \cref{qgs:sec:gauge} entire (charge rigidity, action gap, torus example, Kuranishi map, gauge slice, obstruction, action floor, torus obstruction); boundary discussions; the recursion appendix.\\[3pt]
04 & \emph{Infinitesimal Gauge Holonomy as a Quantum Resource}; \nolinkurl{04-gauge-holonomy-}. Its author line reads ``Mathematical development and exposition: ChatGPT''. & \texttt{934810a} & The general real-closed setting of the finite results; Helstrom, pure-state distance and overlap triangle; the normalized affine chart and exact Wilson defect; conjugacy and regular class observables; \cref{qgs:sec:discrimination} entire (optimal $m$-query overlap, no finite perfect test, query scaling, state separation, unambiguous bound, frontier, valuation law, examples, precision certificate, calibration cost, sensing and control applications, real specializations); a proof of the sharp cost theorem with a second equality family; the quaternionic Gram lemma, the simulation at full real-closed generality, and the bilinear-map proof of the composite obstruction; implementation formulas.\\[3pt]
08 & \emph{Surreal Scales in Quantum Operations, Soft-Mode Reduction, and Non-Abelian Interferometry}; \nolinkurl{08-quantum-operations-} & \texttt{c0a36ee} & The intrinsic form of the leading-Gram theorem, its mixed-output example and amplitude precision; the universal rare embedding; a proof of the sharp cost theorem with the residue threshold; finite-trial accounting; the entanglement filter; \cref{qgs:sec:schur} entire (defect formula, classification, displacement, Feshbach reduction, static window); the general commutator identity; the analytic invariant-germ obstruction; \cref{qgs:sec:dark} entire (dark-port unitary, precision and background thresholds); the Puiseux and rank-two realizations; the decision table and scalar-effect criterion.\\
\end{longtable}

Both pins are ancestors of the collection commit at which this report was placed, and none of the reports the sources inspected changed between those pins and that commit.

\paragraph{Why one report.} Every pair of sources shares at least one theorem: 03, 07 and 08 prove the leading-Gram theorem; 03, 04, 07 and 08 prove a postselection bound and give the same tagged-state equality example; 03, 04 and 07 prove the quaternionic simulation; 04 and 07 prove the composite obstruction; 04, 07 and 08 prove commutator identities for a small loop; 03, 07 and 08 treat effective Hamiltonians. The shared spine is printed once.

\paragraph{Where the merge had to choose.}
\begin{itemize}
\item \emph{Base.} Source 07, because on each duplicated theorem its version is at least as general and sharp as source 03's (the runner-up), and its title is the natural umbrella.
\item \emph{Postselection.} The sharp inequality $\max(p,q)\TD_{\rm out}\le\TD_{\rm in}$ (proved identically, by a failure flag, in 04, 07 and 08) is printed as \cref{qgs:thm:cost}. Source 03's bound $2\TD/p_*$ is true but weaker (a factor two, and $p_*\le\min(p,q)$ in place of $\max(p,q)$); it is kept as \cref{qgs:obs:cor:postselection} together with 03's own proof, which loses the factor by bounding $\abs{p-q}$ separately.
\item \emph{Leading-Gram theorem.} Printed once as \cref{qgs:thm:rare}, in 07's statement, with 08's intrinsic invariance and 03's Gram-representation invariance, and credited to all three. The commutative diagonal case is \texttt{fsp:thm:skeleton} of the finite-probability report \cite{repoProbability}.
\item \emph{Commutator identity.} Source 08's identity for arbitrary unit quaternions is printed as \cref{qgs:ops:thm:commutator}; 04's identity in the normalized affine chart is \cref{qgs:hol:cor:comm}; 07's leading-term statement in the exponential chart is \cref{qgs:prop:plaquette}, whose leading term already appears with a stronger remainder in \texttt{squat:eq:groupcomm} \cite{repoQuaternion}.
\item \emph{Quaternionic simulation.} Stated at 04's generality (any real closed $F\supseteq\R$); the composite obstruction keeps both proofs (07's evaluation matrix, 04's bilinear map) as marked routes.
\item \emph{Effective Hamiltonians.} 07's elimination (unlimited high block, independent scales), 03's block diagonalization (ordinary gap, infinitesimal couplings, support certificates) and 08's soft-mode and Feshbach results are different theorems and are all kept.
\item \emph{Notation.} Each source's symbols were renamed to one convention, listed in \cref{qgs:sec:notation}.
\end{itemize}
Results added in the merge carry the tag \merge\ and have complete proofs here.

\subsection{Notation, and the symbols renamed from the sources}\label{qgs:sec:notation}
\begin{convention}[Scalars]\label{qgs:conv:scalars}
$F$ is a real closed field and $K=F[i]$ its algebraic closure, as in the collection's physics report; when explicit, $F=\R((t^\Gamma))$ and $K=\C((t^\Gamma))$ for a set-sized divisible ordered abelian group $\Gamma$. The valuation is $v$, with $\lc$ the leading coefficient, and $\OO_F,\mm_F,\OO_K,\mm_K$ the finite elements and infinitesimals. The quaternion algebra over $F$ is $\HF$, with units $\qi,\qj,\qk$ in bold; the complex unit is $i$.
\end{convention}

\begin{convention}[Reserved and single-meaning symbols]\label{qgs:conv:symbols}
\begin{itemize}
\item $\epsilon$ is a positive infinitesimal, fixed locally. In the rank-two workspace $\Gamma=\Q^2$ (lexicographic) we always write $\epsilon=t^{(0,1)}$ and $\zeta=t^{(1,0)}$, so $\zeta<\epsilon^N$ for every ordinary $N$. The letter $\eta$ is \emph{not} used: the collection reserves it for the infinitesimal phase residue (see the physics report's phase conventions).
\item $\TD(\rho,\sigma)=\tfrac12\norm{\rho-\sigma}_1$ is trace distance; $\TD_m$ and $\TD_{m,\max}$ are trace distances in \cref{qgs:sec:discrimination}. An upright $D$, $D_0$ is always an eliminated or high-energy matrix block.
\item $\Delta=\diag(\delta_1,\ldots,\delta_s)$ is only the scale matrix of \cref{qgs:sec:elimination}. The ordinary gap matrix of \cref{qgs:sec:block} is $\Gap$; the Schur defect of \cref{qgs:sec:schur} is $\Xi$.
\item $\kappa=v(u_1\times u_2)=v(c-1)$ is only the valuation of a small commutator (\cref{qgs:sec:holonomy,qgs:sec:discrimination,qgs:sec:dark}).
\item $L_j$ are only leading Gram factors of a rare branch. The effective Hamiltonian is $\Heff$; the elliptic operator of \cref{qgs:sec:gauge} is $\Lop$.
\item $C=BD^{-1/2}$ is only the whitened coupling of \cref{qgs:sec:schur}; the correction matrix of \cref{qgs:thm:elimination} is $\Corr$.
\item $\chi$ is the complex representation of quaternions; the letter $J$ is not used for it. $W$ is the graph isometry of an invariant subspace.
\item Kraus operators are $M_j$, $M_{a,j}$; the letter $K$ is the complex field.
\item $S(H)=A-BD^{-1}B^*$ is the Schur complement of a block matrix $H=\bigl(\begin{smallmatrix}A&B\\B^*&D\end{smallmatrix}\bigr)$ with invertible $D$; $\mathcal S$ is a Wilson action and $\SYM$ the Yang--Mills action.
\end{itemize}
\end{convention}

\Cref{qgs:tab:renamed} discloses every symbol renamed from a source. No normalization was changed.

\begin{longtable}{@{}L{0.07\textwidth}L{0.40\textwidth}L{0.45\textwidth}@{}}
\caption{Symbols renamed from the sources.}\label{qgs:tab:renamed}\\
\toprule
Source & Source symbol & Symbol here\\
\midrule
\endfirsthead
\toprule
Source & Source symbol & Symbol here\\
\midrule
\endhead
\bottomrule
\endfoot
07 & real field $K$, complex field $F$, group $G$; $\Tr$ & $F$, $K$, $\Gamma$; $\tr$\\
07 & second scale $\eta=t^{(1,0)}$; Kraus $K_j$; $\exp_f$ & $\zeta$; $M_j$; $\expfin$\\
07 & effective Hamiltonian $L$, leading term $L_0$, correction $C$, $S=I+X^*X$ & $\Heff$, $S(H)$, $\Corr$, $T_X$\\
07 & trace distance $D$; face holonomy $Q_f$, face weight $\kappa_f$ & $\TD$; $q_f$, $\lambda_f$\\
07 & Jacobian $J$; norm bound $M$; generic matrices $M,N$ & $\partial_YP$; $b_0$; $Z_1,Z_2$\\
03 & $F_\Gamma$, $K_\Gamma$, valuation $\nu$; Kraus $K_{ba}$, branch $\mathcal T_b$ & $F$, $K$, $v$; $M_{ba}$, $\Phi_b$\\
03 & gap matrix $\Delta$, scalar energy $E$, perturbation blocks $A,B,C$ & $\Gap$, $e$, $P_{11},P_{12},P_{22}$\\
03 & graph isometries $J$, $J_\perp$; complementary unitary $W$ & $W$, $V$; $U_\perp$\\
03 & generator $K$ (Hermitian); three-level gap $\Delta$; second scale $\eta$ & $Z$; $d$; $\zeta$\\
03 & support set $S$, monoid $S^*$; action $S(A)$; Gibbs index set $J$ & $\Sigma$, $\Sigma^*$; $\SYM(A)$; $\mathcal J$\\
03 & elliptic operator $L$, Green operator $Q$, obstruction $\kappa$, $\kappa_2$ & $\Lop$, $\Green$, $\ob$, $\ob_2$\\
03 & quaternion units $i,j,k$ (italic) & $\qi,\qj,\qk$\\
04 & real closed $K$, complex $K(\mathrm i)$, valuation $\nu$, $\HH_K$ & $F$, $K$, $v$, $\HF$\\
04 & imaginary quaternions $u,v$; commutator $q$ & $u_1,u_2$; $c$\\
04 & Kraus $K_a$, filter $K_{\mathrm s}$; $D_m$, $D_{m,\max}$, $D_{\max}$ & $M_a$, $M_{\mathrm s}$; $\TD_m$, $\TD_{m,\max}$, $\TD_{\max}$\\
04 & plaquette $P$, holonomy $q_P$, weight $\lambda_P$, action $S$; $\chi_n$ & face $f$, $q_f$, $\lambda_f$, $\mathcal S$; $\chi$\\
04 & scales $x=t^{(0,1)}$, $y=t^{(1,0)}$; $\eta$ in Example 8.1 & $\epsilon$, $\zeta$; $\zeta$\\
08 & representation $J$; normalized trace $W$; axes $I,J$ & $\chi$; $\Rea c$; $n_1,n_2$\\
08 & quaternions $q=s+u$, $r=h+v$ & $q_1=s_1+u_1$, $q_2=s_2+u_2$\\
08 & inputs $\rho,\eta$; outputs $\sigma,\tau$ & $\rho,\sigma$; $\rho',\sigma'$\\
08 & Kraus $K_{a,j}$, $K_\pm$; ordinary Kraus $L_j$ and $S=\sum L_j^*L_j$ in Prop.~3.5 & $M_{a,j}$, $M_\pm$; $A_j$ and $E_\Psi$\\
08 & defect $\Delta$; realization factor $L$; $W=BU$; spectral ratio $\kappa$ & $\Xi$; $N$; $\widehat B$; $\xi$\\
08 & quadratic form $Q_H$; smallest eigenvalue $\delta$ of $D$ & $\QH$; $\delta_D$\\
08 & scales $\varepsilon=t^{(1,0)}$, $h=t^{(0,1)}$; densities $\sigma$, $\nu$ of its background model & $\zeta$, $\epsilon$; $\rho_{\rm d}$, $\rho_{\rm f}$\\
\end{longtable}

The verification programs keep their own variable names, which need not match the notation here.

\section{Literature and the collection boundary}\label{qgs:sec:sources}

\subsection{Mathematical and physical precedents}
\TI\ Gonshor's treatment of surreal numbers and the work of van den Dries and Ehrlich provide the background for real closure, normal forms, and ordered exponential structure \cite{Gonshor,vdDE}; Conway's normal forms and real closedness are the standard foundation \cite{Conway}. Berarducci and Mantova construct derivations and transseries structure on the surreals \cite{BM}; related universality results concern differential fields \cite{ADH}. We do not need to choose such a derivation. Our main operations are finite matrix algebra and strongly summable evaluation of formal series at infinitesimals.

Proposals connecting surreal numbers to physics predate this report. Nieto discusses possible links with singularities, supersymmetry, qubits, twistors, and matroids \cite{Nieto2016,Nieto2018}. These motivate questions, but do not establish that extending the scalar field resolves a physical singularity or changes an experimentally observable probability. Non-Archimedean quantum formalisms also exist independently of surreal numbers: Benci's ultrafunction approach uses a different nonstandard construction and a different analytic infrastructure \cite{Benci}. We do not identify a Hahn field with a hyperreal enlargement or assume a transfer principle. Formal non-Archimedean quantum states have a serious precedent in deformation quantization: the formal GNS construction of Bordemann and Waldmann uses ordered extensions of Laurent-series fields and positive states \cite{BW}. It is not used here for a new GNS theorem.

The ordinary theory of effective Hamiltonians is extensive. Bravyi, DiVincenzo, and Loss discuss exact and perturbative Schrieffer--Wolff transformations and the relation between low-energy subspaces \cite{BDL}; the Feshbach--Schur framework explicitly retains an energy-dependent effective operator \cite{DSS}; shorted operators and generalized Schur complements go back at least to Anderson and Trapp \cite{AT}, with the finite positive-matrix formulas collected by Gallier \cite{Gallier}. None of our elimination theorems is a discovery of virtual transitions, invariant graphs, or the Schur complement. The local deformation theory of anti-self-dual connections is organized by elliptic complexes and finite-dimensional obstructions \cite{AHS,Itoh}; our use of those ideas is an extension, not an invention.

Ordinary quantum information is treated systematically by Watrous \cite{Watrous}. The state-separation problem was solved by Chefles and Barnett \cite{CheflesBarnett}, with later geometric treatments of general priors \cite{Bagan}; Ac\'in developed the geometry of discriminating unitaries \cite{Acin}; Vidal's single-copy theorem contains the ordinary counterpart of our restricted filtering problem \cite{Vidal}. Graydon gives complex representations of quaternionic quantum dynamics and measurements \cite{Graydon}. Barnum and Wilce, and Barnum, Graydon and Wilce, analyze how compositional and tomographic hypotheses constrain Jordan-algebraic quantum theories \cite{BarnumWilce,BGW}; Moretti and Oppio obtain a complex structure for suitable quaternionic elementary relativistic systems under additional hypotheses \cite{MorettiOppio}. Consequently, neither doubling a quaternionic Hilbert space into a complex one nor recognizing a quaternionic composition problem is a novelty claim here. Wilson's lattice formulation supplies the established gauge-theoretic setting of the lattice applications \cite{Wilson}; non-Abelian holonomy in quantum control is classical \cite{ZanardiRasetti}. Valuations, Newton polygons, and tropical geometry have also already entered quantum-model analysis, for example in the study of exceptional points \cite{BanerjeeEtAl}. ``Valuations in physics'' is therefore much too broad a phrase to serve as a claim of originality.

\subsection{What the collection already contains}\label{qgs:sec:collection}
\TI\ The sources inspected the repository at two pins, \texttt{934810a} (sources 03 and 04) and \texttt{c0a36ee} (sources 07 and 08). At both pins the inspected material documented actual surreal and surcomplex field constructions, standard-part homomorphisms, positive-order Hahn evaluation, the Neumann support lemma, and formal lifting machinery \cite{repo}. The following reports are the load-bearing prior material; all their labels cited below exist in the collection at the placement commit.

\begin{itemize}
\item \emph{Physics report} \cite{repoPhysics}, \texttt{docs/physics/surreal-scalars-and-spacetime}. It contains a finite Hermitian spectral theorem (\texttt{phys:thm:spectral}) and a finite quantum-shadow theorem (\texttt{phys:thm:quantumshadow}), whose clause (d) reduces conditioning \emph{provided} the outcome probability has nonzero standard part; the rare-outcome example \texttt{phys:ex:rare}; the amplification example \texttt{phys:ex:amplify} ($H=\epsilon\sigma_z$, $T=\epsilon^{-1}$); the three tractable classes of exponentials in \texttt{phys:sub:dynamics}; and the warnings about fine topology, global infinite phases, and noncommuting long-time evolution. That conditioning boundary is the starting point of \cref{qgs:sec:rare}, not an error in the earlier result. \Cref{qgs:sec:physrel} records exactly which of its statements this report confirms, sharpens, answers, or continues.
\item \emph{Finite surreal probability} \cite{repoProbability}, \texttt{docs/surreal/finite-surreal-probability}. Its conditional leading-scale skeleton \texttt{fsp:thm:skeleton} is the commutative diagonal case of \cref{qgs:thm:rare}; its example \texttt{fsp:ex:shadowloss} is the classical analogue of \cref{qgs:thm:alljets}; its conditioning stability theorem \texttt{fsp:thm:stability} proves $\abs{P(A\mid B)-Q(A\mid B)}\le2\delta/b$, which \cref{qgs:thm:cost} improves to $\delta/\max(P(B),Q(B))$ (\cref{qgs:mrg:rem:fsp}); its precision contract \texttt{fsp:eq:precision} is the classical form of \cref{qgs:ops:cor:amplitude-precision,qgs:ops:cor:distance-scale}; and its finite multiscale softmax limit \texttt{fsp:thm:multisoft} is the classical content of \cref{qgs:obs:thm:gibbs} in an eigenbasis. Sources 03, 04 and 07 did not cite this report; source 08 did.
\item \emph{Surcomplex spectral theory} \cite{repoSpectral}, \texttt{docs/surcomplex/spectral-theory}. Its exact effective matrix (\texttt{prop:effective}: the Feshbach determinant identity and the bound $\norm{E_{12}}^2/\delta$) precedes \cref{qgs:ops:prop:feshbach,qgs:ops:thm:window}; its example \texttt{ex:secondorder} ($\lambda_-=-t^2+t^4-2t^6+5t^8+\dots$) is the three-level bright-state reduction of \cref{qgs:obs:sec:threelevel} at $r=d=1$; its separated-subspace bound \texttt{thm:subspace} estimates what \cref{qgs:obs:thm:block} constructs exactly. Source 08 credited this report; sources 03 and 07 did not.
\item \emph{Surquaternions} \cite{repoQuaternion}, \texttt{docs/surquaternions/surquaternions}. It develops the Hamilton algebra over real closed fields, the $2\times2$ complex realization, the Hermitian spectral theorem \texttt{squat:thm:spectral}, filtered commutators and the associated graded rotation algebra \texttt{squat:thm:graded}, and, in \texttt{squat:prop:bch}, the group commutator $e^xe^ye^{-x}e^{-y}=1+[x,y]+R$ with $v(R)\ge\alpha+\beta+\min\{\alpha,\beta\}$ (\texttt{squat:eq:groupcomm}). That remainder is stronger than the one in \cref{qgs:prop:plaquette}. Its guide records ``no physics''; instrument simulation, the composite obstruction and the exact commutator norm identities here are new to the collection.
\item \emph{Vector and tensor fields} \cite{repoVTF}, \texttt{docs/surreal/vector-and-tensor-fields}. It defines finite-rank gauge fields $F_A=dA+A\wedge A$ and Yang--Mills-type equations, proves the exact doubling of the Maxwell energy scale (\texttt{vtf:mink:thm:energyscale}), the abelian analogue of \cref{qgs:obs:thm:action}, and lists among its open directions (\texttt{vtf:subsec:open}) ``support certificates for specific gauge-fixed nonlinear systems''. \Cref{qgs:obs:thm:kuranishi} partly answers that item; see \cref{qgs:sec:vtfrel}.
\end{itemize}
Further neighbours without overlap: the two-scale positivity theorem \texttt{hnd:thm:twoscale} of the hidden-negative-directions report is a neighbour of \cref{qgs:sec:schur}; Part IV of the infinite-dimensional spectral report treats compact self-adjoint operators under noncommuting Hahn perturbations, the setting of the first open question of \cref{qgs:sec:open}; \texttt{e3:thm:holonomy} uses ``holonomy'' for spherical excess, unrelated to the loop holonomies here.

\subsection{Repository statements in the sources, checked}\label{qgs:sec:stale}
The sources' descriptions of the repository were checked against the collection at the placement commit.
\begin{itemize}
\item True as stated: source 03's list of prior material (standard part, strong Hahn summation, finite Hermitian spectral theory, noncommutative support-controlled lifting, spectral estimates at surreal gap scales, a simple long-time example); source 04's description of the physics, surquaternion and spectral reports; source 07's description of the conditioning boundary; source 08's comparison table (\cref{qgs:ops:tab:overlap}).
\item \emph{Stale or incomplete, corrected here.} At its pin \texttt{934810a}, source 04 reported that a code search for ``postselection'' returned no hits (it rightly refused to infer absence): the physics report does use the word, in the status notes of \texttt{phys:ex:rare} and \texttt{phys:ex:amplify}. At its pin \texttt{c0a36ee}, source 08 reported that code searches for Schur-related terms returned no matches (again without inferring absence): the spectral-theory, hidden-negative-directions and Hahn--Hilbert-geometry reports use Schur complements. Sources 03, 04 and 07 did not cite the finite-probability report, whose commutative antecedents are now credited at every point of use.
\end{itemize}

\subsection{What was and was not inspected}
\TA\ Each source's review was targeted, not a full audit of every repository file. Source 07 inspected the README and the physics guide with selected code-index searches; source 03 the root README, the opening of the formalization ledger, and the physics and surquaternion guides; source 04 the physics, surquaternion and spectral guides; source 08 the README, the catalogue, the physics, probability and surquaternion guides, and lines 1--260 of the spectral article. None built the project, audited its Lean source, or cloned it locally (the sources used a connected repository reader), and none inferred the formal status of a theorem from the presence of a related report. The repository reports a pinned Lean/mathlib version and an axiom audit; that report is not a substitute for a fresh build, and none of the results here is an addition to the repository's verified theorem inventory. Detailed source-by-source inspection boundaries are in \cref{qgs:app:audit}.

\section{Set-sized surreal workspaces and strong evaluation}\label{qgs:sec:workspace}

\subsection{Fields, valuation, and residue}\label{qgs:sec:fields}
Let $\Gamma$ be a set-sized divisible totally ordered abelian group, written additively, and nonzero. Put
\begin{equation}\label{qgs:eq:workspace}
 F=\R((t^\Gamma)),\qquad K=\C((t^\Gamma))=F[i].
\end{equation}
An element is a Hahn series $x=\sum_{\gamma\in\Sigma}x_\gamma t^\gamma$ with well-ordered support $\Sigma\subseteq\Gamma$. The order on $F$ is determined by the sign of the coefficient of the least exponent. Thus $t^\gamma$ is positive infinitesimal when $\gamma>0$. \TI\ We use the classical facts that $F$ is real closed and $K$ is algebraically closed; divisibility of $\Gamma$ is a genuine hypothesis of these field assertions, although much of the support calculus does not need it. These are foundational inputs, not new conclusions \cite{Gonshor,vdDE}. Examples with rational exponents are understood in an ordered copy of $\Q$ inside $\Gamma$, and a bare $t$ means $t^1$ in that copy.

For nonzero $x$, set
\begin{equation}
 v(x)=\min\supp(x),\qquad
 \lc(x)=x_{v(x)},\qquad v(0)=+\infty.
\end{equation}
Write $\OO_F=\{x:v(x)\ge0\}$ and $\mm_F=\{x:v(x)>0\}$, and $\OO_K=\OO_F[i]$, $\mm_K=\mm_F[i]$. Then $\OO_F$ is exactly the ring of elements bounded in absolute value by an ordinary real constant (\emph{limited} or \emph{finite} elements), and $\mm_F$ consists exactly of the infinitesimals. For $z=x+iy\in K$, conjugation fixes $F$, $\abs z=\sqrt{x^2+y^2}$, and $v(z)=\min\{v(x),v(y)\}=v(\abs z)$. The coefficient map
\begin{equation}
 \st:\OO_F\longrightarrow\R,\qquad \st(x)=x_0,
\end{equation}
is an ordered ring homomorphism with kernel $\mm_F$, and extends coordinatewise to $\OO_K\to\C$. It is multiplicative because two nonnegative exponents sum to zero only when both are zero. An element of $\OO_F$ is \emph{appreciable} if its standard part is nonzero. For positive limited $p$, $\st p>0$ exactly when $v(p)=0$, and then $p^{-1}$ is limited with $\st(p^{-1})=(\st p)^{-1}$. We write $x\ll y$ for $x/y\in\mm_F$ when $x,y>0$, and $x\sim y$ for nonzero scalars when $x/y\in1+\mm$.

For a finite matrix, vector or quaternion $M$ over $K$ or $F$, define $v(M)$ as the minimum valuation of its entries (or real coordinates), with $v(0)=+\infty$. Matrix standard part is entrywise. The inequalities
\begin{equation}\label{qgs:eq:matrixval}
 v(M+N)\ge\min(v(M),v(N)),\qquad v(MN)\ge v(M)+v(N)
\end{equation}
are immediate. If $D$ and $D^{-1}$ have finite entries, then $v(DM)=v(M)$. In particular, this applies to $D\in M_n(\OO_K)$ with invertible $\st D$. For positive $a\in F$, $M=O(a)$ means that $a^{-1}M$ has finite entries; $M=\Oval{\alpha}$ means $v(M)\ge\alpha$ and $M=o_v(t^\alpha)$ means $v(M)>\alpha$. These are valuation statements, not uniform numerical error estimates for an unspecified real parameter.

\begin{lemma}[\TC\ No cancellation in a sum of squares]\label{qgs:lem:squares}
If $x_1,\ldots,x_m\in K$ are not all zero and $r=\min_jv(x_j)$, then
\[
 v\!\left(\sum_j\overline{x_j}x_j\right)=2r,
 \qquad
 \st\!\left(t^{-2r}\sum_j\overline{x_j}x_j\right)
 =\sum_j\abs{\st(t^{-r}x_j)}^2>0.
\]
The same formula holds for a finite sum of squared matrix entries. For a nonzero finite matrix $R$ with $\gamma=v(R)$, one has $v\norm R_\op=v\norm R_\HS=\gamma$ and $\norm R_\HS^2=\tr(RR^*)$.
\end{lemma}
\begin{proof}
All $t^{-r}x_j$ are finite and at least one has nonzero residue. Standard part commutes with conjugation and finite sums and products. The resulting ordinary sum of nonnegative real squares is strictly positive, so no cancellation raises its valuation. For the norm assertion use this identity for the Hilbert--Schmidt norm and $\norm R_\op\le\norm R_\HS\le\sqrt{\min(m,n)}\norm R_\op$ for an $m\times n$ matrix, which follows from its singular values (\cref{qgs:prop:spectral}).
\end{proof}

\TA\ The absence of cancellation in this lemma is indispensable, and is used repeatedly below. Amplitudes can cancel coherently before a Gram matrix is formed; positive probabilities cannot cancel after it is formed.

\subsection{Arbitrary real closed fields}\label{qgs:hol:sec:scalars}
Many finite results below need no Hahn presentation. Fix a real closed ordered field $F\supseteq\R$. Its ring of finite elements and ideal of infinitesimals are
\[
 \OO=\{x:\abs{x}\leq r\text{ for some }r\in\R_{>0}\},\qquad
 \mm=\{x:\abs{x}<r\text{ for every }r\in\R_{>0}\}.
\]
Every $x\in\OO$ has a unique real number $\st(x)$ with $x-\st(x)\in\mm$: take the real supremum of $\{r\in\R:r\leq x\}$; its defining inequalities give an infinitesimal difference. The map $\st:\OO\to\R$ is an ordered ring homomorphism with kernel $\mm$, applied coordinatewise to complex numbers, quaternions, and matrices. The \emph{natural valuation} has valuation ring $\OO$ and is written additively, $v:F^\times\to\Gamma$, $v(xy)=v(x)+v(y)$, $v(x+y)\ge\min\{v(x),v(y)\}$. Positive valuation means infinitesimal, and a nonzero finite number has valuation zero precisely when it is appreciable. Since $F$ is real closed, positive elements have positive $n$th roots; consequently $\Gamma$ is divisible and expressions such as $\beta/2$ are unambiguous. The group need not embed in $\R$. On a Hahn field the natural valuation is the Hahn valuation.

Positivity prevents cancellation in a sum of squares, so
\begin{equation}\label{qgs:hol:eq:normval}
 v\!\left(\sqrt{x_1^2+\cdots+x_n^2}\right)=\min_jv(x_j)
\end{equation}
when the vector is nonzero: divide by a coordinate of minimal valuation; the sum of squared normalized coordinates is a finite positive unit. A statement $r=O_v(x^k)$ means $v(r)\ge k\,v(x)$; it is a valuation statement, not a claim about a topological limit.

\TC\ The surreal field is real closed, and every finite collection of surreal numbers belongs to a set-sized real closed subfield containing $\R$: take the ordered field generated by the coefficients and then its real algebraic closure inside $\No$. Thus all finite algebra over an arbitrary real closed $F$ applies to full surreal input without forming a proper-class Hilbert space \cite{Conway}. The sections that use only this generality are marked; they do not assume that an arbitrary Hahn field is closed under a global exponential.

\subsection{How the workspace becomes surreal}\label{qgs:sec:surreal}
Choose an additive order embedding $\iota:\Gamma\hookrightarrow\No$. Conway normal forms then identify the Hahn expression with
\begin{equation}\label{qgs:eq:embedding}
 \sum_\gamma a_\gamma t^\gamma\ \longmapsto\ \sum_\gamma a_\gamma\omega^{-\iota(\gamma)}.
\end{equation}
The support on the right is reverse well ordered, as required for a surreal normal form; the minus sign reverses the support convention. Addition, multiplication, order, and conjugation are respected. Thus the results below are about genuine surreal coefficients after this embedding has been chosen, not merely about a suggestive notation. The repository's normal-form bridge provides the framework \cite{repoLedger}. We do not assert a canonical embedding for arbitrary $\Gamma$, and no theorem uses a class-sized vector space or a sum indexed by the proper class. The Conway $\omega$-map, a scalar derivation, a formal time derivative, and a physical time evolution are not identified; the existence of powerful surreal differential structures \cite{BM} does not by itself license such an identification.

A concrete and sufficient example is $\Gamma=\Q^2$ with lexicographic order and $\iota(a,b)=a\omega+b$. Set
\begin{equation}\label{qgs:eq:ranktwo}
 \epsilon=t^{(0,1)},\qquad \zeta=t^{(1,0)}.
\end{equation}
Then $0<\zeta\ll\epsilon^N$ for every ordinary $N\in\N$, while their surreal images are $\omega^{-1}$ and $\omega^{-\omega}$. \TA\ This does \emph{not} identify $\zeta$ with a particular surreal exponential. It records an order hierarchy. A physical tunnelling law or instanton action would have to justify a corresponding hierarchy independently.

Conversely, finitely many surreal inputs have set-sized normal-form supports. The divisible additive subgroup of $\No$ generated by the exponents in those supports is set sized. They therefore fit into a workspace of this form (compare \texttt{found:thm:workspace} of the foundations report). The proofs still concern fixed finite matrices over a set-sized field, not a proper-class Hilbert space. \TA\ The main common ingredient of this report is not the proper-class size of $\No$; it is a set-sized Hahn workspace with strong summation. Once embedded, its theorems are theorems about the corresponding surreal, surcomplex, or surquaternionic quantities; conversely most proofs remain valid without mentioning $\No$ at all, which identifies the actual source of each conclusion.

\subsection{Strong summability is not sequential convergence}\label{qgs:sec:strong}
A family of Hahn series is \emph{strongly summable} if the union of its supports is well ordered and only finitely many members contribute to each exponent. Its sum is then defined coefficientwise. Positive-order substitution into a formal power series is justified by the following form of the Neumann support lemma, which is also part of the repository's documented Hahn infrastructure \cite{repo,repoQuaternion}.

\begin{lemma}[\TC\ Positive-support evaluation]\label{qgs:obs:lem:support}
Let $\Sigma\subset\Gamma_{>0}$ be well ordered and let $\Sigma^*$ be the additive monoid it generates, including $0$. Then $\Sigma^*$ is well ordered, and each $\gamma\in\Gamma$ has only finitely many presentations as an ordered finite word of elements of $\Sigma$. Consequently, every formal power series in finitely many variables with ordinary coefficients can be evaluated strongly at arguments with positive supports contained in $\Sigma$. The same assertion holds for formal series with coefficients in an ordinary vector space, provided each homogeneous degree is a finite polynomial.
\end{lemma}
\begin{proof}
Use the finite-word well-quasi-order theorem on the well-ordered alphabet $\Sigma$, ordered by its numerical order. If one word embeds in another with each letter weakly increased, its sum weakly increases; because all letters are positive, equality forces equal lengths and equality of every corresponding letter. An infinite family of distinct words of one sum would therefore contradict the finite-word theorem. A strictly descending infinite sequence of sums, represented by words, gives the same contradiction without even requiring equality. Thus the generated monoid is well ordered and the finite-presentation assertion follows.

Expand a monomial in the input series into words of support exponents. For any fixed output exponent, the first assertion leaves finitely many word lengths and finitely many words. There are only finitely many variable labels at each letter and finitely many coefficients in each homogeneous polynomial. Hence every coefficient receives finitely many contributions, and all output supports lie in $\Sigma^*$. This proves the strong evaluation claim, including vector-valued coefficients.
\end{proof}

The finite-word theorem is the standard combinatorial ingredient of the Neumann support lemma; the form above, including the need to control all word lengths, is already used in the collection's surquaternion report \cite{repoQuaternion}. It is recalled because it validates the nonlinear constructions of \cref{qgs:sec:block,qgs:sec:gauge}.

The distinction between strong summation and convergence matters when $\Gamma$ has higher rank.

\begin{example}[\TC\ Why positive valuation is not topological convergence]\label{qgs:obs:ex:noconv}
In the rank-two workspace \eqref{qgs:eq:ranktwo}, $n\,v(\epsilon)<v(\zeta)$ for every ordinary $n\in\N$. The series $\sum_{n\ge0}\epsilon^n$ is a legitimate strong Hahn series and equals $(1-\epsilon)^{-1}$. Its finite partial sums do not converge to it in the full valuation topology: the valuation of each nonzero remainder is $(n+1)v(\epsilon)<v(\zeta)$. Thus an argument of the form ``the valuation goes to infinity, so Picard iteration converges'' would be invalid in the generality of this report.
\end{example}

One must not justify the existence of such series by saying that $n\,v(\epsilon)$ eventually exceeds every valuation threshold. It does not. All infinite series used below are justified by strong support conditions or by finite algebra, not by that incorrect cofinality argument: one first solves by formal homogeneous degree, then evaluates the resulting formal series using \cref{qgs:obs:lem:support}, and proves uniqueness by a separate leading-valuation argument. No cofinal positive element of $\Gamma$ is assumed.

\begin{lemma}[\TC\ Positive-order formal implicit evaluation]\label{qgs:lem:implicit}
Let $P(Y,z)$ be a finite system of formal equations in $Y-Y_0$ and $z$ over $\R$ or $\C$, where the unknown $Y$ has finitely many coordinates, $P(Y_0,0)=0$, and the Jacobian $\partial_YP(Y_0,0)$ is invertible. Then there is a unique formal solution $Y(z)$ with constant term $Y_0$. If the parameter tuple $z$ is evaluated at finitely many Hahn infinitesimals, this solution evaluates strongly and satisfies $P=0$, whenever the formal substitutions appearing in $P$ have the usual zero-constant inner parameters. In the polynomial case it is the unique finite solution with residue $Y_0$.
\end{lemma}
\begin{proof}
Write $Y=Y_0+\sum_{n\ge1}Y_n$, where $Y_n$ is homogeneous of total degree $n$ in $z$. At degree $n$ the equation has the form $\partial_YP(Y_0,0)\,Y_n+R_n=0$, where $R_n$ depends only on lower-degree terms. Invertibility of the Jacobian gives existence and uniqueness recursively. Strong positive-order evaluation follows from \cref{qgs:obs:lem:support}, and permits the finite algebra and formal compositions used to verify the identity.

For uniqueness after polynomial evaluation, suppose $Y$ and $Y'$ are finite solutions with the same residue. Telescoping each monomial expresses
\[
 P(Y,z)-P(Y',z)=M(Y,Y',z)(Y-Y'),
\]
where $\st M=\partial_YP(Y_0,0)$. Hence $M$ is invertible over the finite valuation ring and $Y=Y'$. The formal existence argument does not assert computability of arbitrary Hahn supports or convergence at non-infinitesimal parameters.
\end{proof}

\begin{remark}[Coefficients that are themselves finite Hahn elements]\label{qgs:rem:coefficients}
To apply \cref{qgs:lem:implicit} with finite matrix coefficients, regard their real and imaginary coordinate differences from their residues as additional infinitesimal parameters. This gives a finite parameter list. Correlations among those parameters do not invalidate evaluation of the universal formal identity.
\end{remark}

\subsection{Two reductions that must not be confused}\label{qgs:obs:sec:reductions}
There are two operations with different purposes. Standard part discards all positive valuations. Leading-term extraction first divides by a selected monomial and then takes standard part:
\begin{equation}\label{qgs:obs:eq:initial}
 \ini_\alpha(x)=\st(t^{-\alpha}x)
 \quad\text{when }v(x)\ge\alpha.
\end{equation}
\TA\ The first models a fixed-resolution readout. The second models a rescaled experiment or a coefficient-level analysis. For example, $\st(t^\alpha A)=0$ for an ordinary matrix $A$ and $\alpha>0$, but $\ini_\alpha(t^\alpha A)=A$.

\TC\ There is no unital multiplicative extension of standard part to the whole field that still sends a nonzero infinitesimal to zero: $1=t^\alpha t^{-\alpha}$ would be sent to zero. This elementary obstruction explains why normalization by an infinitesimal probability or energy scale must be performed \emph{before} ordinary reduction. It does not prohibit such normalization inside the larger field.

\subsection{Spacetime and differentiation conventions}\label{qgs:obs:sec:conventions}
All manifolds in \cref{qgs:sec:gauge} and all rescaled time variables are ordinary real objects. Differentiation acts on coefficient functions and treats every $t^\gamma$ as constant. \TA\ This is not the Berarducci--Mantova derivation of the surreal field \cite{BM}. That derivation is an important structure on $\No$, but its properties do not automatically make it the derivative in a Schr\"odinger or Yang--Mills equation. Similarly, strong summation is not a claim of continuity in the fine order topology of all surreal numbers; the physics report explains why ordinary real paths and that fine topology are not interchangeable \cite{repoPhysics}. Our choice is ordinary spacetime with controlled non-Archimedean coefficients.

\section{Finite quantum algebra and ordinary shadows}\label{qgs:sec:quantum}
\status{\TC\ Established finite-dimensional algebra, re-proved over the chosen field; sources 03, 04, 07 and 08 each prove the parts they need, and the physics report already contains the core spectral and quantum-shadow results (\texttt{phys:thm:spectral}, \texttt{phys:thm:quantumshadow}). Everything in this section except \cref{qgs:sec:expfin}, which uses strong Hahn evaluation, holds over an arbitrary real closed $F\supseteq\R$ with $K=F[i]$.}

\subsection{Spectral calculus without completeness}
The adjoint of a matrix is its conjugate transpose. A Hermitian matrix $H=H^*$ is positive semidefinite, written $H\ge0$, if $x^*Hx\ge0$ for every column vector $x$; positive definite means strict positivity for $x\neq0$. The inner product $\langle x,y\rangle=x^*y$ is conjugate-linear in the first variable, and Cauchy--Schwarz follows by subtracting the orthogonal projection onto a nonzero vector; it uses positivity of a sum of squares, not completeness.

\begin{proposition}[\TC\ Finite Hermitian spectral theorem]\label{qgs:prop:spectral}
Every Hermitian matrix over $K$ has an orthonormal eigenbasis and eigenvalues in $F$. It is positive semidefinite exactly when all eigenvalues are nonnegative. Every positive semidefinite matrix has a unique positive semidefinite square root and Gram factorizations. Every finite isometry extends to a unitary after enlarging the target if necessary. Every matrix has a singular-value decomposition; the operator norm is its largest singular value and the trace norm the sum of its singular values.
\end{proposition}
\begin{proof}
Algebraic closure of $K$ gives an eigenvector $x\ne0$. Its squared norm $x^*x$ is strictly positive. Since $x^*Hx=\lambda x^*x$ is real, the eigenvalue $\lambda$ belongs to $F$. Divide $x$ by the positive square root of $x^*x$. The orthogonal complement is invariant because $x^*Hy=(Hx)^*y=0$ for $y\perp x$. Induction gives the orthonormal eigenbasis. Positivity is then equivalent to nonnegative diagonal entries. Taking their nonnegative square roots constructs the positive square root. For uniqueness, a positive square root commutes with its square and hence preserves each eigenspace of $H$; diagonalizing its restriction forces the nonnegative root on each such space. Gram--Schmidt, including the choice of a basis of an orthogonal complement, extends isometries. Diagonalizing $M^*M$ gives singular values.
\end{proof}

This is proved in the physics report (\texttt{phys:thm:spectral}), in the surquaternion report over $\HF$ (\texttt{squat:thm:spectral}), and in all four sources; the proof requires no completeness of $F$. \TA\ It is a finite algebraic theorem, not an infinite-dimensional spectral theorem for unbounded operators.

\begin{lemma}[\TC\ Positive matrices and reduction]\label{qgs:lem:finitequantum}
If $X\ge0$ and $\tr X$ is limited, then every entry of $X$ is limited and its residue $X_0$ is positive semidefinite. States ($\rho\ge0$, $\tr\rho=1$), effects ($0\le E\le I$) and every Kraus matrix $M$ of a family with $\sum M^*M\le I$ have finite entries. If $M$ is a limited matrix, its singular values are limited and
\begin{equation}\label{qgs:ops:eq:norm-reduce}
 \st\norm M_1=\norm{\st M}_1,\qquad
 \st\norm M_\op=\norm{\st M}_\op;
\end{equation}
for limited Hermitian $H$ this is $\st\norm H_1=\norm{\st H}_1$, and its eigenvalues, positive and negative parts and absolute value are limited. For limited $X\ge0$, $\st(X^{1/2})=(\st X)^{1/2}$. Standard part sends states, effects and instruments to ordinary complex states, effects and instruments.
\end{lemma}
\begin{proof}
The diagonal entries of $X$ are nonnegative and bounded by $\tr X$. Every two-coordinate restriction gives $\abs{X_{jk}}^2\le X_{jj}X_{kk}$, including the case of a zero diagonal entry. Thus all entries are limited. For a real or complex constant vector $x$, $x^*X_0x=\st(x^*Xx)\ge0$, so $X_0\ge0$. The same argument applies to $E$ and $I-E$. For a Kraus matrix, each column's sum of squared moduli is a diagonal entry of $M^*M\le I$, hence at most one, so its entries are finite; no extra bounded-apparatus assumption is hidden in finite Kraus models.

A unitary matrix has limited entries because the squared moduli in each column sum to one, and its residue is again unitary. For finite Hermitian $H$ and a normalized eigenvector, $\lambda=x^*Hx$ is a finite sum of finite terms; therefore every eigenvalue is finite, and the spectral formulas for $\abs H$ and $H_\pm$ have finite entries. Diagonalize the limited positive matrix $M^*M$; its eigenvalues are bounded by its limited trace. Taking residues of the diagonalization and positive square roots shows that the singular-value multiset reduces to the ordinary singular-value multiset, including coincident or zero residue values; finite sums and the maximum of a finite ordered list commute with standard part, and $\st\abs\lambda=\abs{\st\lambda}$. This proves \eqref{qgs:ops:eq:norm-reduce}. The same diagonalization applied to $X$ proves the square-root statement. Positivity of residues follows by applying quadratic forms to ordinary complex vectors; traces, adjoints, sums, and products are preserved.
\end{proof}

\subsection{Instruments, dilations, and the operational shadow}\label{qgs:sec:instruments}
A finite instrument consists of maps
\begin{equation}\label{qgs:eq:kraus}
 \Phi_a(X)=\sum_{j=1}^{m_a}M_{a,j}XM_{a,j}^*,
 \qquad \sum_{a,j}M_{a,j}^*M_{a,j}=I,
\end{equation}
where $a$ records an observed outcome and $j$ an unobserved Kraus label. Its branch probability and conditional state are $p_a=\tr\Phi_a(\rho)$ and $\rho_a=\Phi_a(\rho)/p_a$ when $p_a>0$. For a single selected event the family satisfies $\sum_jM_j^*M_j\le I$; we call this a \emph{success map} or \emph{branch}. All index sets in this report are ordinary finite sets. The finite Kraus form is taken as the \emph{definition} of the allowed completely positive maps; it remains positive after adjoining any ordinary finite ancilla, since each ampliation has the same Kraus form, and no representation theorem for maps on an infinite operator algebra is used. These definitions fix the operational model; they are not derived from the choice of scalars.

\TC\ A channel has a finite isometric dilation $x\mapsto\sum_aM_ax\otimes\ket a$. A mixed initial state can be purified using its spectral decomposition. Measurements can be deferred by recording their outcomes in orthogonal ancilla states and applying subsequent controls conditionally. These are finite constructions over $K$ by \cref{qgs:prop:spectral}; they justify allowing pure global states in upper bounds for adaptive protocols (\cref{qgs:sec:discrimination}).

\begin{proposition}[\TC\ Finite operational shadow]\label{qgs:prop:shadow}
Consider a protocol over $K$ with ordinary finite Hilbert-space dimensions, finitely many outcomes at every step, and an ordinary finite number of instrument steps, allowing finite ancillas, adaptive choices based on recorded outcomes, and classical randomization by finite probabilities. The standard parts of all \emph{unnormalized} branch states and all outcome probabilities are exactly those of the ordinary complex protocol obtained by applying standard part to the initial state and to every Kraus operator and probability; the reduction commutes with composition and finite tensor products. If an outcome has probability $p$ with $\st p>0$, its normalized conditional output also reduces to the ordinary conditional output,
\begin{equation}\label{qgs:ops:eq:ordinary-conditioning}
 \st\rho_a=\frac{\Phi_{a,0}(\rho_0)}{\tr\Phi_{a,0}(\rho_0)},
\end{equation}
where the subscript $0$ denotes residues. This conclusion need not hold, and the right side need not even be defined, when $p$ is nonzero but infinitesimal.
\end{proposition}
\begin{proof}
The reduced Kraus operators satisfy the same completeness equation because standard part commutes with finite sums, products, and adjoints; by \cref{qgs:lem:finitequantum} they define ordinary complex instruments, and the same holds for states, effects and finite classical probability vectors. Every unnormalized branch output is a finite sum of expressions $L\rho L^*$, where $L$ is a finite product of Kraus matrices, including the chosen ancilla and classical factors, and adaptive choices change the products for each path but not the argument. Reduce term by term; the branch trace reduces in the same way. Since the protocol tree is finite, induction along it proves the statement for all branches and finite sums of branches, using unnormalized branches so that zero-shadow outcomes cause no divisions. If $\st p>0$, then $p$ is a unit of $\OO$ and $\st(p^{-1})=(\st p)^{-1}$, so normalization commutes with reduction. If $p\in\mm$, division by $p$ is not an operation in $\OO$, the residue of $p$ is zero, and the right side of \eqref{qgs:ops:eq:ordinary-conditioning} has a zero denominator. \Cref{qgs:thm:alljets,qgs:thm:cost} show that the failure can be operationally maximal conditional on success.
\end{proof}

\status{\TC\ Proved in sources 03 (with adaptive protocols), 04, 07 and 08; it is the instrument form of \texttt{phys:thm:quantumshadow}(d) and reproduces its conditioning boundary \cite{repoPhysics}. \TA\ It is not a new physical postulate and not a new impossibility theorem against all infinitesimal effects. Reading $\st p$ as an observed probability is an additional modelling choice; it implies that every such unconditioned finite protocol has an ordinary complex counterpart with the same readout statistics. It deliberately says nothing about normalization when $\st p=0$; \cref{qgs:sec:rare} supplies the missing leading-scale information.}

\subsection{Trace distance and distinguishability over the coefficient field}\label{qgs:sec:tracenorm}
For Hermitian $X$ with eigenvalues $\lambda_j\in F$, define
\[
 \norm{X}_1=\sum_j\abs{\lambda_j},\qquad
 \TD(\rho,\sigma)=\tfrac12\norm{\rho-\sigma}_1.
\]
These are $F$-valued quantities; for states $0\le\TD\le1$, and standard part commutes with $\TD$ on states by \cref{qgs:lem:finitequantum}.

\begin{lemma}[\TC\ Finite trace-norm inequalities]\label{qgs:lem:tracenorm}
For Hermitian $X$ there is an attained maximum
\begin{equation}\label{qgs:eq:tracedual}
 \norm{X}_1=\max_{-I\le T\le I}\tr(TX),
\end{equation}
and for trace-zero Hermitian $X$,
\begin{equation}\label{qgs:hol:eq:effect-dual}
 \tfrac12\norm{X}_1=\max_{0\le E\le I}\abs{\tr(EX)}.
\end{equation}
The trace norm satisfies the triangle inequality on Hermitian matrices, and $\norm{X\otimes Y}_1=\norm{X}_1\norm{Y}_1$. A positive trace-nonincreasing linear map contracts it on Hermitian inputs; in particular every finite Kraus branch, every finite channel, and partial trace do so.
\end{lemma}
\begin{proof}
Diagonalize $X=U\diag(\lambda_j)U^*$. For $-I\le T\le I$, every diagonal entry of $U^*TU$ lies in $[-1,1]$, so $\tr(TX)\le\sum_j\abs{\lambda_j}$; equality holds for the spectral sign matrix of $X$, taking any value in $[-1,1]$ on its kernel. This is an attained algebraic maximum, not a compactness argument. For trace-zero $X$ the positive and negative spectral subspaces give the witnesses of \eqref{qgs:hol:eq:effect-dual}. Applying \eqref{qgs:eq:tracedual} to a maximizing sign matrix for a sum proves the triangle inequality, and tensoring eigenbases gives the product formula.

Write $X=X_+-X_-$ with $X_\pm\ge0$ and $\norm{X}_1=\tr X_++\tr X_-$. For a positive trace-nonincreasing $\Psi$ and every $-I\le T\le I$,
\[
 \tr T\Psi(X)\le\tr\Psi(X_+)+\tr\Psi(X_-)\le\tr X_++\tr X_-,
\]
since $\tr AB=\tr(B^{1/2}AB^{1/2})\ge0$ for positive $A,B$; maximize over $T$. A Kraus branch is positive and has $\tr\Phi(Y)\le\tr Y$ for $Y\ge0$ directly from $\sum M_j^*M_j\le I$. Partial trace is a finite Kraus channel in a product basis.
\end{proof}

These are algebraic versions of familiar finite-dimensional facts; the ordinary contraction statement appears, for example, as Corollary~3.40 of \cite{Watrous}. Only the Hermitian case, proved above, is needed; neither an operator algebra over a Banach field nor a supremum over an infinite family is assumed.

\TC\ By \eqref{qgs:hol:eq:effect-dual}, $\TD$ is the largest probability difference of a binary effect. For equally likely hypotheses, using effect $E$ to announce the first state succeeds with probability $\tfrac12[1+\tr(E(\rho-\sigma))]$, so the optimal minimum-error success probability is the finite Helstrom formula
\begin{equation}\label{qgs:hol:eq:helstrom}
 P_{\mathrm{ME}}=\tfrac12\bigl(1+\TD(\rho,\sigma)\bigr).
\end{equation}

\begin{lemma}[\TC\ Pure-state distance]\label{qgs:hol:lem:puredistance}
For unit vectors $x,y$ with $c=\abs{x^*y}$, $\TD(xx^*,yy^*)=\sqrt{1-c^2}$.
\end{lemma}
\begin{proof}
Only the span of $x,y$ matters. After changing a phase, write $y=cx+\sqrt{1-c^2}\,z$, where $z$ is a unit vector perpendicular to $x$; the cases $c=1$ and one-dimensional span are immediate. In the basis $x,z$, the difference $xx^*-yy^*$ has trace zero and determinant $-(1-c^2)$. Its two eigenvalues are $\pm\sqrt{1-c^2}$.
\end{proof}

The next lemma replaces the Fubini--Study angle by an algebraic inequality, avoiding all questions about defining $\arccos$ or a global surreal phase.

\begin{lemma}[\TC\ Overlap triangle]\label{qgs:hol:lem:triangle}
For unit vectors $x,y,z$, set $a=\abs{x^*y}$ and $b=\abs{y^*z}$. Then
\begin{equation}\label{qgs:hol:eq:triangle}
 \abs{x^*z}\geq ab-\sqrt{1-a^2}\sqrt{1-b^2}.
\end{equation}
If $a\geq a_0$, $b\geq b_0$, and $a_0,b_0\in[0,1]$, the right side is at least $a_0b_0-\sqrt{1-a_0^2}\sqrt{1-b_0^2}$.
\end{lemma}
\begin{proof}
Choose phases so the coefficients of $y$ in the orthogonal decompositions of $x$ and $z$ are nonnegative. Write $x=ay+x_\perp$ and $z=by+z_\perp$. Then $x^*z=ab+x_\perp^*z_\perp$, $\norm{x_\perp}=\sqrt{1-a^2}$ and $\norm{z_\perp}=\sqrt{1-b^2}$. The reverse triangle inequality and Cauchy--Schwarz prove \eqref{qgs:hol:eq:triangle}. The right side is increasing in each variable on $[0,1]$: the product term increases while the subtracted nonnegative square-root product decreases.
\end{proof}

\subsection{Finite-time evolution, and the excluded operation}\label{qgs:sec:expfin}
For finite $z\in K$, source 07 defines
\begin{equation}\label{qgs:eq:expf}
 \expfin(z)=e^{\st z}\sum_{n\ge0}\frac{(z-\st z)^n}{n!},
\end{equation}
a strongly summable series. For a finite Hermitian $Z$, diagonalization defines $\expfin(-isZ)$ at ordinary real time $s$. Source 03 defines the same operator directly: for $Z=Z^*\in M_m(\OO_K)$ with $Z_0=\st Z$, the ordinary matrix function $(s,Z)\mapsto\exp(-isZ)$ is real analytic in the real and imaginary entries of $Z$; Taylor-expand it in $Z-Z_0$ and evaluate strongly at the infinitesimal entries. This is not the raw Hahn sum $\sum_n(-isZ)^n/n!$ when $Z$ has nonzero constant part: that sum would put infinitely many terms at exponent zero; the ordinary exponential of the constant part is already contained in the Taylor coefficients.

\begin{lemma}[\TC\ Reduction of finite-angle evolution]\label{qgs:obs:lem:exponential}
For ordinary $s$ the matrix $V(s)=\expfin(-isZ)$ is unitary and
\begin{equation}\label{qgs:eq:exp-shadow}
 \st V(s)=e^{-isZ_0},\qquad
 \frac{dV}{ds}=-iZV,\qquad V(s_1+s_2)=V(s_1)V(s_2),
\end{equation}
where the derivative is coefficientwise on ordinary smooth functions of $s$ and $s_1,s_2$ are ordinary. The two definitions above agree.
\end{lemma}
\begin{proof}
There are finitely many infinitesimal entry variables, so \cref{qgs:obs:lem:support} validates the Taylor evaluation, with support in their generated monoid. The ordinary analytic identities for unitarity, the differential equation and the group law give identities of formal Taylor series at $Z_0$, and evaluation preserves them. Every fixed Hahn coefficient is a finite sum of ordinary analytic coefficient functions of $s$, so termwise differentiation is legitimate; the constant term is $e^{-isZ_0}$.

\merge\ For the agreement, let $V_{07}(s)=U\diag(\expfin(-is\lambda_j))U^*$ from a unitary diagonalization of $Z$. Each Hahn coefficient of $\expfin(-is\lambda_j)$ is a finite sum of polynomials in $s$ times $e^{-is\st\lambda_j}$, and Hahn products have finite coefficient sums, so $V_{07}$ is coefficientwise smooth with $dV_{07}/ds=-iZV_{07}$ and $V_{07}(0)=I$; it is unitary. Hence $\frac{d}{ds}\bigl(V_{07}^*V\bigr)=iV_{07}^*ZV-iV_{07}^*ZV=0$ coefficientwise, so every Hahn coefficient of $V_{07}^*V$ is a constant ordinary function, equal to its value $I$ at $s=0$. Thus $V=V_{07}$. For the shadow, normalized eigenvectors have finite entries, their residues form a unitary matrix, finite eigenvalues reduce, and the scalar identity holds on each diagonal entry.
\end{proof}

This construction is sufficient for every finite effective Hamiltonian below. It is an instance of the first two tractable classes of the physics report (\texttt{phys:sub:dynamics}), and it is distinct from that report's global exponential $\operatorname{Exp}_{\mathrm{EK}}$. \TA\ We do not use $\exp(i\Lambda)$ for an arbitrary unlimited $\Lambda$, and do not claim that a global oscillatory calculus follows from real closure; in accordance with \texttt{phys:sub:dynamics}, no global unitary evolution for an unlimited Hamiltonian is inferred from the spectral theorem alone.

\section{Rare quantum outcomes and their initial states}\label{qgs:sec:rare}

\subsection{A normalization theorem below the ordinary shadow}
A branch with infinitesimal probability is not a zero branch. Its normalized state can be appreciable, but taking standard part before normalizing discards all the relevant data. The correct operation is to remove the leading common monomial first. In this section $F=\R((t^\Gamma))$ is a Hahn field.

\begin{theorem}[\TC\ Leading-Gram formula for a rare branch]\label{qgs:thm:rare}
Let $\rho=RR^*$ be a state over $K$, where $R$ has finitely many columns (for example $R=\rho^{1/2}$). Let
\[
 A=\Phi(\rho)=\sum_{j=1}^{m}M_j\rho M_j^*,\qquad p=\tr A>0,
\]
be a branch of a finite quantum instrument. Define
\begin{equation}\label{qgs:eq:rare-initial}
 r=\min_jv(M_jR),\qquad
 L_j=\st(t^{-r}M_jR),\qquad
 G=\sum_jL_jL_j^*,\qquad
 c=\tr G=\sum_j\tr(L_jL_j^*).
\end{equation}
Then $r\ge0$, $c>0$ is an ordinary real number, and
\begin{align}
 v(p)&=2r, & \lc(p)&=\st(t^{-2r}p)=c,\label{qgs:eq:rare-prob}\\
 \st(A/p)&=\frac{G}{c}=\frac{\sum_jL_jL_j^*}{\sum_j\tr(L_jL_j^*)}.\label{qgs:eq:rare-state}
\end{align}
In particular, the formula is valid whether $p$ is appreciable or infinitesimal. The quantities $r=v(p)/2$ and $G=\st(t^{-v(p)}A)$ are intrinsic to the branch output: they are independent of the chosen finite Gram factor $R$ of the input, of the Kraus representation of the branch map, and of any Gram representation $A=CC^*$ of the output. The conditional ordinary state is determined by the leading \emph{matrix} data, not merely by the leading scalar probability.
\end{theorem}
\begin{proof}
Since $\tr RR^*=1$, all entries of $R$ are finite, and every Kraus entry is finite (\cref{qgs:lem:finitequantum}); hence $r\ge0$. At least one product $M_jR$ is nonzero, or else $p=0$. Therefore $r$ is defined, all matrices $t^{-r}M_jR$ have finite entries, and at least one of their residues is nonzero. The ordinary matrix $G$ is positive and nonzero, so its trace $c$ is strictly positive.

The identity
\[
 p=\sum_{j,a,b}\abs{(M_jR)_{ab}}^2
\]
and \cref{qgs:lem:squares} give \eqref{qgs:eq:rare-prob}. Moreover,
\[
 t^{-2r}A=\sum_j(t^{-r}M_jR)(t^{-r}M_jR)^*
\]
has residue $G$, and its trace is appreciable after this rescaling. Dividing by $t^{-2r}p$ is now an operation inside the finite valuation ring, and taking standard part proves \eqref{qgs:eq:rare-state}. Finally, $r=\tfrac12v(p)$ and $G=\st(t^{-2r}A)$ depend only on the branch output $A$. For any other representation $A=CC^*$ the same argument gives $2v(C)=v(p)$, so $v(C)=r$, and its leading product is $\st(t^{-2r}A)=G$; the normalized result is identical.
\end{proof}

\status{\TC\ Proved in sources 03 (Theorem 4.3 there, with $R=\rho^{1/2}$ and the block row $[M_1R\cdots M_mR]$), 07 (Theorem 5.1) and 08 (Theorem 3.2, in the intrinsic form), with the same leading-rescaling proof. \TA\ The mechanism is elementary; the proposed contribution is its arbitrary-rank, arbitrary-Kraus-rank formulation as a normalization rule that precedes reduction. It is not a free experimental amplification (\cref{qgs:sec:cost}).}

\begin{remark}[\merge\ The commutative case is the probability report's skeleton]\label{qgs:mrg:rem:skeleton}
Let $\rho=\diag(p_1,\ldots,p_n)$ with $p_i=c_it^{a_i}(1+\epsilon_i)$, $c_i>0$, $v(\epsilon_i)>0$, and let the branch be the projection $P_B$ onto the coordinates in a nonempty set $B$. With $R=\rho^{1/2}$, the matrix $P_BR$ has diagonal entries $\sqrt{p_i}$ for $i\in B$, so $r=a(B)/2$ with $a(B)=\min_{i\in B}a_i$, and $L=\st(t^{-r}P_BR)$ has diagonal entries $\sqrt{c_i}$ exactly for $i\in B$ with $a_i=a(B)$. Then \eqref{qgs:eq:rare-state} is the diagonal matrix with entries $c_i/\sum_{a_j=a(B)}c_j$ on that leading group: the conditional leading-scale skeleton \texttt{fsp:thm:skeleton} of the finite-probability report \cite{repoProbability}. \Cref{qgs:thm:rare} is its noncommutative generalization, not a duplicate: it covers mixed inputs, several Kraus operators per branch, coherences, and cancellation among amplitudes before the branch norm is formed.
\end{remark}

\TA\ The equality $v(p)=2r$ is the positivity mechanism that makes the formula more informative than a generic quotient rule. It says that the first nonzero \emph{amplitude matrix}, not a guessed dominant term of the probability expression, determines the entire conditional residue. Mixed states do not spoil the rule because a Gram factor converts the branch to a Gram matrix. The distinction between coherent and incoherent information is substantial: if two amplitude contributions inside a single $M_jR$ have the same scale and opposite leading coefficient, they cancel before \cref{qgs:thm:rare} is applied, and the relevant $r$ increases. In contrast, different discarded Kraus labels contribute by the sum $\sum_jL_jL_j^*$, and their leading probabilities do not cancel. Replacing that sum by $(\sum_jL_j)(\sum_jL_j)^*$ would describe a different, coherently recombined experiment.

\begin{corollary}[\TC\ Stability and amplitude precision of the initial conditional state]\label{qgs:cor:initialstability}\label{qgs:ops:cor:amplitude-precision}
Suppose every product $M_jR$ is changed by a matrix of valuation strictly greater than the common minimum $r$, that is, another family $R'_j$ of the same finite shapes has $v(R'_j-M_jR)>r$ for every $j$. If the changed products define a nonzero branch output, it has the same probability valuation, the same leading probability coefficient $c$ and the same conditional shadow \eqref{qgs:eq:rare-state}. The strict inequality cannot be replaced by $\ge r$ in general. The same statements are preserved by an order-preserving additive embedding of the exponent group that leaves the coefficient field fixed.
\end{corollary}
\begin{proof}
After multiplication by $t^{-r}$, the changes have zero residue, so the $L_j$ are unchanged; apply \cref{qgs:thm:rare}. For sharpness compare the single column $R=t^r\ket0$ with $R'=t^r\ket1$: their difference has valuation $r$, their success probabilities both equal $t^{2r}$, and their conditional shadows are orthogonal. An exponent embedding preserves the least exponent, products, and coefficients at that exponent.
\end{proof}

\TA\ When the same input $\rho$ is used and $v(M'_j-M_j)>r$, the hypothesis follows because $\rho^{1/2}$ is limited. If the input also changes, the Gram factors must be compared, or the trace-distance theorem of \cref{qgs:sec:cost} can be used. No uniform Lipschitz estimate for the square root at a zero eigenvalue is assumed.

\subsection{Three worked branches}
\begin{example}[\TE\ A nontrivial conditional qubit]\label{qgs:obs:ex:qubit}
For $\epsilon=t^\gamma$, $\gamma>0$, let the input be $\ket0$ and take
\[
 M=\epsilon\begin{pmatrix}1&0\\2&0\end{pmatrix}.
\]
This is a valid trace-nonincreasing branch since $M^*M=\diag(5\epsilon^2,0)\le I$; it is completed to an instrument by $(I-M^*M)^{1/2}$. The success probability is $5\epsilon^2$ and the conditional state is exactly
\[
 \frac15\begin{pmatrix}1&2\\2&4\end{pmatrix}.
\]
The reduced success map is zero, yet the leading-Gram reduction gives a pure state with ordinary nonzero coherence. Taking standard part first would have erased the data required for that calculation. (Source 03.)
\end{example}

\begin{example}[\TE\ A mixed leading output]\label{qgs:ops:ex:mixed}
Take two output factors
\[
 R_1=t^2\begin{pmatrix}1&0\\1&t\end{pmatrix},\qquad
 R_2=t^2\begin{pmatrix}0&0\\1&0\end{pmatrix}.
\]
They give
\[
 X=R_1R_1^*+R_2R_2^*=t^4\begin{pmatrix}1&1\\1&2+t^2\end{pmatrix},\qquad
 p=t^4(3+t^2),\qquad
 \st(X/p)=\frac13\begin{pmatrix}1&1\\1&2\end{pmatrix}.
\]
This is a Gram-output example, not an assertion that arbitrary unscaled $R_j$ already form a complete instrument. To realize it as a branch, take input $I/2$, Kraus operators $\sqrt2R_j$, and add the positive square-root failure operator; their total effect is infinitesimal and therefore strictly below $I$. The conditional residue is genuinely mixed, retaining a leading coherence that no scalar list of outcome weights specifies. (Source 08.)
\end{example}

\begin{example}[\TE\ A rare component that the ordinary state forgets]\label{qgs:ex:rare3}
In the rank-two field \eqref{qgs:eq:ranktwo}, let
\[
 \psi=\frac{\ket{0}+\epsilon\ket{1}+\zeta\ket{2}}
 {\sqrt{1+\epsilon^2+\zeta^2}},\qquad
 P=\proj{1}+\proj{2}.
\]
The ordinary shadow of $\proj{\psi}$ is $\proj{0}$, so the shadow state assigns zero probability to $P$. Nevertheless the valued model gives
\[
 p=\frac{\epsilon^2+\zeta^2}{1+\epsilon^2+\zeta^2}>0,
 \qquad
 \st\!\left(\frac{P\proj{\psi}P}{p}\right)=\proj{1}.
\]
A further projection onto $\ket{2}$ has a still rarer nonzero branch, with conditional state $\proj{2}$. There is no contradiction: each conditioning step divides by a quantity that the preceding ordinary reduction had sent to zero. (Source 07.)
\end{example}

\subsection{Adaptive protocols and state-dependent orders}
For a path through a finite adaptive instrument, replace $M_j$ in \cref{qgs:thm:rare} by the appropriate product of Kraus matrices along that path. For a coarse-grained collection of recorded paths, retain all those products in the sum. Positivity prevents cancellation \emph{between} distinct recorded alternatives. It does not prevent coherent cancellation \emph{inside} an individual product or sum of amplitudes.

The order is determined by $M_jR$, not by $M_j$ alone. For example,
\[
 M=\begin{pmatrix}\zeta&0\\0&\epsilon\end{pmatrix}
\]
is a valid contraction when $0<\zeta\ll\epsilon\ll1$. On input $\ket{0}$ its success weight is $\zeta^2$; on input $\ket{1}$ it is $\epsilon^2$. A state annihilated by the leading Kraus matrix can therefore expose a much deeper scale.

\TC\ More generally, if $Z_1,Z_2$ are nonzero finite matrices, equality in $v(Z_1Z_2)\ge v(Z_1)+v(Z_2)$ holds precisely when
\[
 \st(t^{-v(Z_1)}Z_1)\,\st(t^{-v(Z_2)}Z_2)\ne0.
\]
If this product vanishes, the product's order rises or the product is zero. \TA\ Consequently, replacing a quantum circuit by a rule that merely adds the valuations of its gates loses essential information about kernels and interference. A useful scale calculation must propagate initial matrices, not only scalar orders.

\subsection{All power-order data can be insufficient}
The next obstruction is stronger than the failure of one finite truncation. In higher rank, there can be nonzero data invisible modulo \emph{every} power of one chosen parameter.

\begin{theorem}[\TC\ No conditional reconstruction from all one-scale power data]\label{qgs:thm:alljets}
Let $\epsilon,\zeta$ be as in \eqref{qgs:eq:ranktwo}, and let
\begin{align*}
 \rho_1&=(1-\zeta^2)\proj{0}+\zeta^2\proj{1},\\
 \rho_2&=(1-\zeta^2)\proj{0}+\zeta^2\proj{2}.
\end{align*}
For every ordinary positive integer $N$, the two states are equal modulo $\epsilon^N M_3(\OO_K)$. They therefore determine identical compatible data in all quotients $M_3(\OO_K)/\epsilon^N M_3(\OO_K)$. For the same projection $P=\proj{1}+\proj{2}$, however, both success probabilities are exactly $\zeta^2$, while the conditional ordinary states are the orthogonal states $\proj{1}$ and $\proj{2}$.

Thus no rule using only those one-scale quotient data, the fixed projector, and even the exact common success probability can determine the conditional ordinary state for every input state.
\end{theorem}
\begin{proof}
The only nonzero entries of $\rho_1-\rho_2$ are $\zeta^2$ and $-\zeta^2$. Since
\[
 v(\zeta^2/\epsilon^N)=(2,-N)>0
\]
in lexicographic order, this difference belongs to $\epsilon^N M_3(\OO_K)$ for every $N$. On the other hand,
\[
 P\rho_1P=\zeta^2\proj{1},\qquad
 P\rho_2P=\zeta^2\proj{2}.
\]
Their traces coincide and are nonzero. Normalization gives the claimed distinct outputs. A reconstruction rule would have identical inputs in the two cases but would have to return two different states.
\end{proof}

The algebra behind this example is the nonzero ideal
\begin{equation}
 I_\epsilon=\bigcap_{N\ge1}\epsilon^N\OO_K.
\end{equation}
The element $\zeta^2$ belongs to it. This does not happen in the same way for an ordinary one-variable formal power series ring with its separated power filtration. The higher-rank workspace retains information that a chosen one-scale completion erases. The same pair of states, with $\zeta^2$ replaced by an arbitrary infinitesimal $p$, is the tagged example used for sharpness in \cref{qgs:sec:cost}; it appears in all four sources. Its classical analogue, identical shadows with different rare-event answers, is \texttt{fsp:ex:shadowloss} \cite{repoProbability}.

\TA\ A physically suggestive real family is $\epsilon=h$, $\zeta=e^{-1/h}$ as $h\downarrow0$. Its exponentially small terms are invisible to every algebraic order in $h$. That analogy does not by itself produce a surreal embedding of all functions or a summation prescription for divergent expansions. What \cref{qgs:thm:alljets} proves is the exact algebraic obstruction once the hierarchy has been specified.

\subsection{Any ordinary channel can be hidden in a rare branch}
The failure of reconstruction from a zero-shadow branch is as broad as finite channel theory itself.

\begin{proposition}[\TC\ Universal rare embedding]\label{qgs:ops:prop:rare-embedding}
Let $\Psi(X)=\sum_jA_jXA_j^*$ be any ordinary complex finite Kraus map from a system to itself, and let $\gamma>0$. Put $E_\Psi=\sum_jA_j^*A_j$ and define
\begin{equation}\label{qgs:ops:eq:rare-embedding}
 M_j=t^\gamma A_j,\qquad
 M_f=(I-t^{2\gamma}E_\Psi)^{1/2}.
\end{equation}
These define a two-outcome instrument. The rare outcome has map $t^{2\gamma}\Psi$, its shadow map is zero, and the failure outcome has shadow the identity channel. If $\Psi$ is trace preserving, the rare success probability is exactly $t^{2\gamma}$ for every input, and its conditional map is exactly $\Psi$.
\end{proposition}
\begin{proof}
$E_\Psi$ is an ordinary positive matrix. Each eigenvalue of $t^{2\gamma}E_\Psi$ is infinitesimal, so $I-t^{2\gamma}E_\Psi>0$ and the stated positive square root exists. The sum of all effects is $t^{2\gamma}E_\Psi+M_f^*M_f=I$. Moreover $\st M_f=I$ by \cref{qgs:lem:finitequantum}. The branch formulas follow directly, and trace preservation gives the last assertion.
\end{proof}

\TA\ In particular, identical complete instrument shadows can conceal arbitrary different ordinary conditional channels. This is a negative answer to reconstruction from the shadow alone, and a positive recipe once leading data are retained. It is not a method for executing an arbitrary channel without implementing that channel: the Kraus operators $A_j$ are part of the required apparatus. Together with \cref{qgs:thm:alljets}, it confirms and sharpens the physics report's exclusion (\cref{qgs:sec:physrel}).

\subsection{Potential application and its boundary}\label{qgs:obs:sec:rareapp}
\TA\ A symbolic rare-event engine could retain a pair consisting of a success scale and a normalized residue state, $(v(p),\lc(p),\st(X/p))$, with enough leading amplitude data to verify it. The pair distinguishes a genuinely impossible branch, $p=0$, from a branch of nonzero but infinitesimal weight, and identifies when two small signals remain indistinguishable after postselection and when their conditional residues separate. This is a plausible application to multiscale measurement models and postselected perturbative calculations. It is not a mechanism for obtaining additional finite-resource information from a zero-readout-probability event. Neither the theorem nor its proof establishes an experimentally realizable apparatus with literal surreal success probabilities. A physical use needs an explicit family of ordinary systems and a rule relating its asymptotic regime to the chosen Hahn parameters.

\section{Visibility, postselection cost, and finite resources}\label{qgs:sec:cost}
This section holds over an arbitrary real closed $F\supseteq\R$ except where a Hahn monomial is displayed.

\subsection{A sharp success-weighted inequality}
\begin{theorem}[\TC\ Postselection cannot amplify distinguishability for free]\label{qgs:thm:cost}
Let $\rho,\sigma$ be states, and let $\Phi$ be the \emph{same} finite completely positive trace-nonincreasing map for both inputs, with $E_\Phi=\sum_jM_j^*M_j\le I$. Suppose
\[
 p=\tr\Phi(\rho)>0,\quad q=\tr\Phi(\sigma)>0,
 \quad \rho'=\frac{\Phi(\rho)}p,\quad \sigma'=\frac{\Phi(\sigma)}q.
\]
Then
\begin{equation}\label{qgs:eq:postcost}
 \boxed{\ \max(p,q)\,\TD(\rho',\sigma')\le \TD(\rho,\sigma).\ }
\end{equation}
The constant one and the denominator $\max(p,q)$ are optimal, even for diagonal three-level states and a projective branch with $p=q$. The theorem remains valid with an arbitrary fixed finite ancilla.
\end{theorem}
\begin{proof}
Put $A=\Phi(\rho)$ and $B=\Phi(\sigma)$. Append a one-dimensional failure flag:
\[
 \widehat\Phi(X)=\Phi(X)\oplus\tr\bigl((I-E_\Phi)X\bigr).
\]
This map is positive, since $\tr((I-E_\Phi)X)=\tr(X^{1/2}(I-E_\Phi)X^{1/2})\ge0$ for $X\ge0$, and trace preserving. Its output difference is block diagonal, so \cref{qgs:lem:tracenorm} gives
\begin{equation}\label{qgs:eq:flagbound}
 \norm{A-B}_1+|p-q|\le\norm{\rho-\sigma}_1.
\end{equation}
Now
\[
 p(\rho'-\sigma')=(A-B)+(q-p)\sigma',\qquad
 q(\rho'-\sigma')=(A-B)+(q-p)\rho'.
\]
Both $\rho'$ and $\sigma'$ have trace norm one, so the triangle inequality and \eqref{qgs:eq:flagbound} bound each of $p\norm{\rho'-\sigma'}_1$ and $q\norm{\rho'-\sigma'}_1$ by $\norm{\rho-\sigma}_1$. Dividing by two proves \eqref{qgs:eq:postcost}. The proof applies unchanged after tensoring each Kraus operator with a finite identity.

For optimality, take any $0<p<1$ and the states of \cref{qgs:thm:alljets} with $\zeta^2$ replaced by $p$,
\[
 \rho=(1-p)\proj0+p\proj1,\qquad \sigma=(1-p)\proj0+p\proj2,
\]
and the branch $\Phi(X)=PXP$ with $P=\proj1+\proj2$. Their input trace distance is $p$, both success probabilities are $p$, and their conditional trace distance is one: equality holds. More generally (source 04), for $0<d\le p<1$ the states $\rho=\diag(1-p,p,0)$ and $\sigma=\diag(1-p,p-d,d)$ with the same branch have input distance $d$, both success probabilities $p$, and conditional distance exactly $d/p$.
\end{proof}

\status{\TC\ Proved identically, by the failure flag, in sources 04 (Theorem 8.5 there), 07 (Theorem 6.2) and 08 (Theorem 4.1); the tagged equality example appears in all four sources. \TA\ The proof is a finite data-processing argument with the discarded flag retained until the end; it is the standard method, re-proved over $F$, given here in a sharp symmetric form with an explicit valuation consequence.}

\begin{corollary}[\TC\ The weaker bound of source 03]\label{qgs:obs:cor:postselection}
If $p,q\ge p_*>0$, then
\begin{equation}\label{qgs:obs:eq:postselect-bound}
 \TD(\rho',\sigma')\le\frac{2\,\TD(\rho,\sigma)}{p_*},
\end{equation}
and, whenever the factors are nonzero,
\begin{equation}\label{qgs:obs:eq:postselect-value}
 v\bigl(\TD(\rho',\sigma')\bigr)\ge v\bigl(\TD(\rho,\sigma)\bigr)-v(p_*).
\end{equation}
If the latter difference is positive, the conditional states have identical standard parts.
\end{corollary}
\begin{proof}
\Cref{qgs:thm:cost} gives $\TD(\rho',\sigma')\le\TD(\rho,\sigma)/\max(p,q)\le\TD(\rho,\sigma)/p_*$, which is stronger than \eqref{qgs:obs:eq:postselect-bound}. Source 03's own proof is the algebraic identity $A/p-B/q=(A-B)/p+B(1/p-1/q)$, which with positivity of $B$ gives $\norm{A/p-B/q}_1\le(\norm{A-B}_1+\abs{p-q})/p\le2\norm{A-B}_1/p$, using $\abs{\tr(A-B)}\le\norm{A-B}_1$; contraction (\cref{qgs:lem:tracenorm}) and $p\ge p_*$ finish it. The factor two is lost in that second step, where the flag argument instead absorbs $\abs{p-q}$ into \eqref{qgs:eq:flagbound}. An inequality between positive elements reverses their valuations, and nonzero ordinary constants have valuation zero, proving \eqref{qgs:obs:eq:postselect-value}; positive valuation of the output distance gives zero reduced distance by \cref{qgs:lem:finitequantum}.
\end{proof}

Source 03 stated explicitly that its constant two was not claimed optimal; \cref{qgs:thm:cost} shows it is not, and that the exponent loss \eqref{qgs:obs:eq:postselect-value} is attained (by the tagged example).

\begin{remark}[\merge\ The classical case improves \texttt{fsp:thm:stability} by a factor two]\label{qgs:mrg:rem:fsp}
Let $P,Q$ be finite probability laws with total-variation distance $\delta=\max_A\abs{P(A)-Q(A)}$, and let $B$ be an event with $P(B)>0$ and $Q(B)>0$. Apply \cref{qgs:thm:cost} to the diagonal states $\rho=\diag(p_i)$, $\sigma=\diag(q_i)$ and the projection onto $B$: $\TD(\rho,\sigma)=\delta$, and the conditional trace distance is the total variation of the conditional laws, which dominates $\abs{P(A\mid B)-Q(A\mid B)}$ for every event $A$. Hence
\[
 \abs{P(A\mid B)-Q(A\mid B)}\le\frac{\delta}{\max\{P(B),Q(B)\}}.
\]
Under the hypothesis $P(B)\ge b>\delta$ of \texttt{fsp:thm:stability} \cite{repoProbability}, which in particular makes $Q(B)>0$, this is at most $\delta/b$, half the bound $2\delta/b$ printed there. The probability report does not claim that its factor two is sharp, so this is an improvement, not a correction; the constant one is attained by the tagged example. Its separate valuation contract \texttt{fsp:eq:precision} is unaffected.
\end{remark}

\begin{corollary}[\TC\ Valuation budget and residue precision threshold]\label{qgs:cor:budget}\label{qgs:ops:cor:distance-scale}
If $\TD(\rho,\sigma)$ is infinitesimal and either success probability is appreciable, the conditional outputs have the same ordinary shadow. More generally, assume $\TD(\rho,\sigma)>0$ and put
\[
 \alpha=v\bigl(\TD(\rho,\sigma)\bigr),\qquad
 \beta=v(\max(p,q))=\min\{v(p),v(q)\}.
\]
If $\TD(\rho',\sigma')>0$, then
\begin{equation}\label{qgs:ops:eq:precision}
 v\bigl(\TD(\rho',\sigma')\bigr)\ge\alpha-\beta.
\end{equation}
In particular, $\alpha>\beta$ implies $\st\rho'=\st\sigma'$, and if the conditional distance has positive standard part then $\min\{v(p),v(q)\}\ge v(\TD(\rho,\sigma))$. At the boundary $\alpha=\beta$ the shadows can be orthogonal.
\end{corollary}
\begin{proof}
For positive elements, $a\le b$ implies $v(a)\ge v(b)$, and $v(\max(p,q))=\min(v(p),v(q))$. Apply this to \eqref{qgs:eq:postcost}. If $\alpha>\beta$, the output distance is infinitesimal or zero, and its ordinary trace distance vanishes by \eqref{qgs:ops:eq:norm-reduce}. The equality example of \cref{qgs:thm:cost} with $p=t^\beta$, $\beta>0$, gives the boundary statement.
\end{proof}

\TA\ A branch of Gram-amplitude scale $t^r$ has probability scale $t^{2r}$. Two sufficient stability certificates are then available: amplitude precision beyond $r$ (\cref{qgs:ops:cor:amplitude-precision}), and, for inputs compared through a fixed branch, trace-distance precision beyond the success scale (\cref{qgs:ops:cor:distance-scale}). These are distinct sufficient contracts, not necessary conditions for every individual protocol, and not interchangeable error metrics. Their classical form is the precision contract \texttt{fsp:eq:precision} of the probability report.

\begin{corollary}[\TC\ Calibration cost of a rare amplifier]\label{qgs:hol:cor:calibration}
Suppose a nominal state and a perturbed state differ by trace distance $d$, and a nominal success probability has valuation $\beta$. If $v(d)>\beta$, the perturbed success probability has the same leading term, and the conditional outputs differ infinitesimally; their distance has valuation at least $v(d)-\beta$. In particular, a worst-case guarantee for a heralding rate of valuation $2\kappa$ (\cref{qgs:sec:discrimination}) requires input trace-distance error of valuation strictly greater than $2\kappa$.
\end{corollary}
\begin{proof}
The effect variational formula \eqref{qgs:hol:eq:effect-dual} gives $\abs{p-q}\le d$, so $q/p\in1+\mm$ under the hypothesis. Apply \eqref{qgs:eq:postcost} and the order-reversing relation between magnitude and valuation. The final statement is a sufficient worst-case requirement; the equality examples show why no better uniform dependence on $p$ is possible.
\end{proof}

\TA\ Source 04 wrote that the final statement ``requires'' the stated precision; the requirement is worst-case and sufficient, and its necessity is only in the worst case, as the equality example shows. It does not say that every particular coherent calibration error is that destructive: structured noise can be less harmful. It says that a claimed rare-event amplifier needs a \emph{noise model}, not only a large conditional signal.

\TC\ For completeness, if $m$ unitary query slots are replaced by nearby unitaries with operator-norm errors $e_1,\ldots,e_m$, telescoping products and unitary invariance give a final pure-vector error at most $\sum_je_j$. Pure-state trace distance is no larger than vector distance, since $1-\abs{x^*y}^2\le2-2\Rea(x^*y)$. After purification this bounds the output trace-distance error as well. Feeding that bound into \cref{qgs:thm:cost} supplies a finite circuit-level certification rule. No small numerical query count is required, but the count must remain ordinary finite.

\TA\ \Cref{qgs:thm:cost} is a finite idealized discrimination statement. It is not a blanket claim that postselection can never help with a particular detector's saturation or technical noise. Those are different resource models, discussed in the weak-value literature \cite{FerrieCombes,JordanEtAl,Torres}, where critiques explicitly contest the scope of blanket negative claims \cite{Vaidman}. The theorem's assumptions are explicit: the same branch map, fixed finite matrices, positivity, and success-weighted comparison. It excludes treating a large normalized signal as a cost-free increase of input distinguishability; it does not adjudicate the experimental debate.

\subsection{Ordinary finite repetitions do not expose infinitesimal differences}\label{qgs:sec:copies}
\begin{proposition}[\TC\ Finite-copy invisibility]\label{qgs:prop:copies}
For every ordinary positive integer $N$,
\[
 \TD(\rho^{\otimes N},\sigma^{\otimes N})\le N\,\TD(\rho,\sigma).
\]
Consequently, a fixed finite quantum protocol on a fixed finite number of copies, including adaptive processing, cannot turn an infinitesimal input trace distance into an appreciable unconditioned trace distance or an appreciable minimum-error advantage.
\end{proposition}
\begin{proof}
Expand the difference of tensor powers by telescoping,
\[
 \rho^{\otimes N}-\sigma^{\otimes N}=\sum_{j=1}^N\rho^{\otimes(j-1)}\otimes(\rho-\sigma)\otimes\sigma^{\otimes(N-j)}.
\]
Each of its $N$ terms is a tensor product of states and one factor $\rho-\sigma$; \cref{qgs:lem:tracenorm} bounds its trace norm by $\norm{\rho-\sigma}_1$. Sum the bounds and divide by two. Any subsequent finite adaptive processing can be written as a channel with a finite classical register and contracts the distance; apply \eqref{qgs:hol:eq:helstrom}.
\end{proof}

\TA\ There is no hidden quantification over a ``surreal number of copies'' here. An ordinary finite tensor product has been constructed; an infinite or hyperfinite product requires a separate theory. This is a substantive resource restriction, not a merely linguistic one. (Sources 03 and 07.)

\TC\ For a fixed ordinary integer $N$, the probability of at least one success in $N$ trials, each with probability at most $p$, is at most $Np$, whether or not the trials are independent: this is finite subadditivity. When independent trials are specified, the exact finite formula is
\[
 1-(1-p)^N=Np+p^2S_N(p),\qquad S_N\in\Z[p],
\]
which has the same valuation as $p$ when $p\in\mm$. The assertion is a finite binomial identity and a finite product probability calculation; it does not posit a countably additive surreal probability space. Thus infinitesimal success weights remain invisible to ordinary finite repetition counts, and a finite number of repetitions cannot make an exceptional branch appreciable. \TA\ Likewise, $v(p)>0$ does not say how many ordinary trials are needed until a success occurs; it says that every fixed ordinary trial count has zero success probability after the chosen standard-part reduction.

\TA\ Writing $p^{-1}$ in a larger field does not define a laboratory protocol consisting of an unlimited number of trials; the informal phrase ``wait for $1/p$ trials'' is not a finite operation when $p$ is a literal infinitesimal, and we do not define a transfinite sequence of trials by writing that reciprocal. One would have to specify an additional index structure, probability theory, and operational interpretation; the probability report studies some of the additional structures required \cite{repoProbability}. In a real family $p(s)>0$ one can instead choose an ordinary integer $N(s)$ and analyze $1-(1-p(s))^{N(s)}$: if $p(s)\to0$ and $N(s)p(s)\to c$, elementary logarithmic estimates give the limit $1-e^{-c}$, and the mean number of independent real Bernoulli trials before success is $1/p(s)$. This is a trial-count statement, not a universal quantum gate-complexity lower bound: a protocol that coherently reuses a preparation and its inverse must be analyzed separately rather than counted as independent repetitions. None of these observations turns an element $t^{-\gamma}$ into a finite matrix dimension or defines an infinite surreal-valued path measure. Conversely, for a real family with small but nonzero probability, \cref{qgs:thm:cost} supplies an ordinary resource tradeoff at every finite parameter value. This is the appropriate route from exact scale algebra back to quantitative physics: specialize a validated family, rather than interpret an infinite surreal directly as an experimental budget.

\subsection{Infinitesimal entanglement as a resource with a price}\label{qgs:ops:sec:entanglement}
Let $d\ge2$ be an ordinary integer, and consider the full-Schmidt-rank pure state
\begin{equation}\label{qgs:ops:eq:schmidt}
 \ket\psi=\sum_{j=1}^d\sqrt{\lambda_j}\ket{jj},\qquad
 \lambda_1\ge\cdots\ge\lambda_d>0,\qquad
 \sum_j\lambda_j=1,
\end{equation}
over an arbitrary real closed $F$. Its Schmidt representation follows from the finite singular-value decomposition; here the representation itself is given as data. A successful one-sided filter is a single contraction $M$ on the first system, with unnormalized output $(M\otimes I)\ket\psi$; a fixed local unitary correction of the successful output is allowed.

\begin{theorem}[\TC\ Exact concentration probability]\label{qgs:ops:thm:entanglement}
Among one-sided single-branch filters taking \eqref{qgs:ops:eq:schmidt} to a maximally entangled state of Schmidt rank $d$, the largest success probability is
\begin{equation}\label{qgs:ops:eq:ent-probability}
 p_{\max}=d\lambda_d.
\end{equation}
It is attained by
\begin{equation}\label{qgs:ops:eq:ent-filter}
 M=\diag\left(\sqrt{\frac{\lambda_d}{\lambda_1}},\ldots,
                 \sqrt{\frac{\lambda_d}{\lambda_d}}\right).
\end{equation}
Hence $v(p_{\max})=v(\lambda_d)$.
\end{theorem}
\begin{proof}
Put $\Lambda=\diag(\lambda_j)$. A successful maximally entangled output with probability $p$ has first reduced matrix $(p/d)I$, so $M\Lambda M^*=(p/d)I$. Since $p>0$ and $\Lambda>0$, $M$ is invertible. It follows that $M\Lambda^{1/2}=\sqrt{p/d}\,U$ for a unitary $U$, whence
\[
 M^*M=(p/d)\Lambda^{-1}\le I.
\]
Thus $p\le d\lambda_d$. For the diagonal filter \eqref{qgs:ops:eq:ent-filter}, all diagonal entries lie in $[0,1]$, and the output is $\sqrt{\lambda_d}\sum_j\ket{jj}$, whose squared norm is $d\lambda_d$. The failure filter $(I-M^*M)^{1/2}$ completes the local instrument. A positive ordinary integer has valuation zero, proving the last assertion.
\end{proof}

\TI\ The ordinary complex result is a special case of the broader optimal single-copy transformation theorem of Vidal \cite{Vidal}. \TA\ We make no new claim about optimality over an unrestricted surreal analogue of arbitrary LOCC protocols. The restriction here makes the proof finite, explicit, and sufficient for a sharp infinitesimal resource law.

\begin{example}[\TE\ A product shadow concealing a Bell pair]\label{qgs:ops:ex:bell}
Fix $a>0$ in $\Gamma$ and set
\begin{equation}\label{qgs:ops:eq:tiny-entanglement}
 \ket{\psi_a}=\frac{\ket{00}+t^a\ket{11}}{\sqrt{1+t^{2a}}}.
\end{equation}
Its ordinary density-matrix shadow is the product state $\proj{00}$. Nevertheless, applying $\diag(t^a,1)$ to the first qubit produces a Bell state on success, with
\begin{equation}\label{qgs:ops:eq:bell-cost}
 p_{\rm Bell}=\frac{2t^{2a}}{1+t^{2a}},\qquad
 v(p_{\rm Bell})=2a,\qquad \lc(p_{\rm Bell})=2.
\end{equation}
There is no contradiction with the finite shadow theorem: the successful branch has zero ordinary probability. One can quantify the unreduced entanglement without logarithms. Define negativity by $\mathcal N(\rho)=(\norm{\rho^{T_B}}_1-1)/2$ using the finite trace norm. For a two-qubit Schmidt state with weights $\lambda_1,\lambda_2$, the four partial-transpose eigenvalues are $\lambda_1,\lambda_2,\sqrt{\lambda_1\lambda_2},-\sqrt{\lambda_1\lambda_2}$. Therefore
\begin{equation}\label{qgs:ops:eq:negativity}
 \mathcal N(\proj{\psi_a})=\frac{t^a}{1+t^{2a}}.
\end{equation}
The valuation of the entanglement indicator is $a$, while the optimal one-sided Bell-production probability has valuation $2a$. Positivity of the indicator is an exact mathematical distinction from separability; its vanishing standard part means that ordinary shadow tomography forgets that distinction.
\end{example}

\TC\ For any fixed number of copies, the tensor-product input shadow remains a product state across the bipartition, and finite local instruments reduce to ordinary local instruments. Thus an entangled conditional ordinary output cannot arise with positive ordinary success probability from that fixed-copy product shadow. \TA\ This conclusion does not exclude families whose number of copies grows as a real small parameter tends to zero; those have a different resource accounting. In an asymptotic entanglement calculation one often discards small Schmidt weights too early. Keeping their valuations gives an exact warning: a state can have zero leading entanglement and still admit concentration, but the omitted smallest weight controls the yield. For the real family obtained by replacing $t^a$ with $s^a$, formula \eqref{qgs:ops:eq:bell-cost} is an ordinary two-qubit calculation valid at every $s>0$. Surreal notation organizes its limiting hierarchy; it is not required to build the filter or to establish a new experimental effect.

\section{Exact elimination of independently separated virtual sectors}\label{qgs:sec:elimination}
This section and the next two treat three different effective-Hamiltonian problems, from three sources, and all three are kept: here an \emph{unlimited} high block with critically scaled, possibly unlimited couplings and independently separated scales (source 07); in \cref{qgs:sec:block} an \emph{ordinary} gap with infinitesimal couplings and explicit support certificates (source 03); in \cref{qgs:sec:schur} a positive block whose eliminated part becomes \emph{soft} in the shadow (source 08). For a block matrix $H=\bigl(\begin{smallmatrix}A&B\\B^*&D\end{smallmatrix}\bigr)$ with invertible $D$ we write $S(H)=A-BD^{-1}B^*$ for its Schur complement.

\subsection{The model and the critical scaling}
Let $r,s$ be ordinary positive integers. Consider a low sector of dimension $r$ and a high sector of dimension $s$. Let
\[
 A=A^*\in M_r(\OO_K),\quad
 D=D^*\in M_s(\OO_K),\quad
 B\in M_{r\times s}(\OO_K),
\]
and assume $D_0=\st D$ is positive definite. Set
\begin{equation}\label{qgs:eq:Delta}
 \Delta=\diag(\delta_1,\ldots,\delta_s),
 \qquad 0<\delta_j\in\mm_F,
 \qquad \tau=\sum_{j=1}^s\delta_j^2.
\end{equation}
There is no comparison hypothesis on the ratios $\delta_j/\delta_k$. In particular, one may be smaller than every power of another.

The Hamiltonian is
\begin{equation}\label{qgs:eq:criticalH}
 H=\begin{pmatrix}
 A&B\Delta^{-1}\\
 \Delta^{-1}B^*&\Delta^{-1}D\Delta^{-1}
 \end{pmatrix}.
\end{equation}
Both the high block and some couplings may be unlimited. Nevertheless, coupling divided by the corresponding energy scale is infinitesimal. The Schur complement of $H$ is the finite Hermitian matrix
\begin{equation}\label{qgs:eq:Kzero}
 S(H)=A-(B\Delta^{-1})(\Delta D^{-1}\Delta)(\Delta^{-1}B^*)=A-BD^{-1}B^*.
\end{equation}
The adjective ``critical'' refers to the fact that the second-order virtual shift is finite rather than infinitesimal.

For positive $a\in F$, the notation $M=O(a)$ means that $a^{-1}M$ has finite entries. In a fixed finite dimension this is equivalent to an operator-norm bound by an ordinary real constant times $a$. Remainders $O(\tau^2)$ below are statements in this precise sense, not unexplained analytic limits.

\subsection{The invariant-graph theorem}
\begin{theorem}[\TC\ Exact arbitrary-rank virtual-sector elimination]\label{qgs:thm:elimination}
Under the hypotheses above, the following conclusions hold.
\begin{enumerate}[label=\textup{(\roman*)},leftmargin=2em,itemsep=3pt]
\item There is a unique finite matrix $Y\in M_{s\times r}(\OO_K)$ satisfying
\begin{equation}\label{qgs:eq:riccati}
 B^*+DY=\Delta^2Y(A+BY).
\end{equation}
Its residue is $\st Y=-D_0^{-1}B_0^*$, where $B_0=\st B$.
\item Set $X=\Delta Y$, $T_X=I+X^*X$, and
\begin{equation}\label{qgs:eq:W}
 W=\begin{pmatrix}I\\X\end{pmatrix}T_X^{-1/2}.
\end{equation}
Then $W^*W=I$ and $HW=W\Heff$, where the finite matrix
\begin{equation}\label{qgs:eq:Kexact}
 \Heff=T_X^{1/2}(A+BY)T_X^{-1/2}=W^*HW
\end{equation}
is Hermitian. Its ordinary shadow is
\begin{equation}\label{qgs:eq:Kshadow}
 \st\Heff=A_0-B_0D_0^{-1}B_0^*.
\end{equation}
\item The complementary orthogonal subspace is invariant, and its spectrum is positive and unlimited. Thus $H$ has exactly $r$ finite eigenvalues, counted with multiplicity. Its low and high subspaces are related to the original coordinate sectors by a unitary matrix differing infinitesimally from the identity.
\item Define
\begin{equation}\label{qgs:eq:C}
 \Corr=BD^{-1}\Delta^2D^{-1}B^*.
\end{equation}
Then the finite effective Hamiltonian has the controlled expansion
\begin{equation}\label{qgs:eq:Kcorrection}
 \boxed{\quad \Heff=S(H)-\tfrac12\bigl(\Corr\,S(H)+S(H)\,\Corr\bigr)+O(\tau^2).\quad}
\end{equation}
All statements allow arbitrary rank of $\Gamma$ and arbitrary relative sizes of the $\delta_j$.
\end{enumerate}
\end{theorem}

\begin{proof}
\textit{Existence.}
At $\Delta^2=0$, the equation is $B^*+DY=0$. Treat the entries of $\Delta^2$ and all finite coefficient deviations from their residues as finitely many formal infinitesimal parameters (\cref{qgs:rem:coefficients}). The base solution is $-D_0^{-1}B_0^*$, and the derivative with respect to $Y$ is left multiplication by $D_0$, which is invertible. \Cref{qgs:lem:implicit} supplies a strongly evaluated finite solution. No assumption that powers of a valuation are cofinal is used.

\textit{Uniqueness among all finite solutions.}
Let $Y,Y'$ be finite solutions, put $Z=Y-Y'$, and write $T=A+BY$. Subtracting their equations gives
\begin{equation}\label{qgs:eq:uniqueY}
 DZ=\Delta^2\bigl(ZT+Y'BZ\bigr).
\end{equation}
Since $D$ and $D^{-1}$ are finite, $v(DZ)=v(Z)$. If $Z\ne0$, the right side of \eqref{qgs:eq:uniqueY} has valuation at least $v(\Delta^2)+v(Z)>v(Z)$, because all the other matrices are finite. This is impossible. Taking residues of \eqref{qgs:eq:riccati} gives the stated residue of $Y$.

\textit{The exact graph.}
The entries of $X=\Delta Y$ are infinitesimal. Direct block multiplication using \eqref{qgs:eq:riccati} gives
\begin{equation}\label{qgs:eq:graphidentity}
 H\begin{pmatrix}I\\X\end{pmatrix}
 =\begin{pmatrix}I\\X\end{pmatrix}(A+BY).
\end{equation}
The positive matrices $T_X^{\pm1/2}$ exist either by the spectral theorem or by the strongly summable binomial series at $X^*X$. They differ infinitesimally from $I$. Hence $W^*W=I$, and \eqref{qgs:eq:graphidentity} gives $HW=W\Heff$ with \eqref{qgs:eq:Kexact}. Multiplying by $W^*$ proves $\Heff=W^*HW$, so $\Heff$ is Hermitian despite the nonsymmetric appearance of its similarity formula. All its entries are finite. Reducing $T_X$, $Y$, and $A+BY$ proves \eqref{qgs:eq:Kshadow}.

\textit{The complementary unitary frame.}
Define
\[
 V=\begin{pmatrix}-X^*\\I\end{pmatrix}(I+XX^*)^{-1/2}.
\]
Then $V^*V=I$ and $W^*V=0$. Thus $U=[\,W\ V\,]$ is a square unitary matrix, $U-I$ has infinitesimal entries, and $U^*HU$ is block diagonal. The second subspace is invariant because $H$ is Hermitian and the first is invariant.

\textit{Unlimited positivity of the complementary block.}
There is an ordinary real $c>0$ with $D\ge cI$, because $D_0$ is positive definite and $D-D_0$ has infinitesimal entries. The exact completed-square identity is
\begin{align}\label{qgs:eq:completesquare}
 \begin{pmatrix}u\\z\end{pmatrix}^{\!*}
 H\begin{pmatrix}u\\z\end{pmatrix}
 &=u^*S(H)u\\
 &\quad+(\Delta^{-1}z+D^{-1}B^*u)^*
 D(\Delta^{-1}z+D^{-1}B^*u).\nonumber
\end{align}
On the orthogonal complement of the graph one has $u=-X^*z$. Choose an ordinary real constant $b_0$ bounding the operator norms of $S(H)$ and $D^{-1}B^*$. Let $\delta_{\max}=\max_j\delta_j$. The elementary inequality $\norm{a+b}^2\ge\tfrac12\norm a^2-\norm b^2$ gives
\begin{align*}
 \begin{pmatrix}u\\z\end{pmatrix}^{\!*}
 H\begin{pmatrix}u\\z\end{pmatrix}
 &\ge \frac c2\delta_{\max}^{-2}\norm z^2
 -(b_0+cb_0^2)\norm u^2.
\end{align*}
Here $\norm u\le\norm X\norm z$, and $\norm X$ is infinitesimal. Since $\delta_{\max}^{-2}$ is positive unlimited, the right side is at least $\tfrac c4\delta_{\max}^{-2}\norm z^2$. Also $\norm{(u,z)}^2\le2\norm z^2$. Thus the restricted quadratic form is bounded below by
\[
 \frac c8\delta_{\max}^{-2}\norm{(u,z)}^2.
\]
Every eigenvalue of the complementary block is positive unlimited. Every eigenvalue of $\Heff$ is finite by the finite-entry Rayleigh-quotient argument of \cref{qgs:lem:finitequantum}. This proves the precise eigenvalue count.

\textit{The correction.}
Write $E=\Delta^2$ and $Y_0=-D^{-1}B^*$. Since $E=O(\tau)$, equation \eqref{qgs:eq:riccati} first gives $Y-Y_0=O(\tau)$, and substituting this once more gives
\begin{equation}\label{qgs:eq:Yexpansion}
 Y=-D^{-1}B^*-D^{-1}ED^{-1}B^*S(H)+O(\tau^2).
\end{equation}
Consequently,
\[
 A+BY=S(H)-\Corr\,S(H)+O(\tau^2),\qquad
 T_X=I+Y^*EY=I+\Corr+O(\tau^2).
\]
The binomial identities, with finite remainders after the quadratic term, give $T_X^{1/2}=I+\tfrac12\Corr+O(\tau^2)$ and $T_X^{-1/2}=I-\tfrac12\Corr+O(\tau^2)$. Inserting these into \eqref{qgs:eq:Kexact} yields
\[
 \Heff=S(H)-\Corr S(H)+\tfrac12\Corr S(H)-\tfrac12S(H)\Corr+O(\tau^2),
\]
which is \eqref{qgs:eq:Kcorrection}. All remainder estimates use products of finite matrices and the positive scalar $\tau$, so none assumes comparability of distinct scales.
\end{proof}

\status{\TC\ Source 07. \TI\ The invariant-graph, Schur-complement and Schrieffer--Wolff mechanisms are classical \cite{BDL}; the finite Feshbach identity and the bound $\norm{E_{12}}^2/\delta$ for a fixed eliminated block are \texttt{prop:effective} of the spectral-theory report \cite{repoSpectral}. \TA\ The proposed contribution is the arbitrary-rank coefficient field, independently separated scales, and an exact ordinary-shadow statement with a controlled finite correction.}

\begin{remark}[A uniform remainder is not a deep-scale certificate]\label{qgs:rem:uniform}
The bound $O(\tau^2)$ is uniform but deliberately coarse. If $\delta_2$ is smaller than every power of $\delta_1$, a displayed $\delta_2^2$ contribution inside $\Corr$ can be smaller than omitted $\delta_1^4$ terms. The truncated formula alone therefore does not certify that deeper contribution against its remainder. To resolve it, one must use the exact graph equation or a support-aware expansion retaining all relevant lower-valued terms. The strongly evaluated formal solution keeps this information; a single total-degree error bound does not.
\end{remark}

\subsection{Why an ordinary gap estimate is not the proof}
\TA\ If the $\delta_j$ are widely separated, the norm of the full coupling $B\Delta^{-1}$ can be much larger than the smallest high-sector energy gap. A sufficient perturbative hypothesis based only on ``coupling norm divided by the smallest gap'' may therefore be useless even though every high-sector mixing coordinate is small. The theorem instead rescales the invariant graph as $X=\Delta Y$ before solving its equation. This removes the negative powers from the implicit problem and keeps the full finite matrix $D$, including its off-diagonal couplings.

The complementary spectral estimate is likewise not obtained by separately diagonalizing each original high coordinate. It comes from the completed square \eqref{qgs:eq:completesquare} on the \emph{exact} orthogonal complement of the graph. This is why the theorem allows arbitrary relative scales without assuming that $D$ commutes with $\Delta$. These points distinguish the actual hypotheses from a loose statement that the usual perturbation series ``must still converge.'' In arbitrary rank, that statement would leave both a support problem and a scale problem unresolved.

\subsection{A real-parameter counterpart}
\TA\ The theorem has a direct interpretation for ordinary small real parameters. This is important: a useful valued calculation should point to an ordinary finite-scale statement, rather than require treating an infinitesimal as a literal laboratory knob.

\begin{corollary}[\TC\ Uniform ordinary realization]\label{qgs:cor:realrealization}
Let $A(h),B(h),D(h)$ be ordinary real or complex matrices of the same fixed dimensions as above, with $A(h),D(h)$ Hermitian. Suppose $D(h)\ge cI$ for a fixed real $c>0$, and that $A,B,D^{-1}$ are uniformly bounded. Let $\delta_j(h)>0$ satisfy $\max_j\delta_j(h)\to0$. There is a bounded solution of \eqref{qgs:eq:riccati} for all sufficiently small $h$, unique among uniformly bounded solution branches near $h=0$. It gives an exact orthonormal invariant graph and effective Hamiltonian satisfying \eqref{qgs:eq:Kcorrection}, with a remainder bounded in operator norm by a constant times $\tau(h)^2$, independent of the ratios of the $\delta_j(h)$.

If $A(h),B(h),D(h)$ have limits $A_0,B_0,D_0$ with $D_0>0$, then $\Heff(h)\to A_0-B_0D_0^{-1}B_0^*$, while the complementary eigenvalues tend to $+\infty$.
\end{corollary}
\begin{proof}
Choose $b_0\ge1$ uniformly bounding $\norm A$, $\norm B$, and $\norm{D^{-1}}$. Rewrite the equation as the fixed-point problem
\[
 Y=\mathcal T(Y):=-D^{-1}B^*+D^{-1}\Delta^2Y(A+BY).
\]
Put $R=2b_0^2+1$. On the closed ball $\norm Y\le R$, the first term has norm at most $b_0^2$ and the second at most $b_0^2\tau R(1+R)$. For sufficiently small $\tau$, this maps the ball into itself. For two matrices in the ball,
\[
 \norm{\mathcal T(Y)-\mathcal T(Y')}
 \le b_0^2\tau(1+2R)\norm{Y-Y'}.
\]
The multiplier is less than $1/2$ for sufficiently small $h$. The ordinary contraction theorem therefore supplies a unique fixed point in that ball. Any uniformly bounded solution branch satisfies $Y+D^{-1}B^*=O(\tau)$ by its equation, so it eventually belongs to the same ball and is the same solution.

The exact graph identities are finite algebra. The substitutions yielding \eqref{qgs:eq:Yexpansion} and \eqref{qgs:eq:Kcorrection} now have uniform ordinary remainder bounds because all displayed finite factors are uniformly bounded. The positive square-root expansion is uniform near the identity. The completed-square estimate applies with the same fixed lower bound $c$, so the high-sector eigenvalues diverge. Taking limits proves the final assertion.
\end{proof}

\TA\ This corollary permits, for example, $\delta_1(h)=h$ and $\delta_2(h)=e^{-1/h}$ without treating the second as a finite power of the first. It does not prove that a given continuum theory produces those functions. It proves that the finite matrix reduction remains valid if a derivation of the model has supplied them.

\subsection{Dynamics on the dressed low sector}
At an ordinary real time $s$, a low-sector initial vector $\psi_0$ can be evolved by
\begin{equation}\label{qgs:eq:lowdynamics}
 \psi(s)=W\expfin(-is\Heff)\psi_0.
\end{equation}
This is an exact evolution inside the invariant graph and, by \cref{qgs:obs:lem:exponential}, has the ordinary shadow
\[
 \st\psi(s)=
 \begin{pmatrix}I\\0\end{pmatrix}
 e^{-is(A_0-B_0D_0^{-1}B_0^*)}\st\psi_0.
\]
\TA\ Equation \eqref{qgs:eq:lowdynamics} uses only the finite $\Heff$. It does not assign a global phase to every eigenvalue of the unlimited high block. Nor does it say that a generic initial state, including arbitrary high-sector components, admits such an evolution in the present formalism.

Time-dependent $A,B,D,\Delta$ would introduce an additional problem: a moving invariant graph is not automatically invariant under time evolution. Derivatives of $W$ generate geometric and off-block terms. The theorem is a static spectral and exact fixed-subspace result; a time-dependent adiabatic theorem would require additional regularity and error bounds.

\section{Support-controlled block diagonalization and slow dynamics}\label{qgs:sec:block}
\status{\TC\ Source 03. The invariant-graph and unitary-normalization mechanisms are classical \cite{BDL}; the collection already studies spectral projectors at infinitesimal gaps \cite{repoQuaternion} and bounds eigenvector rotation at a separated cluster (\texttt{thm:subspace} of \cite{repoSpectral}). Here an exact block is constructed, with a support certificate valid even when the exponent group has no positive order unit.}

\subsection{An isolated residue eigenspace}
Degenerate perturbation theory asks for the dynamics inside a spectral cluster rather than for a termwise correction to an arbitrarily chosen eigenvector. The Schrieffer--Wolff method makes this precise by a unitary transformation that removes the off-diagonal coupling \cite{BDL}.

Fix ordinary finite dimensions $m,r$. Let
\begin{equation}\label{qgs:obs:eq:blockH}
 H=\begin{pmatrix} eI_m+P_{11}&P_{12}\\P_{12}^*&D_0+P_{22}\end{pmatrix},
 \qquad \Gap=D_0-eI_r,
\end{equation}
where $e\in\R$, $D_0=D_0^*\in M_r(\C)$, and the ordinary gap matrix $\Gap$ is invertible. The perturbations $P_{11}=P_{11}^*$, $P_{22}=P_{22}^*$, and $P_{12}$ have entries in $\mm_K$. For an actual low-energy block one can assume $\Gap>0$; invertibility is enough for the algebraic construction. In this section the word ``gap'' means an ordinary nonzero separation in the residue spectrum. A gap that is itself infinitesimal is handled by normalization at a later stage (\cref{qgs:obs:thm:hierarchy}), not by silently treating its inverse as finite; contrast \cref{qgs:sec:elimination}, where the high block is unlimited, and \cref{qgs:sec:schur}, where it becomes soft.

An invariant graph has the form $\{(u,Xu):u\in K^m\}$ with $X\in M_{r\times m}(K)$. Its invariance is equivalent to the matrix Riccati equation
\begin{equation}\label{qgs:obs:eq:riccati}
 P_{12}^*+(\Gap+P_{22})X-XP_{11}-XP_{12}X=0.
\end{equation}
The noncommutative factor order in this equation is essential.

\begin{theorem}[\TC\ Exact Hahn block diagonalization]\label{qgs:obs:thm:block}
For \eqref{qgs:obs:eq:blockH} there is a unique infinitesimal solution $X$ of \eqref{qgs:obs:eq:riccati}. If $\Sigma$ is the union of the positive supports of the perturbation entries, then
\begin{equation}\label{qgs:obs:eq:X-support}
 \supp X\subseteq\Sigma^*\setminus\{0\}.
\end{equation}
Define
\begin{align}
 W&=\binom{I_m}{X}(I_m+X^*X)^{-1/2},\label{qgs:obs:eq:J}\\
 V&=\binom{-X^*}{I_r}(I_r+XX^*)^{-1/2},\qquad
 U=[W\ V].\label{qgs:obs:eq:U}
\end{align}
Then $U$ is unitary, $U-I\in M_{m+r}(\mm_K)$, and
\begin{equation}\label{qgs:obs:eq:unitary-block}
 U^*HU=\diag(\Heff,H_\perp),\qquad \Heff=W^*HW.
\end{equation}
The corrections in $U$, $\Heff-eI$, and $H_\perp-D_0$ have supports in $\Sigma^*\setminus\{0\}$. If the perturbation entries have valuation at least $\alpha>0$, then
\begin{align}
 X&=-\Gap^{-1}P_{12}^*+\Oval{2\alpha},\label{qgs:obs:eq:X-leading}\\
 \Heff&=eI_m+P_{11}-P_{12}\Gap^{-1}P_{12}^*+\Oval{3\alpha}.
 \label{qgs:obs:eq:Heff-leading}
\end{align}
\end{theorem}
\begin{proof}
First regard the real and imaginary parts of the finitely many independent entries of $P_{11},P_{12},P_{22}$ as formal variables of degree one. Rewrite \eqref{qgs:obs:eq:riccati} as
\begin{equation}\label{qgs:obs:eq:graph-recursion}
 X=-\Gap^{-1}P_{12}^*-\Gap^{-1}P_{22}X
       +\Gap^{-1}XP_{11}+\Gap^{-1}XP_{12}X.
\end{equation}
The homogeneous component of degree one is $X_1=-\Gap^{-1}P_{12}^*$. At degree $n\ge2$, the right side depends only on components $X_j$ with $j<n$. Thus there is a unique formal solution with no constant term. Every component is a finite polynomial in the formal entry variables. This construction takes place over ordinary finite matrix spaces; no convergence theorem has been used.

Evaluate the variables at the given Hahn perturbations. Their finitely many supports have well-ordered positive union $\Sigma$. \Cref{qgs:obs:lem:support} proves that the evaluation is strong, commutes with the finite polynomial equation, and has support \eqref{qgs:obs:eq:X-support}. This gives an actual solution in the Hahn field.

For uniqueness, let $X,Y\in M_{r\times m}(\mm_K)$ be two solutions and put $Z=X-Y$. Subtraction gives
\[
 \Gap Z=-P_{22}Z+ZP_{11}+ZP_{12}X+YP_{12}Z.
\]
If $Z\ne0$, multiplication by the ordinary invertible matrix $\Gap$ preserves its valuation: the leading coefficient cannot be annihilated. Every term on the right has valuation strictly larger than $v(Z)$, since it contains $Z$ and at least one positive-valuation factor. This is impossible. Hence $Z=0$.

The inverse square roots in \eqref{qgs:obs:eq:J}--\eqref{qgs:obs:eq:U} can be obtained by the binomial series in $X^*X$ and $XX^*$. Their arguments have positive valuation, so the series are strongly evaluable. They agree with the positive inverse square roots supplied by \cref{qgs:prop:spectral}. Direct multiplication shows $W^*W=I$, $V^*V=I$, and $W^*V=0$. Thus the square matrix $U$ is unitary. It has residue $I$. Equation \eqref{qgs:obs:eq:riccati} says that the image of $W$ is invariant under $H$. Hermiticity then makes its orthogonal complement invariant, proving \eqref{qgs:obs:eq:unitary-block} and Hermiticity of both reduced blocks. All constructions just used are formal power-series expressions with ordinary constant coefficients, proving the remaining support claims.

For the leading terms, the recursion gives $X=X_1+\Oval{2\alpha}$ and $v(X)\ge\alpha$. Write $T=I_m+X^*X$. Graph invariance gives
\begin{equation}\label{qgs:obs:eq:Heff-exact}
 \Heff=T^{1/2}(eI_m+P_{11}+P_{12}X)T^{-1/2}.
\end{equation}
Since $T^{\pm1/2}=I+\Oval{2\alpha}$, conjugating the scalar $eI_m$ has no effect, conjugating $P_{11}$ changes it only at valuation at least $3\alpha$, and conjugating $P_{12}X$ changes it only at valuation at least $4\alpha$. Finally $P_{12}X=-P_{12}\Gap^{-1}P_{12}^*+\Oval{3\alpha}$, proving \eqref{qgs:obs:eq:Heff-leading}.
\end{proof}

\TA\ The similarity expression \eqref{qgs:obs:eq:Heff-exact} need not look Hermitian term by term. Its equality to $W^*HW$ proves exact Hermiticity. Omitting the normalization of the graph would instead produce a matrix representing the same restricted operator in a nonorthonormal basis. Confusing these two matrices is a common source of incorrect higher-order formulas, and the same point appears in \cref{qgs:thm:elimination}, where the normalization produces the anticommutator in \eqref{qgs:eq:Kcorrection}.

\subsection{Why the support statement matters}
\TA\ In the rank-two workspace \eqref{qgs:eq:ranktwo}, a perturbation can have ordinary power corrections in $\epsilon$ and also a separate $\zeta$ contribution that lies beyond all of those powers. A truncation justified only by a finite power count in $\epsilon$ cannot determine the $\zeta$ term. \Cref{qgs:obs:thm:block} does not erase that second scale. Every coefficient comes from a finite set of words in the input supports, even if infinitely many lower-rank scales precede a chosen exponent in the order. The graph is an exact Hahn object, and the error estimate is a separate valuation statement. The method neither assumes a numerical convergence radius nor claims that a finite numerical expansion reveals every scale.

\status{\TA\ The extension proved here is the exact positive-support Hahn evaluation with all-rank uniqueness and explicit support containment. It is not a new Schrieffer--Wolff principle, nor does it by itself establish many-body locality or an error bound uniform in system size.}

\subsection{Inverse-splitting-scale dynamics}
\begin{theorem}[\TC\ Inverse-splitting-scale dynamics]\label{qgs:obs:thm:dynamics}
Suppose an exactly isolated effective block has the form
\begin{equation}\label{qgs:obs:eq:slowH}
 \Heff=eI_m+t^\gamma Z,\qquad \gamma>0,\quad Z=Z^*\in M_m(\OO_K).
\end{equation}
Remove the scalar phase generated by $eI_m$, and use rescaled time $s\,t^{-\gamma}$ with ordinary $s$. The resulting relative propagator is well defined by
\begin{equation}\label{qgs:obs:eq:slow-prop}
 V(s)=\expfin(-isZ).
\end{equation}
For any finite density matrix $\rho$ and effect $E$ in this block,
\begin{equation}\label{qgs:obs:eq:slow-prob}
 \st\tr\bigl(EV(s)\rho V(s)^*\bigr)
 =\tr\bigl((\st E)e^{-isZ_0}(\st\rho)e^{isZ_0}\bigr).
\end{equation}
In contrast, at ordinary unscaled time, the same relative evolution has standard part $I_m$. Thus reduction before rescaling can discard a complete nontrivial effective quantum dynamics.
\end{theorem}
\begin{proof}
After subtracting $eI_m$, multiplication of the generator by the rescaled time $s t^{-\gamma}$ gives the finite matrix $sZ$. Define its evolution by \cref{qgs:obs:lem:exponential}, not by assigning a phase to an arbitrary infinite imaginary argument. That lemma and multiplicativity of standard part give \eqref{qgs:obs:eq:slow-prob}. At an ordinary time $s$, the generator is $s t^\gamma Z\in M_m(\mm_K)$, whose exponential has residue $I_m$.
\end{proof}

\TA\ This generalizes the physics report's example \texttt{phys:ex:amplify} ($H=\epsilon\sigma_z$, $T=\epsilon^{-1}$) to any exactly isolated effective block. Like that example, it demonstrates a mathematical failure of interchange and supplies no physical realization of the unlimited duration. Removing the scalar energy does not require us to define its possibly infinite absolute phase: probabilities inside a single isolated block are invariant under that phase. Coherent interference between different blocks with an unlimited relative phase would be a different question and is not solved by this construction.

\begin{corollary}[\TC\ Insensitivity to a bounded high-block completion]\label{qgs:obs:cor:insensitive}
For the unitary $U$ of \cref{qgs:obs:thm:block}, let $U_\perp$ be any unitary matrix on the complementary block. At a fixed ordinary rescaled time form $\mathcal V=U\diag(V(s),U_\perp)U^*$. For inputs and measurements supported on the original first $m$ coordinate positions, their standard-part statistics depend only on $e^{-isZ_0}$ and not on $U_\perp$.
\end{corollary}
\begin{proof}
Every entry of $U_\perp$ is finite and $\st U=I$. Hence $\st\mathcal V=\diag(e^{-isZ_0},\st U_\perp)$. Multiplication by first-block inputs and effects removes the second block exactly in the reduced formula.
\end{proof}

\TA\ This is an independence statement, not an existence theorem for a particular high-frequency Schr\"odinger evolution. It shows why a choice among bounded unitary high-block phases cannot contaminate this low-block residue calculation.

\subsection{A fully worked mediated-coupling example}\label{qgs:obs:sec:threelevel}
Let $\epsilon=t^\alpha$, $\alpha>0$, $d>0$ ordinary, and $a,b\in\C$ not both zero. Consider
\begin{equation}\label{qgs:obs:eq:threeH}
 H=\begin{pmatrix}
 0&0&\epsilon a\\
 0&0&\epsilon b\\
 \epsilon\bar a&\epsilon\bar b&d
 \end{pmatrix},\qquad
 x=\binom a b,\quad r^2=\abs a^2+\abs b^2.
\end{equation}
The dark low vector is orthogonal to $x$ and has eigenvalue zero. On the span of the bright vector $x/r$ and the third coordinate, the Hamiltonian is $\bigl(\begin{smallmatrix}0&\epsilon r\\\epsilon r&d\end{smallmatrix}\bigr)$. \TE\ Its lower eigenvalue is
\begin{equation}\label{qgs:obs:eq:lambda}
 \lambda_-=\frac{d-\sqrt{d^2+4\epsilon^2r^2}}2
 =-\frac{\epsilon^2r^2}{d}
  +\frac{\epsilon^4r^4}{d^3}
  -\frac{2\epsilon^6r^6}{d^5}+\Oval{8\alpha}.
\end{equation}
The square root is the positive Hahn square root; its binomial expansion is strongly summable because $\epsilon$ has positive valuation. At $r=d=1$ this is the example \texttt{ex:secondorder} of the spectral-theory report ($\lambda_-=-t^2+t^4-2t^6+5t^8+\dots$), which source 03 did not cite.

For the normalized graph identification in \eqref{qgs:obs:eq:J}, the dark vector is fixed and the bright vector maps to the low eigenvector with positive bright component. Therefore the exact effective Hamiltonian is
\begin{equation}\label{qgs:obs:eq:threeheff}
 \Heff=\frac{\lambda_-}{r^2}xx^*
 =-\frac{\epsilon^2}{d}xx^*+\Oval{4\alpha}.
\end{equation}
The sharper even-power remainder here comes from the exact bright/dark solution, not from an unproved improvement to the general theorem. At time $s\epsilon^{-2}$, the reduced generator is $-xx^*/d$. \TE\ For $a=b=1$, the transition between the two low coordinate states has
\begin{equation}\label{qgs:obs:eq:rabi}
 P_{1\to2}^{\rm shadow}(s)=\sin^2(s/d).
\end{equation}
Indeed, $xx^*$ has eigenvalues $0,2$, and $e^{isxx^*/d}=I+(e^{2is/d}-1)xx^*/2$; the off-diagonal squared modulus is \eqref{qgs:obs:eq:rabi}.

\TA\ The full residue of \eqref{qgs:obs:eq:threeH} has a completely degenerate zero low-energy block with no coupling. Nevertheless its infinitesimal virtual coupling produces the nontrivial ordinary oscillation \eqref{qgs:obs:eq:rabi} on the inverse-square scale. The states in the exact low block are the dressed states $W\ket1,W\ket2$; \cref{qgs:obs:cor:insensitive} explains the standard-part independence of a bounded complementary evolution when one uses the bare low coordinates. No infinite-angle global surcomplex exponential was required.

\subsection{An exact finite-depth hierarchy}\label{qgs:obs:sec:hierarchy}
The isolated-residue theorem can be repeated after normalizing an infinitesimal splitting. This leads to a finite tree of scales, even when the entries themselves have transfinite supports.

\begin{theorem}[\TC\ Finite-depth spectral selection]\label{qgs:obs:thm:hierarchy}
Every Hermitian $n\times n$ matrix over $K$ admits a recursive exact unitary block decomposition with the following properties. At each non-scalar block one subtracts a scalar, extracts its leading valuation, diagonalizes an ordinary Hermitian residue matrix, and exactly separates its distinct residue eigenspaces using \cref{qgs:obs:thm:block}. The recursion terminates in scalar blocks. There are at most $n-1$ nontrivial splitting nodes in the tree and at most $n-1$ along any branch. Each split is equipped with its leading gap scale, and each normalized infinitesimal correction has the local support certificate of \cref{qgs:obs:thm:block}. The statement is independent of the rank of $\Gamma$.
\end{theorem}
\begin{proof}
Consider a block $H'$ of dimension $d'$. Put $c=\tr H'/d'$. If $H'=cI$, stop. Otherwise let $\gamma=v(H'-cI)$ and
\[
 Y'=t^{-\gamma}(H'-cI),\qquad Y'_0=\st Y'.
\]
Then $Y'_0$ is a nonzero traceless ordinary Hermitian matrix. It cannot be scalar: a scalar traceless matrix is zero in characteristic zero. Thus it has at least two distinct real eigenvalues. Use an ordinary unitary to write its eigenspaces as blocks. The normalized matrix has that ordinary block diagonal residue and an infinitesimal perturbation.

Separate one residue eigenvalue from all the others by \cref{qgs:obs:thm:block}; its ordinary gap matrix is invertible. Repeat on the remaining complement for each other residue eigenvalue. The resulting unitary makes the decomposition exact. Rescale by $t^\gamma$ and add $cI$; these operations preserve the exact invariant subspaces. Every new block has dimension strictly less than $d'$.

Continue on each block. Dimension decreases by at least one along each branch, so the process terminates. At any splitting node there are at least two children. A finite rooted tree with at most $n$ terminal blocks and all internal nodes having at least two children has at most $n-1$ internal nodes. This proves the bounds. At a split, eigenvalues in distinct residue children differ by $t^\gamma$ times a nonzero ordinary residue difference; their gap valuation in that block is exactly $\gamma$. The local support assertions come directly from the evaluated graph constructions.
\end{proof}

\TA\ Finite depth is not finite computational cost. A graph at one node can contain infinitely many support terms, and deciding whether a block is exactly scalar can be undecidable for a chosen notation language. The theorem is an algebraic construction with access to exact Hahn operations, not an algorithm for arbitrary finite strings purporting to denote surreals. Nor is the final support automatically the monoid of the original unscaled entry supports: division by small gaps can shift exponents. The certified monoid statement is local to each normalized split.

\subsection{Finite thermal states at surreal scales}\label{qgs:obs:sec:gibbs}
The previous results need only a Hahn field and finite-angle analytic lifting. Real thermal exponentials are different: $e^{-x}$ is meaningful for an unlimited positive surreal $x$ using the Gonshor exponential, but it may leave a previously chosen Hahn workspace. We therefore state the next theorem in $\No$ itself, or in an explicitly chosen exponential-closed subfield containing the finite data. Exponential-closed surreal subfields are studied in \cite{BournezGuilmant}; compatibility of the real exponential with finite infinitesimal expansions belongs to the standard surreal exponential theory \cite{Gonshor}, also used in \cite{BM}.

\begin{theorem}[\TC\ Gibbs reduction at multiple energy scales]\label{qgs:obs:thm:gibbs}
Let $H$ be a finite Hermitian matrix over $\No[i]$. Choose rank-one orthogonal spectral projections $P_1,\ldots,P_n$ and real surreal energies $E_1,\ldots,E_n$. Let $\beta>0$, set $E_*=\min_jE_j$, and define
\begin{equation}\label{qgs:obs:eq:gibbs}
 x_j=\beta(E_j-E_*)\ge0,\qquad
 \rho_\beta=\frac{\sum_{j=1}^n e^{-x_j}P_j}
                  {\sum_{j=1}^n e^{-x_j}}.
\end{equation}
Let $\mathcal J$ be the nonempty set of indices for which $x_j$ is finite, and write $a_j=\st x_j$ for $j\in\mathcal J$. Then $\rho_\beta$ is a finite density matrix and
\begin{equation}\label{qgs:obs:eq:gibbs-shadow}
 \st\rho_\beta=
 \frac{\sum_{j\in\mathcal J}e^{-a_j}\st P_j}
      {\sum_{j\in\mathcal J}e^{-a_j}}.
\end{equation}
The expression is independent of choices of rank-one decomposition inside degenerate eigenspaces.
\end{theorem}
\begin{proof}
A normalized eigenvector has finite coordinates because each coordinate has modulus at most one. Hence every $P_j$ is finite. Each weight in \eqref{qgs:obs:eq:gibbs} lies in $(0,1]$, and at least one equals one. The denominator is between $1$ and $n$, so normalization is by a finite unit. Positivity and trace one follow directly from the orthogonal spectral formula.

If $x_j$ is finite, write $x_j=a_j+\varepsilon_j$ with $\varepsilon_j$ infinitesimal. The real exponential addition law and its infinitesimal expansion give $\st e^{-x_j}=e^{-a_j}$. If $x_j$ is positive unlimited, then for every ordinary real $M>0$ one has $x_j>\log M$, so $e^{-x_j}<1/M$; that weight has standard part zero. Reduce the finite sums and the unit denominator to obtain \eqref{qgs:obs:eq:gibbs-shadow}. Equal energies have equal weights, so regrouping their rank-one projections does not change either side.
\end{proof}

\status{\TC\ Source 03. \TI\ In an eigenbasis, the diagonal of $\rho_\beta$ is a finite softmax of separated scales, whose classical content is the finite multiscale softmax limit \texttt{fsp:thm:multisoft} of the probability report \cite{repoProbability}; source 03 did not cite it. The theorem here adds the spectral projections and arbitrary surreal energy differences.}

\begin{example}[\TC\ A three-level thermal scale diagram]\label{qgs:obs:ex:thermal}
Let $H=\diag(0,\epsilon,\zeta)$, where $0<\zeta\ll\epsilon\ll1$. The following are exact reductions.
\begin{center}\small
\begin{tabular}{p{0.43\textwidth}p{0.43\textwidth}}
\toprule
Inverse-temperature regime & Reduced diagonal weights\\
\midrule
$\beta\epsilon\ll1$ & $(1,1,1)/3$\\[3pt]
$\beta=\epsilon^{-1}$ & $(1,e^{-1},1)/(2+e^{-1})$\\[3pt]
$\epsilon^{-1}\ll\beta\ll\zeta^{-1}$ & $(1,0,1)/2$\\[3pt]
$\beta=\zeta^{-1}$ & $(1,0,e^{-1})/(1+e^{-1})$\\[3pt]
$\zeta^{-1}\ll\beta$ & $(1,0,0)$\\
\bottomrule
\end{tabular}
\end{center}
An intermediate regime really exists: $\beta=(\epsilon\zeta)^{-1/2}$ makes $\beta\epsilon$ unlimited and $\beta\zeta$ infinitesimal. The order of the scale limits determines which ordinary effective thermal state is selected.
\end{example}

\TA\ The theorem packages a hierarchy of finite-system low-temperature limits; it does not prove a thermodynamic phase transition. A fixed finite sum has none of the volume-limit questions of an infinite statistical-mechanical system. Its possible utility is to record several competing degeneracy scales in one exact object, including weights beyond all powers of a previously selected parameter.

\section{Virtual selection and the limits of direct truncation}\label{qgs:sec:virtual}

\subsection{A scalar scale-balance law}
The simplest matrix makes the different regimes explicit. Let $a,b\in\Gamma$ with $a>0$ and $a>b$, and let $d>0$ and $c\ne0$ be ordinary real constants. Consider
\[
 H_{a,b}=\begin{pmatrix}0&c t^{-b}\\c t^{-b}&d t^{-a}\end{pmatrix}.
\]
\TE\ Its lower eigenvalue is
\begin{equation}\label{qgs:eq:scalarlambda}
 \lambda_-=
 \frac{d t^{-a}-\sqrt{d^2t^{-2a}+4c^2t^{-2b}}}{2}
 =-\frac{c^2}{d}t^{a-2b}
 \left(1+O\bigl(t^{2(a-b)}\bigr)\right).
\end{equation}
The expansion follows by factoring $d^2t^{-2a}$ from the square root and evaluating the binomial series at the positive-order parameter $4c^2t^{2(a-b)}/d^2$.

\TC\ The high-component to low-component ratio of the low eigenvector is $O(t^{a-b})$, hence infinitesimal. Nevertheless the virtual energy shift is infinitesimal, appreciable, or unlimited according as $a-2b$ is positive, zero, or negative. Thus small mixing alone does not imply a negligible energy correction. The critical case $a=2b$ is exactly the scaling built into \cref{qgs:thm:elimination}. \TA\ This is ordinary second-order perturbative physics written as an exact ordered-group trichotomy. The spectral-theory example \texttt{ex:secondorder} is the neighbouring regime of an ordinary high level with an infinitesimal coupling (\cref{qgs:sec:block}), and the soft counterpart, with a vanishing rather than unlimited eliminated stiffness, is \eqref{qgs:ops:eq:scalar-soft}.

\subsection{An explicit beyond-all-orders selection counterexample}
\begin{proposition}[\TC\ Virtual splitting defeats a smaller direct bias]\label{qgs:prop:virtual}
Let $\epsilon,\zeta$ be as in \eqref{qgs:eq:ranktwo}, and define
\begin{equation}\label{qgs:eq:threelevel}
 H_{\epsilon,\zeta}=
 \begin{pmatrix}
 0&0&\epsilon\\
 0&\zeta&\epsilon\\
 \epsilon&\epsilon&1
 \end{pmatrix}.
\end{equation}
Its ground eigenvalue $\lambda_{\min}$ is simple and satisfies
\begin{equation}\label{qgs:eq:groundenergy}
 \st(\lambda_{\min}/\epsilon^2)=-2.
\end{equation}
If $P_{\min}$ is its ground-state orthogonal projector, then
\begin{equation}\label{qgs:eq:groundstate}
 \st P_{\min}=\proj{+},\qquad
 \ket{+}=\frac{\ket{0}+\ket{1}}{\sqrt2},
\end{equation}
with zero third component. In contrast, the direct compression of $H_{\epsilon,\zeta}$ to $\Span\{\ket{0},\ket{1}\}$ has the unique ground state $\ket{0}$.
\end{proposition}
\begin{proof}
Divide the Hamiltonian by the positive scalar $\epsilon^2$, which does not change its eigenvectors or the ordering of its eigenvalues. It takes the form of \eqref{qgs:eq:criticalH} with
\[
 A=\diag(0,\zeta/\epsilon^2),\quad
 B=\begin{pmatrix}1\\1\end{pmatrix},\quad D=1,\quad\Delta=\epsilon.
\]
The entry $\zeta/\epsilon^2$ is infinitesimal. By \cref{qgs:thm:elimination}, the finite effective matrix has residue
\[
 \st\Heff=-\begin{pmatrix}1&1\\1&1\end{pmatrix}.
\]
This ordinary matrix has the distinct eigenvalues $-2$ and $0$, with the $-2$ eigenspace spanned by $\ket{+}$. Reducing a unitary diagonalization of $\Heff$ shows that exactly one of its eigenvalues has residue $-2$ and the other has residue $0$. The first is therefore the simple smaller eigenvalue. The complementary eigenvalue is positive unlimited in the rescaled matrix, so cannot be the ground eigenvalue. This proves \eqref{qgs:eq:groundenergy}.

A normalized ground eigenvector of $\Heff$ has nonzero normalized residue in the $-2$ eigenspace. Its projector therefore reduces to $\proj{+}$, independently of its phase. Since the dressing matrix $W$ reduces to the coordinate inclusion, the same is true in three dimensions. Finally, direct compression yields $\diag(0,\zeta)$, whose ground vector is $\ket{0}$ because $\zeta>0$.
\end{proof}

\TE\ At $\zeta=0$ the exact eigenvalues are
\[
 0,\qquad \frac{1+\sqrt{1+8\epsilon^2}}2,\qquad
 \frac{1-\sqrt{1+8\epsilon^2}}2,
\]
and the ground energy begins
\[
 -2\epsilon^2+4\epsilon^4-16\epsilon^6+80\epsilon^8+O(\epsilon^{10}).
\]
\TA\ The role of the tiny positive $\zeta$ is not to generate this shift, but to expose a failure of naive scale-by-scale direct restriction. Eliminating the high level first produces a negative $\epsilon^2$ matrix on the low plane; ignoring that virtual contribution and inspecting the direct $\zeta$ bias selects the wrong shadow state. This is the three-level mediated coupling of \cref{qgs:obs:sec:threelevel} with $a=b=1$, $d=1$, plus a direct bias.

This is a useful benchmark for symbolic effective-Hamiltonian algorithms. A method that retains only direct coefficients inside each degenerate subspace will fail it, even if it correctly orders $\zeta$ below every power of $\epsilon$. The relevant data include the paths through the eliminated sector. The mechanism is well known in ordinary degenerate perturbation theory \cite{BDL}; the higher-rank example makes the truncation failure exact and impossible to repair by assigning a sufficiently large but finite power to $\zeta$.

\subsection{Potential model applications and what must be supplied}
\TA\ The same algebra can organize off-resonant virtual couplings, several detuned auxiliary levels, and nested low-energy manifolds. In those settings $\Delta^{-1}D\Delta^{-1}$ represents a high-energy block and $B\Delta^{-1}$ a critically scaled coupling. The theorem supplies the exact finite shift and its first normalization correction once those matrices have been derived.

For tunnelling or metastability, an exponentially small matrix element is usually the output of a separate semiclassical calculation. The present framework can compare it with algebraic virtual shifts without erasing it; it does not derive an instanton action, a tunnelling prefactor, or a resummation prescription. A mathematically honest application must retain that division of labor.

The ordinary counterpart in \cref{qgs:cor:realrealization} also prevents an overly strong claim of uniquely surreal predictions. For finite matrices, the same limits can be established with real-parameter estimates. The gain is a compact exact language and reusable algebraic certificates, not a proof that real or complex mathematics is inadequate.

\section{Complete classification of soft-mode Schur defects}\label{qgs:sec:schur}
\status{\TC\ Source 08; arbitrary real closed $F$ unless a Hahn monomial is displayed. \TI\ The generalized Schur complement and shorted operators are classical \cite{AT,Gallier}. \TA\ The fixed-shadow classification with its rank obstruction and one-parameter realization is proposed as a candidate contribution, not a certified first-in-literature result.}

\subsection{The fixed-shadow problem}
Let
\begin{equation}\label{qgs:ops:eq:block}
 H=\begin{pmatrix}A&B\\B^*&D\end{pmatrix}\ge0,
 \qquad D>0,
\end{equation}
where $n,m$ are positive ordinary integers, $A$ is $n\times n$, $D$ is $m\times m$, and every entry is limited. Its Schur complement is
\begin{equation}\label{qgs:ops:eq:schur}
 S(H)=A-BD^{-1}B^*.
\end{equation}
The inverse $D^{-1}$ need not be limited. Nevertheless $S(H)$ is limited and positive semidefinite, as we will prove. Denote the residue blocks by $A_0,B_0,D_0$ and write $H_0=\st H$.

For a positive ordinary matrix $D_0$, its Moore--Penrose inverse $D_0^\dagger$ acts as the reciprocal on each positive eigenspace and as zero on its kernel. Define the generalized Schur complement
\begin{equation}\label{qgs:ops:eq:generalized-schur}
 S_0=A_0-B_0D_0^\dagger B_0^*,\qquad r=\dim\ker D_0.
\end{equation}
This is the ordinary quadratic form obtained after minimizing over the eliminated coordinates of $H_0$. It is closely related to the shorted operator construction \cite{AT,Gallier}. The question here is:
\begin{quote}
For a fixed ordinary block matrix $H_0\ge0$, which matrices can occur as $\st S(H)$ among all limited positive-semidefinite lifts with $\st H=H_0$ and $D>0$?
\end{quote}
The answer will be a positive interval with an exact rank restriction. Contrast \cref{qgs:sec:elimination}, where $D_0>0$, the eliminated block is unlimited after scaling, and the shadow of the effective matrix is the ordinary Schur complement: there are no soft columns.

\begin{lemma}[\TC\ Completion of the square and whitening]\label{qgs:ops:lem:schur}
For \eqref{qgs:ops:eq:block},
\begin{align}\label{qgs:ops:eq:square}
 \begin{pmatrix}x\\y\end{pmatrix}^*
 H\begin{pmatrix}x\\y\end{pmatrix}
 &=x^*S(H)x
 +(y+D^{-1}B^*x)^*D(y+D^{-1}B^*x).
\end{align}
Consequently $S(H)\ge0$ and
\begin{equation}\label{qgs:ops:eq:whiten}
 C=BD^{-1/2}\quad\Longrightarrow\quad
 CC^*=BD^{-1}B^*\le A.
\end{equation}
Both the whitened coupling $C$ and $S(H)$ are limited.
\end{lemma}
\begin{proof}
Expand the right side of \eqref{qgs:ops:eq:square}; the added quadratic term $x^*BD^{-1}B^*x$ cancels the subtraction in $x^*S(H)x$, and the two cross terms are the off-diagonal terms of the left side. Choosing $y=-D^{-1}B^*x$ proves $S(H)\ge0$. Formula \eqref{qgs:ops:eq:whiten} follows, and $A-CC^*=S(H)\ge0$. Each row squared norm of $C$ is a diagonal entry of $CC^*$ and is bounded by the corresponding limited diagonal entry of $A$. Therefore every entry of $C$ is limited. It follows that $S(H)=A-CC^*$ is limited.
\end{proof}

\TA\ Notice the role of positivity. A merely small off-diagonal matrix is not enough to control $BD^{-1}B^*$ near a singular $D$. Positivity bounds the \emph{whitened coupling} $BD^{-1/2}$, which is the correct invariant quantity.

\begin{lemma}[\TC\ Ordinary degenerate block]\label{qgs:ops:lem:ordinary-schur}
If $H_0\ge0$, then $B_0\ker D_0=0$ and $S_0\ge0$. After an ordinary unitary change in the eliminated coordinates,
\begin{equation}\label{qgs:ops:eq:ordinary-split}
 D_0=\diag(D_h,0_r),\qquad B_0=(B_h\ \ 0),\qquad D_h>0,
\end{equation}
and $S_0=A_0-B_hD_h^{-1}B_h^*$.
\end{lemma}
\begin{proof}
Diagonalize $D_0$. For every zero diagonal entry in a positive matrix, all entries in its row and column vanish, by the two-coordinate inequality $\abs{H_{jk}}^2\le H_{jj}H_{kk}$. This proves the asserted form of $B_0$. Complete the square using the positive block $D_h$ to obtain $S_0\ge0$. The empty $D_h$ case means $S_0=A_0$.
\end{proof}

\subsection{The exact defect formula}
Diagonalize the actual eliminated block, not just its residue:
\begin{equation}\label{qgs:ops:eq:D-diag}
 D=U\diag(d_1,\ldots,d_m)U^*,\qquad d_j>0.
\end{equation}
The unitary $U$ is limited. Split its eigenvalue indices into the \emph{regular} set $\mathcal J_h=\{j:\st d_j>0\}$ and the \emph{soft} set $\mathcal J_s=\{j:\st d_j=0\}$. Reduction of \eqref{qgs:ops:eq:D-diag} shows $\abs{\mathcal J_s}=r$. Put
\begin{equation}\label{qgs:ops:eq:C-columns}
 \widehat B=BU,\qquad
 \widetilde C=\widehat B\diag(d_j^{-1/2}),
\end{equation}
and let $\widetilde C_s$ contain the soft columns. This is $C$ of \cref{qgs:ops:lem:schur} expressed in an eigenbasis and is limited.

\begin{theorem}[\TC\ Soft-mode defect formula]\label{qgs:ops:thm:defect}
Under \eqref{qgs:ops:eq:block}, define
\begin{equation}\label{qgs:ops:eq:Delta}
 \Xi=\st(\widetilde C_s)\,\st(\widetilde C_s)^*.
\end{equation}
Then
\begin{equation}\label{qgs:ops:eq:defect-formula}
 \boxed{\st S(H)=S_0-\Xi,\qquad
 0\le\Xi\le S_0,\qquad \rank\Xi\le r.}
\end{equation}
The matrix $\Xi$ is independent of the chosen diagonalization. Standard part commutes with the generalized Schur reduction exactly when every soft column of the whitened coupling has zero residue.
\end{theorem}
\begin{proof}
The regular columns of $\widetilde C$ have residues $\st(\widehat B_j)/\sqrt{\st d_j}$. The soft columns satisfy $\widehat B_j=\widetilde C_j\sqrt{d_j}$, hence $\st\widehat B_j=0$. Since $\st U$ is an ordinary unitary diagonalizing $D_0$, we obtain
\begin{align*}
 \st(BD^{-1}B^*)
 &=\sum_{j\in\mathcal J_h}\frac{\st\widehat B_j(\st\widehat B_j)^*}{\st d_j}
    +\sum_{j\in\mathcal J_s}\st\widetilde C_j(\st\widetilde C_j)^*\\
 &=B_0D_0^\dagger B_0^*+\Xi.
\end{align*}
Subtract from $A_0$ to prove the equality in \eqref{qgs:ops:eq:defect-formula}. The definition makes $\Xi\ge0$, and $\st S(H)\ge0$ gives $\Xi\le S_0$. Its Gram factor has $r$ columns, so its rank is at most $r$. The equality $\Xi=S_0-\st S(H)$ makes it intrinsic. A positive Gram sum is zero exactly when every column is zero, proving the final criterion.
\end{proof}

\TA\ When $D_0>0$, no soft columns exist and the familiar commuting formula is recovered. At a singular residue, the correct extra datum is not $B_0$, which vanishes on the soft subspace, but the residue of its coupling divided by the \emph{square root} of the soft stiffness. This is a matrix version of the leading-Gram mechanism of \cref{qgs:thm:rare}.

\subsection{Every permitted defect occurs}
\begin{theorem}[\TC\ Complete fixed-shadow classification]\label{qgs:ops:thm:classification}
Fix any ordinary complex Hermitian block matrix
\[
 H_0=\begin{pmatrix}A_0&B_0\\B_0^*&D_0\end{pmatrix}\ge0
\]
of the given finite sizes. With $S_0$ and $r$ as in \eqref{qgs:ops:eq:generalized-schur}, the set of all possible $\st S(H)$ for limited lifts $H\ge0$, $\st H=H_0$, and $D>0$ is exactly
\begin{equation}\label{qgs:ops:eq:classification-set}
 \boxed{\left\{S_0-\Xi:\ 0\le\Xi\le S_0,
                           \ \rank\Xi\le r\right\}.}
\end{equation}
Every member is realized using just one arbitrarily prescribed positive infinitesimal parameter $\varepsilon$ and ordinary complex matrices.
\end{theorem}
\begin{proof}
\Cref{qgs:ops:thm:defect} proves necessity. For sufficiency, perform the ordinary unitary change of eliminated coordinates in \cref{qgs:ops:lem:ordinary-schur}. Given an admissible $\Xi$, an ordinary spectral factorization, padded with zero columns, gives an $n\times r$ matrix $N$ with $NN^*=\Xi$. Choose $\varepsilon>0$ infinitesimal and define
\begin{equation}\label{qgs:ops:eq:realize}
 D=\begin{pmatrix}D_h&0\\0&\varepsilon^2 I_r\end{pmatrix},\qquad
 B=(B_h\ \ \varepsilon N),\qquad A=A_0.
\end{equation}
Then $D>0$, all entries are limited, and the block residue is exactly $H_0$ in these coordinates. Direct calculation gives
\[
 A-BD^{-1}B^*=A_0-B_hD_h^{-1}B_h^*-NN^*
 =S_0-\Xi\ge0.
\]
The completion-of-the-square identity shows that the full $H$ is positive semidefinite. Undo the ordinary unitary coordinate change. This preserves its shadow and the Schur complement and finishes the construction.
\end{proof}

\begin{corollary}[\TC\ Sharp number of hidden directions]\label{qgs:ops:cor:rank-cost}
A Schur defect of rank $k$ requires at least $k$ eliminated modes whose stiffnesses vanish in the ordinary shadow. Conversely, $k$ such modes suffice for every admissible rank-$k$ defect. If $r\ge\rank S_0$, every matrix $E$ with $0\le E\le S_0$ is a possible reduced shadow.
\end{corollary}
\begin{proof}
The first two assertions follow from the rank bound and construction. For the last, $\Xi=S_0-E$ is positive and bounded by $S_0$. Its range lies in the range of $S_0$: if $x\in\ker S_0$, positivity gives $x^*\Xi x=0$, hence $\Xi x=0$ by a positive square root. Thus $\rank\Xi\le\rank S_0\le r$.
\end{proof}

\TA\ This classification supplies both a no-go result and a realization principle. Given only $H_0$, a unique relaxed shadow cannot in general be inferred. But the ambiguity is not arbitrary: it can only \emph{decrease} the generalized Schur complement, remain positive, and use no more rank than the available soft subspace. An indefinite proposed defect or an increase of stiffness cannot be explained by this particular positive soft-mode mechanism. Other mechanisms, nonpositive forms, changed constraints, or different scalings are not excluded.

\begin{example}[\TE\ A fixed shadow with an entire interval of reductions]\label{qgs:ops:ex:interval}
Let $A_0=I_n$, $B_0=0$, and $D_0=0_n$. For every $0\le E\le I_n$, choose $N$ with $NN^*=I_n-E$ and let
\[
 H_E=\begin{pmatrix}I_n&\varepsilon N\\
                    \varepsilon N^*&\varepsilon^2I_n\end{pmatrix}.
\]
All these matrices have the same ordinary shadow $\diag(I_n,0_n)$, but $S(H_E)=E$ exactly. Thus the residue of the full quadratic form alone does not determine even whether its unrestricted relaxation has a positive stiffness in a given direction.
\end{example}

\subsection{Scale thresholds and why positivity matters}
\TE\ For a single soft mode, take ordinary $a,d>0$ and $b\in\C\setminus\{0\}$, and exponents $\alpha,\beta>0$:
\begin{equation}\label{qgs:ops:eq:scalar-soft}
 H=\begin{pmatrix}a&t^\beta b\\t^\beta\overline b&t^{2\alpha}d\end{pmatrix},
 \qquad
 S(H)=a-\frac{\abs b^2}{d}t^{2(\beta-\alpha)}.
\end{equation}
If $\beta>\alpha$, the soft mode leaves no defect. If $\beta=\alpha$, it leaves the ordinary defect $\abs b^2/d$, and positivity holds exactly when $a\ge\abs b^2/d$. If $\beta<\alpha$, the subtracted term is positive unlimited, contradicting $H\ge0$ with limited $a$. \TA\ This gives an exact three-way threshold rather than the unqualified instruction to discard an infinitesimal off-diagonal entry; compare the scalar trichotomy \eqref{qgs:eq:scalarlambda}, where the eliminated level is unlimited instead of soft.

For multiple soft modes, entrywise exponents of $B$ and $D$ in an arbitrary coordinate system need not decide the answer: the soft eigenvectors themselves can depend on the scale. The intrinsic object is the whitened coupling in \cref{qgs:ops:thm:defect}. Its positive Gram sum rules out cancellation between different normalized eliminated modes, while allowing cancellations within the coupling to one mode.

\section{Quadratic physics: relaxation, displacement, and energy windows}\label{qgs:ops:sec:quadratic-physics}

\subsection{What the Schur complement describes}
\TA\ Interpret \eqref{qgs:ops:eq:block} as a dimensionless positive quadratic energy, a finite stiffness matrix, or a finite-mode Euclidean action, with retained coordinate $x$ and eliminated coordinate $y$. \TC\ The function
\begin{equation}\label{qgs:ops:eq:Q}
 \QH(x,y)=\begin{pmatrix}x\\y\end{pmatrix}^*
             H\begin{pmatrix}x\\y\end{pmatrix}
\end{equation}
has the unique minimizing $y$ at fixed $x$, $y_*=-D^{-1}B^*x$, and its minimum is $x^*S(H)x$. This is an exact finite algebraic statement. It does not require a Gaussian measure. An ordinary finite Gaussian integral, when separately defined over real variables, also has a determinant normalization and other hypotheses; no Hahn-valued integration rule is inferred here.

There is an important physical qualification. The minimizer $y_*$ may be unlimited. A harmonic approximation derived only for bounded amplitudes may then be invalid at its formal minimum. The next proposition identifies this obstruction without an appeal to any particular material or field model.

\begin{proposition}[\TC\ Bounded-displacement obstruction]\label{qgs:ops:prop:displacement}
Let $x\in\C^n$ be an ordinary retained vector, and let $\Xi$ be the defect of $H$ in \cref{qgs:ops:thm:defect}. For every limited $y\in K^m$,
\begin{equation}\label{qgs:ops:eq:bounded-displacement}
 \st\QH(x,y)\ge x^*S_0x.
\end{equation}
Equality is attained by a constant ordinary choice of $y$. In contrast,
\begin{equation}\label{qgs:ops:eq:unbounded-displacement}
 \st\min_{y\in K^m}\QH(x,y)
 =x^*(S_0-\Xi)x.
\end{equation}
Consequently, if $x^*\Xi x>0$, the unrestricted minimizer is unlimited, and no limited displacement attains, even to infinitesimal energy error, the unrestricted minimum. If the unrestricted minimizer is limited, then $x^*\Xi x=0$.
\end{proposition}
\begin{proof}
For limited $y$, reduction of the finite quadratic expression gives $\st\QH(x,y)=\mathcal Q_{H_0}(x,\st y)$. Ordinary minimization of this positive quadratic form over its regular eliminated subspace yields the lower bound $x^*S_0x$, attained at $y_0=-D_0^\dagger B_0^*x$, with a zero kernel component. Taking $y=y_0$ as a constant achieves equality in \eqref{qgs:ops:eq:bounded-displacement}. Equation \eqref{qgs:ops:eq:unbounded-displacement} follows from completion of the square and \cref{qgs:ops:thm:defect}. A positive ordinary difference $x^*\Xi x$ separates the two minimum residues. Thus the unrestricted minimizer cannot be limited, nor can a limited $y$ have an infinitesimal gap above its energy. The last assertion is the contrapositive.
\end{proof}

In the construction \eqref{qgs:ops:eq:realize}, the soft part of the minimizer is especially transparent:
\begin{equation}\label{qgs:ops:eq:soft-displacement}
 y_{*,s}=-\varepsilon^{-1}N^*x.
\end{equation}
\TA\ Its divergence is the mechanism by which the weak coupling $\varepsilon N$ leaves a finite energy correction. A meaningful physical application must either allow these displacements, rescale the soft coordinate before constructing the model, or include nonlinear terms that control them. The theorem does not justify extending a small-displacement theory outside its regime.

\subsection{Static reduction is not a fixed-energy effective Hamiltonian}
For a real spectral parameter $z\in F$ with $D-zI$ invertible, define
\begin{equation}\label{qgs:ops:eq:feshbach}
 F_H(z)=A-zI-B(D-zI)^{-1}B^*.
\end{equation}
This is an energy-dependent Feshbach--Schur pencil, not the single matrix $S(H)$. The distinction is classical \cite{DSS}, and it is especially important when an eliminated gap is infinitesimal.

\begin{proposition}[\TC\ Finite isospectral reduction]\label{qgs:ops:prop:feshbach}
Whenever $D-zI$ is invertible,
\begin{equation}\label{qgs:ops:eq:det-feshbach}
 \det(H-zI)=\det(D-zI)\det F_H(z).
\end{equation}
The map
\begin{equation}\label{qgs:ops:eq:kernel-feshbach}
 x\longmapsto\begin{pmatrix}x\\-(D-zI)^{-1}B^*x\end{pmatrix}
\end{equation}
is a bijection from $\ker F_H(z)$ to $\ker(H-zI)$.
\end{proposition}
\begin{proof}
Block elimination with diagonal blocks $I$ and $D-zI$ gives the determinant formula. For the kernel statement, the lower block equation uniquely determines the displayed $y$ from $x$, and substitution into the upper equation gives exactly $F_H(z)x=0$. A kernel vector with $x=0$ has $y=0$ because $D-zI$ is invertible, so no eigenvector is lost.
\end{proof}

\status{\TI\ The determinant identity and the bound $\norm{E_{12}}_\op^2/\delta$ for a shifted inverse are already \texttt{prop:effective} of the spectral-theory report \cite{repoSpectral}, which source 08 credited; the kernel bijection is classical \cite{DSS}.}

Let $\delta_D>0$ be the smallest eigenvalue of $D$, and use the bounded whitened coupling $C=BD^{-1/2}$.

\begin{theorem}[\TC\ A gap-relative static window]\label{qgs:ops:thm:window}
For real $z$ with $\abs z<\delta_D$,
\begin{equation}\label{qgs:ops:eq:window-bound}
 \norm{F_H(z)-(S(H)-zI)}_\op
 \le \norm C_\op^2\,\frac{\abs z}{\delta_D-\abs z}.
\end{equation}
In particular, if $\abs z/\delta_D$ is infinitesimal, the two sides inside the norm have identical residues. The condition is relative to the actual eliminated gap, not to a fixed real energy unit.
\end{theorem}
\begin{proof}
Since $D$ commutes with $D-zI$, finite spectral calculus gives
\[
 B(D-zI)^{-1}B^*-BD^{-1}B^*
 =C\,[z(D-zI)^{-1}]\,C^*.
\]
Diagonalizing $D$ shows $\norm{(D-zI)^{-1}}_\op\le1/(\delta_D-\abs z)$. Submultiplicativity of the finite operator norm now gives the bound. If $\abs z/\delta_D$ is infinitesimal, then $\abs z/(\delta_D-\abs z)$ is infinitesimal, while $\norm C_\op$ is limited. Every entry of the difference is therefore infinitesimal.
\end{proof}

\TA\ No assertion is made for $z$ at or near an eigenvalue of $D$. That is precisely where the inverse in the reduced pencil becomes singular. An eliminated \emph{soft} mode is not a heavy mode that can be discarded at every fixed energy. \merge\ Relative to \texttt{prop:effective}, which bounds the whole correction $B(D-zI)^{-1}B^*$ by $\norm{B}_\op^2/\operatorname{dist}(z,\operatorname{spec}D)$, the window bounds only its \emph{difference} from the static term, with the constant $\norm C_\op^2=\norm{BD^{-1}B^*}_\op$; by \cref{qgs:ops:lem:schur} this constant is limited even when $B$ and the gap are both infinitesimal.

\begin{example}[\TE\ Three distinct energy regimes]\label{qgs:ops:ex:three-regimes}
Let $a>b^2>0$ be ordinary reals and let $d>0$ be infinitesimal. Set
\begin{equation}\label{qgs:ops:eq:two-level-soft}
 H=\begin{pmatrix}a&b\sqrt d\\b\sqrt d&d\end{pmatrix},\qquad
 F_H(z)=a-z-\frac{b^2d}{d-z}.
\end{equation}
At zero energy the reduced shadow is $a-b^2$. At a fixed ordinary $z\ne0$ it is $a-z$. On the transitional scale $z=\xi d$, with ordinary $\xi\ne1$, it is
\begin{equation}\label{qgs:ops:eq:transition}
 \st F_H(\xi d)=a-\frac{b^2}{1-\xi}.
\end{equation}
At $\xi=1$ the displayed Feshbach chart is not defined. These are not conflicting approximations: they evaluate different ratios $z/d$.

For comparison, the two eigenvalues of $H$ have sum $a+d$ and product $d(a-b^2)$. Reduction of their positive spectral decomposition gives one eigenvalue with shadow $a$ and one with shadow zero. If they are called $\lambda_h,\lambda_s$ accordingly, then
\begin{equation}\label{qgs:ops:eq:soft-eigenvalue}
 \st\lambda_h=a,\qquad
 \st\left(\frac{\lambda_s}{d}\right)=\frac{a-b^2}{a}.
\end{equation}
The second equality follows from the product formula and the nonzero residue of $\lambda_h$. In particular, the number $a-b^2$ is a static relaxed stiffness, not the shadow of the hard eigenvalue.
\end{example}

\subsection{Possible uses and their limits}
\TA\ The defect classification can organize finite-mode models with nearly flat directions, such as coupled harmonic coordinates or a truncated positive Euclidean fluctuation action. Given a measured or computed ordinary block shadow, it describes exactly how much its unrestricted relaxed energy can depend on invisible couplings. The rank bound gives a lower bound on the number of soft directions required by a proposed correction, while the realization theorem supplies a minimal algebraic model. The spectral-window theorem then tests whether a static approximation is legitimate at the energy being studied.

These applications are conditional on positivity, finite truncation, and an admissible displacement regime. Gauge constraints must be dealt with before imposing $D>0$; an exact gauge zero mode is not made invertible by this theorem. Minkowski-signature actions, non-Hermitian open-system generators, or unbounded continuum operators need different hypotheses. An ultraviolet cutoff cannot be removed merely because the remaining coefficients have surreal names.

\section{Surquaternions: useful representation and a composition obstruction}\label{qgs:sec:quaternions}
\status{\TC\ Sources 03, 04 and 07; every statement holds over an arbitrary real closed $F\supseteq\R$ (source 04's generality), except the rare-branch clause, which uses the Hahn setting of \cref{qgs:thm:rare}. \TI\ The algebra is classical, and the collection's surquaternion report already develops it \cite{repoQuaternion}; complex simulation is due to Graydon \cite{Graydon}.}

\subsection{The division algebra and its complex representation}
Define
\[
 \HF=F\oplus F\qi\oplus F\qj\oplus F\qk,
 \qquad \qi^2=\qj^2=\qk^2=\qi\qj\qk=-1.
\]
This is the surquaternion algebra after embedding $F$ in $\No$. Conjugation changes the sign of the three imaginary components, and
\[
 N(q)=\bar q q=q\bar q=a^2+b^2+c^2+d^2,\qquad \abs q=\sqrt{N(q)},
\]
for $q=a+b\qi+c\qj+d\qk$. The norm is multiplicative. If $q\ne0$, then $N(q)>0$ and $q^{-1}=\bar q/N(q)$. In particular, adding surreal infinitesimals does not introduce zero divisors into this quaternion algebra. Imaginary quaternions identify with $F^3$ and satisfy
\begin{equation}\label{qgs:hol:eq:vectorproduct}
 uv=-u\cdot v+u\times v.
\end{equation}
The valuation of a quaternion is the minimum valuation of its four coordinates; positivity of the sum of squares implies $v(N(q))=2v(q)$ and $v(qr)=v(q)+v(r)$.

Identify $K$ with $F+F\qi$, the quaternion $\qi$ with the complex unit $i$, and write $q=z+w\qj$ with $z,w\in K$. Define
\begin{equation}\label{qgs:eq:chi}
 \chi(q)=\begin{pmatrix}z&w\\-\bar w&\bar z\end{pmatrix},
 \qquad\text{so}\qquad
 \chi(s+x\qi+y\qj+z'\qk)=\begin{pmatrix}s+ix&y+iz'\\-y+iz'&s-ix\end{pmatrix}.
\end{equation}
For a quaternion matrix $A=Z_1+Z_2\qj$, with $Z_1,Z_2$ complex matrices, use the corresponding block matrix
\[
 \chi(A)=\begin{pmatrix}Z_1&Z_2\\-\bar Z_2&\bar Z_1\end{pmatrix},
\]
with the corresponding block sizes for rectangular matrices.

\begin{proposition}[\TC\ Faithful ordered $*$-representation]\label{qgs:prop:chi}
The map $\chi$ is injective, real-linear, multiplicative on compatible rectangular matrices, and satisfies $\chi(A^*)=\chi(A)^*$. For a scalar quaternion,
\[
 \chi(q)^*\chi(q)=N(q)I_2,\qquad
 \det\chi(q)=N(q),\qquad
 \tfrac12\tr\chi(q)=\Rea q.
\]
It identifies the norm-one quaternion group $\Sp(1,F)$ with $\SU(2,K)$. For a square quaternion matrix,
\begin{equation}\label{qgs:eq:quattrace}
 \tr_K\chi(A)=2\Rea\tr_{\HH}(A).
\end{equation}
For a quaternion Hermitian matrix $A$, positivity on quaternion columns is equivalent to complex positivity of $\chi(A)$.
\end{proposition}
\begin{proof}
Use $\qj z=\bar z\qj$. Multiplication gives
\[
 (z+w\qj)(z'+w'\qj)=(zz'-w\bar w')+(zw'+w\bar z')\qj,
\]
which agrees with block multiplication in \eqref{qgs:eq:chi}; the same calculation with finite matrix sums proves multiplicativity. Quaternionic conjugation sends $z+w\qj$ to $\bar z-w\qj$, yielding the adjoint identity. Injectivity, the trace identity, and the scalar norm and determinant identities are direct calculations from the blocks. Conversely, a two-by-two unitary matrix of determinant one has the displayed quaternionic form: comparison of its inverse computed by its adjoint and by its determinant gives the required entries. No assertion about a topological covering space is needed.

For positivity, identify a quaternionic column $x=u+u'\qj$ with the complex column $\phi(x)=(u,-\bar u')^T$. This is a bijection and satisfies $\phi(Ax)=\chi(A)\phi(x)$. For Hermitian $A$, direct expansion identifies the real quantity $x^*Ax$ with $\phi(x)^*\chi(A)\phi(x)$. Thus the two positivity tests quantify over corresponding sets of vectors and are equivalent.
\end{proof}

\begin{lemma}[\TC\ Positive quaternionic Gram factorization]\label{qgs:hol:lem:qgram}
A finite quaternionic Hermitian matrix which is nonnegative on every quaternionic vector has a factorization $A=B^*B$. Its complex block representation is positive semidefinite.
\end{lemma}
\begin{proof}
If every diagonal entry is zero, positivity on vectors supported in two coordinates forces every off-diagonal entry to vanish: otherwise a quaternionic coefficient can make the cross term negative. If a positive diagonal entry exists, permute it into the first position and write
\[
 A=\begin{pmatrix}a&x^*\\x&C'\end{pmatrix},\qquad a>0.
\]
Completion of the square, with real central $a$, shows that $C'-xa^{-1}x^*$ is positive semidefinite. Induction factors this Schur complement and then the whole matrix using $\sqrt a$. Applying the multiplicative adjoint-preserving block representation gives $\chi(A)=\chi(B)^*\chi(B)$.
\end{proof}

\TA\ These algebraic facts are classical (sources 03, 04 and 07 each prove the parts they use). Their use over $F$ makes the scale dependence exact but does not create a new gauge group or a new type of finite quaternionic measurement.

\subsection{Finite quaternionic instruments have complex shadows}
For a quaternionic matrix $A$ write $\tr_\R A=\Rea\sum_jA_{jj}$. This real trace is cyclic, $\tr_\R(AB)=\tr_\R(BA)$, although quaternion-valued matrix traces are not generally cyclic. A quaternionic state is a positive Hermitian matrix $\rho$ on right quaternionic vectors with real quaternionic trace one; effects and Kraus maps $X\mapsto\sum_jM_jXM_j^*$ are defined analogously, with the order of multiplication retained. Probabilities are
\begin{equation}\label{qgs:eq:quatborn}
 p(E)=\Rea\tr_{\HH}(\rho E)=\tr_\R(\rho E).
\end{equation}

\begin{theorem}[\TC\ Complex simulation of finite surquaternion instruments]\label{qgs:thm:quatsimulation}
Let a fixed finite quaternionic protocol on $\HF^n$ use states, effects, and Kraus maps with the usual completeness condition. The assignment
\begin{equation}\label{qgs:eq:quatsim}
 \rho\longmapsto\frac{\chi(\rho)}2,
 \qquad E\longmapsto\chi(E),
 \qquad M_j\longmapsto\chi(M_j)
\end{equation}
gives a valid complex state, effect and channel on $K^{2n}$, and is an exact probability-preserving simulation over $K$ on twice the dimension:
\begin{equation}\label{qgs:hol:eq:qprob}
 \tr\Bigl(\frac{\chi(\rho)}2\,\chi(E)\Bigr)=\tr_\R(\rho E).
\end{equation}
It preserves branch probabilities and normalized postselected states of every finite protocol. All operational matrices are finite, and the construction commutes with standard part; standard part therefore yields an ordinary complex simulation of their unnormalized outcome statistics and of all appreciably probable conditioned outcomes. In the Hahn setting, infinitesimally probable branches are handled by \cref{qgs:thm:rare} applied to the complex simulation.
\end{theorem}
\begin{proof}
\Cref{qgs:prop:chi,qgs:hol:lem:qgram} give $\chi(\rho)/2\ge0$ and $0\le\chi(E)\le I$ (apply the lemma to $I-E$), and multiplicativity and the adjoint rule carry $\sum_jM_j^*M_j=I$ to the complex completeness equation. Equation \eqref{qgs:eq:quattrace} shows that $\chi(\rho)/2$ has trace one and gives \eqref{qgs:hol:eq:qprob}. Multiplicativity and the adjoint rule intertwine each Kraus branch,
\[
 \frac12\chi\Bigl(\sum_jM_j\rho M_j^*\Bigr)=\sum_j\chi(M_j)\frac{\chi(\rho)}2\chi(M_j)^*,
\]
so the factor $1/2$ makes every output of a Kraus branch exactly the image of its quaternionic output. A nonzero branch probability is a central real scalar, so its division commutes with $\chi$ as well. Iterating proves equality for finite protocols. The complex states and Kraus operators are finite by \cref{qgs:lem:finitequantum}; their blocks then show finiteness of the quaternionic coordinates. Standard part commutes with $\chi$, products, and adjoints, since $\chi$ consists of finite real and imaginary coordinate operations. The ordinary and rare-conditioning conclusions follow from \cref{qgs:prop:shadow,qgs:thm:rare}.
\end{proof}

\status{\TC\ Proved in sources 03 (Theorem 8.2 there), 04 (Theorem 10.2, over any real closed field) and 07 (Theorem 9.2). \TI\ The ordinary complex-simulation phenomenon is Graydon's \cite{Graydon}; the theorem is its direct ordered-field version, not a new simulation principle. \TA\ Its point here is compatibility with valuation and with the rare-branch formula.}

\TA\ The simulation embeds a restricted family of doubled complex states and transformations; its image is a structured subset of complex matrices, not the entire complex theory, and it does not identify every doubled complex experiment with a quaternionic one. A quaternionic pure state generally maps to a complex state of rank two. The theorem does \emph{not} assert that quaternionic and complex quantum mechanics are equivalent under every possible formulation, and a doubled-dimensional simulation of a whole finite system need not preserve a stipulated factorization into spatially separated subsystems or a locality constraint; naive tensor-product composition cannot be ignored (\cref{qgs:thm:quatcomposite}). Published reduction theorems based on relativistic symmetry make additional hypotheses: Moretti and Oppio obtain a complex structure for suitable quaternionic elementary relativistic systems \cite{MorettiOppio}. That is distinct from the elementary representation argument just proved, and neither result licenses an unrestricted claim that every quaternionic proposal is physically redundant.

What is excluded here is narrower and useful: replacing complex scalars by surquaternionic scalars inside the specified finite whole-system instrument model is not by itself a proof of statistics beyond finite complex models. To obtain a substantive physical distinction, one must identify an additional axiom, composition rule, field equation, or admissible operation that changes the model.

\subsection{Why the naive tensor-dimension prescription fails}
\begin{lemma}[\TC\ Dimension of a quaternionic Hermitian space]\label{qgs:lem:quatdim}
The $F$-dimension of $\Herm_n(\HF)$ is
\[
 d_n=n+4\binom n2=2n^2-n.
\]
Density matrices span this vector space, and effects span its dual through the nondegenerate pairing $\tr_\R(AE)$.
\end{lemma}
\begin{proof}
Diagonal entries are real and contribute $n$ coordinates. Each strictly upper-triangular entry is an arbitrary quaternion and determines its lower-triangular conjugate, contributing four coordinates. For spanning, use the maximally mixed state together with sufficiently small real perturbations in a fixed ordinary Hermitian basis; they are positive and span the space (equivalently, for any Hermitian $A$ choose a positive scalar $r$ with $A+rI\ge0$ and normalize by the positive real trace). Likewise, small perturbations of $I/2$ in that basis are effects and span. The real trace pairing is nondegenerate because $\tr_\R(A^2)>0$ for a nonzero Hermitian $A$, as is immediate from its coordinates or from $\chi$.
\end{proof}

\begin{theorem}[\TC\ Obstruction to a naive quaternionic composite]\label{qgs:thm:quatcomposite}
Let $m,n>1$. A theory cannot simultaneously satisfy all of the following:
\begin{enumerate}[label=\textup{(\alph*)},leftmargin=2em,itemsep=2pt]
\item each local system has the full states and effects of $\HF^m$ and $\HF^n$;
\item every pair of local preparations and every pair of local effects has an independently implementable joint preparation or effect, with factorized product probabilities;
\item the joint system's states and effects are represented in $\Herm_{mn}(\HF)$ with the usual real-trace probability pairing.
\end{enumerate}
This obstruction survives replacement of real scalars by surreal scalars.
\end{theorem}
\begin{proof}[First route: the evaluation matrix (source 07)]
Choose $d_m$ linearly independent local states and $d_m$ local effects whose evaluation matrix on those states is invertible. Such choices exist by \cref{qgs:lem:quatdim}; do the same for the second system. The joint evaluation matrix of the $d_md_n$ product preparations against the $d_md_n$ product effects is the Kronecker product of the two invertible local matrices. It is therefore invertible. Consequently, those joint preparations must be linearly independent in the joint Hermitian vector space, forcing $d_md_n\le d_{mn}$. But direct calculation gives
\begin{equation}\label{qgs:eq:quatdimensiongap}
 d_m d_n-d_{mn}=(2m^2-m)(2n^2-n)-(2m^2n^2-mn)=2mn(m-1)(n-1)>0,
\end{equation}
a contradiction.
\end{proof}

Source 04 proves a variant with a bilinear hypothesis in place of (b), by a different argument; it is kept as a second route.

\begin{theorem}[\TC\ Independent preparation obstructs the naive quaternionic size; second route, source 04]\label{qgs:hol:thm:composite}
Let $n,m>1$. There is no model with composite vector space $\Herm_{nm}(\HF)$ satisfying all of the following requirements:
\begin{enumerate}[label=\textup{(\arabic*)},leftmargin=2em,itemsep=2pt]
\item product preparation extends to an $F$-bilinear map $\mathcal B:\Herm_n(\HF)\times\Herm_m(\HF)\to\Herm_{nm}(\HF)$;
\item every pair of local effects $E,E'$ has a joint effect $G_{E,E'}$;
\item for all local normalized states, $\tr_\R\bigl(G_{E,E'}\mathcal B(\rho,\sigma)\bigr)=\tr_\R(E\rho)\tr_\R(E'\sigma)$.
\end{enumerate}
In particular, simply replacing complex amplitudes by quaternionic amplitudes while retaining the usual product dimension and independent local preparation and measurement is inconsistent with these assumptions.
\end{theorem}
\begin{proof}
Normalized positive states span each Hermitian space over $F$, and effects span as well (\cref{qgs:lem:quatdim}). The bilinear map induces a linear map
\[
 \Lambda:\Herm_n(\HF)\otimes_F\Herm_m(\HF)\longrightarrow\Herm_{nm}(\HF).
\]
The factorization identity extends from states to all Hermitian pairs by linearity. Hence, if $\Lambda(z)=0$, every tensor product of local effect functionals annihilates $z$. The real-trace pairing is nondegenerate, and since effects span, their functionals span the dual; products of such functionals span the dual tensor product, forcing $z=0$. Therefore $\Lambda$ must be injective, contradicting \eqref{qgs:eq:quatdimensiongap}.
\end{proof}

\TE\ For two quaternionic two-level systems, the local dimensions require $6\cdot6=36$ independent product directions, whereas a quaternionic four-level Hermitian space has dimension $28$. \TA\ The first argument does not mistakenly infer a contradiction merely from a lack of local tomography: an injective map from a 28-dimensional space into 36 statistics would be possible. The contradiction instead comes from \emph{independent preparations with factorized product effects}, which force the 36-dimensional lower bound; the second proof does not even need the assumption that local measurements determine every composite state. Larger composite spaces, restrictions on local effects or preparations, or different probabilistic postulates are not excluded. The theorems only reject the naive prescription under their explicit independence assumptions. They are consistent with the broader compositional issues studied in \cite{BarnumWilce}, and they do not exclude the more subtle Jordan-algebraic composites of \cite{BGW}, which may sacrifice familiar dimension or purity rules. They explain why a successful single-system complex simulation should not be advertised as a complete local tensor-product theory. Extending the center from $\R$ to a surreal field leaves the dimension contradiction unchanged.

\subsection{Why gauge theory is a more natural use}
\TA\ The imaginary quaternions already form the Lie algebra of $\SU(2)$ under commutator, and unit quaternions provide its group representation. Their noncommutativity is therefore useful without postulating a new probability rule. Allowing central Hahn coefficients provides multiscale gauge fields with exactly the same bracket identities. The remaining sections exploit that structure while retaining ordinary spacetime, a fixed bundle, and an ordinary geometric Hodge operator, and while keeping the quantum probes explicitly surcomplex: surquaternions are used only to represent $\SU(2)$ transformations, so the composition issue above does not arise. We are \emph{not} replacing complex quantum amplitudes by quaternionic amplitudes or claiming a new tensor product.

\section{Small non-Abelian loops: commutators, invariants, and Wilson actions}\label{qgs:sec:holonomy}
\status{\TC\ Sources 04, 07 and 08, over an arbitrary real closed $F$ except where Hahn evaluation of a formal series is used (\cref{qgs:ops:thm:invariant,qgs:prop:plaquette}). \TI\ Group-valued links and Wilson traces are classical \cite{Wilson}; the leading group commutator is already \texttt{squat:eq:groupcomm} \cite{repoQuaternion}.}

\subsection{An exact commutator identity}\label{qgs:sec:commutator}
\begin{theorem}[\TC\ Commutator norm and Wilson trace]\label{qgs:ops:thm:commutator}
Let $q_1=s_1+u_1$ and $q_2=s_2+u_2$ be unit quaternions over $F$, with scalar parts $s_i$ and imaginary parts $u_i$, and let $c=q_1q_2q_1^{-1}q_2^{-1}$. Then
\begin{equation}\label{qgs:ops:eq:commutator-identity}
 N(c-1)=4\abs{u_1\times u_2}^2,
 \qquad 1-\Rea c=2\abs{u_1\times u_2}^2,
 \qquad \abs{\Ima c}^2=2w-w^2,
\end{equation}
where $w=1-\Rea c=1-\tfrac12\tr\chi(c)$ is the normalized Wilson defect. In particular $c=1$ if and only if $u_1\times u_2=0$. If $u_1\times u_2\ne0$ and $\kappa=v(u_1\times u_2)$, then
\begin{equation}\label{qgs:ops:eq:loop-valuations}
 v(c-1)=\kappa,\qquad v(1-\Rea c)=2\kappa.
\end{equation}
\end{theorem}
\begin{proof}
Quaternion multiplication gives $q_1q_2-q_2q_1=u_1u_2-u_2u_1=2u_1\times u_2$ by \eqref{qgs:hol:eq:vectorproduct}. Keeping factor order, one has
\[
 c-1=(q_1q_2-q_2q_1)q_1^{-1}q_2^{-1}.
\]
The inverses of $q_1,q_2$ have norm one. Taking multiplicative norms proves the first identity. Since $N(c)=1$,
\[
 N(c-1)=(c-1)(\overline c-1)=2-c-\overline c=2(1-\Rea c),
\]
which proves the second, and the unit-norm identity gives $\abs{\Ima c}^2=1-(\Rea c)^2=2w-w^2$. A sum of squares vanishes only when every coordinate vanishes; this proves the exact identity criterion. Positivity of the coordinate norm gives the valuation statements, and \cref{qgs:prop:chi} identifies the normalized trace.
\end{proof}

The affine chart used by source 04 makes the dependence on two control vectors explicit. For an imaginary quaternion $u$, define the \emph{normalized affine chart}
\begin{equation}\label{qgs:hol:eq:Q}
 Q(u)=\frac{1+u}{\sqrt{1+\norm u^2}}.
\end{equation}
This is not the Cayley transform: the square-root normalization is intentional. If $u$ is infinitesimal, $Q(u)$ is a unit quaternion with standard part $1$. Its inverse is obtained by replacing $u$ with $-u$.

\begin{corollary}[\TC\ Wilson defect of a commutator in the affine chart]\label{qgs:hol:cor:comm}\label{qgs:hol:thm:comm}
Let $u_1,u_2\in\Ima\HF$, $a=Q(u_1)$, $b=Q(u_2)$, $c=aba^{-1}b^{-1}$ and $w(c)=1-\Rea c$. Then
\begin{align}
 w(c)&=\frac{2\norm{u_1\times u_2}^2}
 {(1+\norm{u_1}^2)(1+\norm{u_2}^2)},\label{qgs:hol:eq:comm-w}\\
 \norm{c-1}^2&=2w(c),\qquad \norm{\Ima c}^2=2w(c)-w(c)^2.\label{qgs:hol:eq:comm-norm}
\end{align}
In particular $c=1$ if and only if $u_1\times u_2=0$. If $u_1,u_2$ are infinitesimal and $u_1\times u_2\neq0$, then
\begin{equation}\label{qgs:hol:eq:comm-val}
 v(c-1)=v(u_1\times u_2)=\kappa,
 \qquad v(w(c))=2\kappa,
 \qquad w(c)\sim2\norm{u_1\times u_2}^2.
\end{equation}
\end{corollary}
\begin{proof}
The imaginary part of $Q(u_i)$ is $u_i/\sqrt{1+\norm{u_i}^2}$; apply \cref{qgs:ops:thm:commutator}. When $u_1,u_2$ are infinitesimal, the denominator in \eqref{qgs:hol:eq:comm-w} belongs to $1+\mm$; apply \eqref{qgs:hol:eq:normval}.
\end{proof}

\status{\TC\ Source 04 proves \cref{qgs:hol:cor:comm} directly (its Theorem 4.1, by computing $ab-ba$ in the chart), source 08 proves \cref{qgs:ops:thm:commutator} for arbitrary unit quaternions; the proofs are the same factor-ordered identity. \TA\ The exact quaternion identity is elementary and is not advanced as a new quaternion theorem. Its role is to eliminate truncation error from the resource calculations below. In particular, it does not discard a higher-order cross product when the leading directions happen to be parallel.}

\begin{example}[\TE\ Two rotation scales and an axis scale]\label{qgs:ops:ex:axes}
Let $n_1,n_2$ be unit imaginary quaternions, and let $a,b>0$ be infinitesimal. Set
\begin{equation}\label{qgs:ops:eq:axis-rotations}
 q_a=Q(an_1)=\frac{1+an_1}{\sqrt{1+a^2}},\qquad
 q_b=Q(bn_2)=\frac{1+bn_2}{\sqrt{1+b^2}}.
\end{equation}
Their commutator satisfies the exact formula
\begin{equation}\label{qgs:ops:eq:axis-trace}
 1-\Rea c=
 \frac{2a^2b^2\abs{n_1\times n_2}^2}{(1+a^2)(1+b^2)}.
\end{equation}
If $n_1\times n_2\ne0$, its valuation is twice $v(a)+v(b)+v(n_1\times n_2)$. Thus near-parallel axes can create an additional scale, independent of the two rotation strengths. For ordinary perpendicular axes $n_1=\qi$, $n_2=\qj$, and $a=t^\alpha,b=t^\beta$,
\begin{equation}\label{qgs:ops:eq:leading-comm}
 \st\!\left(t^{-(\alpha+\beta)}(c-1)\right)=2\qk,
 \qquad v(1-\Rea c)=2(\alpha+\beta).
\end{equation}
The leading identity follows either by expanding the four finite factors or by the factor-ordered identity in the proof of \cref{qgs:ops:thm:commutator}, since $q_a^{-1}q_b^{-1}$ has residue one.
\end{example}

\TA\ This finite identity explains an important difference between matrix information and a scalar traced loop observable. A commutator can start at one infinitesimal scale while its traced deficit starts at the square of that scale. It is not enough to read the scale of a non-Abelian rotation from a first-order expansion of its trace. Traced loops are classical gauge observables \cite{Wilson}; the theorem gives their exact scale behaviour in this finite $\SU(2)$ setting.

\subsection{Gauge invariance and the quadratic loss for regular class functions}
\TC\ A unit quaternion has limited coordinates, so standard part maps its finite products to ordinary unit-quaternion products. If $q_1$ and $q_2$ are holonomies of two loops with a common base point, a gauge change at that point conjugates both by the same unit quaternion; their commutator is consequently conjugated, $u_1\times u_2$ is rotated, and $w(c)$ and the valuations in \eqref{qgs:hol:eq:comm-val} are independent of that change of gauge frame. \TA\ For unrelated open links, this simultaneous-conjugation hypothesis is not automatic. We always mean based loops or calibrated spin rotations.

\begin{lemma}[\TC\ Conjugacy of imaginary directions]\label{qgs:hol:lem:conjugacy}
Two imaginary quaternions of the same norm are conjugate by a unit quaternion. Consequently two unit quaternions with the same real part are conjugate.
\end{lemma}
\begin{proof}
For nonzero norm, divide by that norm. Let $a,b$ be unit imaginary quaternions. If $b\neq-a$, the nonzero quaternion $g=b+a$ satisfies $ga=bg$, so its unit normalization conjugates $a$ to $b$. If $b=-a$, choose a unit imaginary vector perpendicular to $a$; it anticommutes with $a$ and gives the required conjugation. Zero imaginary parts require no construction. A unit quaternion's imaginary norm is determined by its real part.
\end{proof}

\begin{theorem}[\TC\ Regular Wilson observables have quadratic valuation suppression]\label{qgs:hol:thm:classfunctions}
Let $f$ be the restriction to unit quaternions of a polynomial in their four real coordinates with coefficients in $\OO$. Suppose $f$ is invariant under unit-quaternion conjugation. Then there is $P\in\OO[X]$ such that
\[
 f(q)=P(\Rea q)\qquad(\norm q=1).
\]
If $q\sim1$, $q\neq1$, and $w=1-\Rea q$, then
\begin{equation}\label{qgs:hol:eq:class-bound}
 v\bigl(f(q)-f(1)\bigr)\geq v(w).
\end{equation}
The bound is attained by $f(q)=\Rea q$.
\end{theorem}
\begin{proof}
Write $q=s+u$ and $r=\norm u$. By \cref{qgs:hol:lem:conjugacy}, $f(q)$ equals both $f(s+r\qi)$ and $f(s-r\qi)$. Let $G(s,r)$ be the polynomial obtained by restricting the given coordinate polynomial to $s+r\qi$. Its even part $\tfrac12\bigl(G(s,r)+G(s,-r)\bigr)=H(s,r^2)$ has coefficients in $\OO$. On the unit group, $r^2=1-s^2$. Therefore $P(s)=H(s,1-s^2)$ has the asserted property. The polynomial $P(s)-P(1)$ is divisible by $s-1$ over $\OO$. Since $s$ is finite, evaluating the remaining factor cannot give negative valuation. This proves \eqref{qgs:hol:eq:class-bound}.
\end{proof}

\TA\ For a commutator with scale $\kappa$, such class observables change only at valuation at least $2\kappa$, while a coherent probe achieves trace distance at valuation $\kappa$ (\cref{qgs:hol:prop:queryscale}). There is no conflict with gauge invariance: the latter experiment includes a probe state and a measurement reference which transform along with the holonomy; it is not a scalar function of the holonomy alone. The regularity restriction matters. A function such as $\sqrt{1-\Rea q}$ is not a polynomial regular at the identity and has valuation $\kappa$; dividing a Wilson defect by an infinitesimal coefficient similarly leaves the bounded-coefficient class. The theorem completely treats polynomial functions of one $\SU(2)$ element; it does not classify observables of a network with several loops and reference systems.

Source 08 proves a complementary statement for analytic germs with ordinary coefficients, using Hahn evaluation.

\begin{theorem}[\TC\ Analytic single-loop first-order obstruction]\label{qgs:ops:thm:invariant}
Let $f$ be an ordinary real-coefficient formal or real-analytic scalar germ on $\Sp(1,\R)$ at $1$, invariant under conjugation by every ordinary unit quaternion. Evaluate its Taylor germ at a unit surquaternion $c\ne1$ over a Hahn field with $c-1$ infinitesimal. Then the Hahn evaluation is defined and
\begin{equation}\label{qgs:ops:eq:invariant-order}
 v(f(c)-f(1))\ge2v(c-1).
\end{equation}
The bound is attained by $f(c)=1-\Rea c$. It does not hold for all nonsmooth invariant observables.
\end{theorem}
\begin{proof}
Use imaginary coordinates $x=(x_1,x_2,x_3)$ near the identity:
\[
 c=\sqrt{1-\abs x^2}+x,
 \qquad h(x)=f(\sqrt{1-\abs x^2}+x).
\]
Conjugation acts on $x$ by ordinary rotations. Invariance under the half-turns about the coordinate axes forces every coefficient of the linear part of $h$ to vanish: each coordinate is negated by at least one of these rotations. Thus $h(x)-h(0)$ has only terms of total degree at least two.

If $\gamma=\min_jv(x_j)>0$, the real-coefficient formal series evaluates as a strong Hahn sum by \cref{qgs:obs:lem:support}. Every term of degree at least two has valuation at least $2\gamma$. Moreover
\[
 \sqrt{1-\abs x^2}-1=-\frac{\abs x^2}{1+\sqrt{1-\abs x^2}}
\]
has valuation $2\gamma$, so $v(c-1)=\gamma$. This proves the bound. The displayed identity also proves sharpness for $1-\Rea c$. Finally the invariant function $\abs{c-1}$ has valuation exactly $v(c-1)$; it is not a real-analytic germ at the identity and supplies the promised counterexample without the regularity hypothesis.
\end{proof}

\TA\ The two results are complementary and both are kept: source 04's covers polynomial class functions with finite (possibly non-ordinary) coefficients over any real closed field; source 08's covers analytic germs with ordinary coefficients in the Hahn setting. Source 08's theorem concerns one loop and ordinary analytic coefficients. It does not exclude a singularly normalized observable, a nonsmooth norm, an external reference axis, or an observable of several loops with additional data. For example a calibrated spin direction can couple linearly to the imaginary part of a small loop, but that probe contains a reference transforming with the gauge choice. The theorem should not be read as an absolute ban on detecting a small rotation.

\subsection{Loops on a finite gauge lattice}
Let a finite oriented graph, with a finite specified family of closed faces, carry edge variables $U_e\in\Sp(1,F)$, with $U_{e^{-1}}=U_e^{-1}$. A gauge transformation assigns $g_x\in\Sp(1,F)$ to each vertex $x$ and acts by
\[
 U_e\longmapsto g_{s(e)}U_e g_{t(e)}^{-1}.
\]
For a face $f$, let $q_f$ be the ordered product around its boundary. It transforms by conjugation at its chosen base vertex. The normalized Wilson weight is
\begin{equation}\label{qgs:eq:WilsonWeight}
 w_f=1-\Rea q_f
 =1-\frac12\tr\chi(q_f)=\tfrac12\norm{q_f-1}^2\ge0.
\end{equation}
The inequality follows from $N(q_f)=1$, which implies $-1\le\Rea q_f\le1$. These are the usual finite $\SU(2)$ loop variables, written quaternionically \cite{Wilson}; no integration over a gauge-field space is required for the statements below. \TA\ This is a specified finite function, not a path integral or a continuum quantum field theory.

\begin{lemma}[\TC\ Small unit quaternions have quadratic Wilson weight]\label{qgs:lem:smallq}
Let $q=s+u\in\Sp(1,F)$, where $s\in F$ is real and $u$ is purely imaginary, over a Hahn field. Assume $\st q=1$ and $q\ne1$. Then $u\ne0$. If $r=v(u)$ and $u_0=\st(t^{-r}u)$, then $r>0$ and
\begin{equation}\label{qgs:eq:smallq}
 v(1-s)=2r,\qquad
 \st\bigl(t^{-2r}(1-s)\bigr)=\frac12 N(u_0).
\end{equation}
\end{lemma}
\begin{proof}
The norm condition is $s^2+N(u)=1$. Since $\st s=1$, the factor $1+s$ is appreciable with residue two, and
\begin{equation}\label{qgs:eq:qnormid}
 1-s=\frac{N(u)}{1+s}.
\end{equation}
If $u=0$, then $s=\pm1$, and $\st s=1$ forces $s=1$, contrary to $q\ne1$. Therefore $r$ is defined and positive. \Cref{qgs:lem:squares} gives $v(N(u))=2r$ and leading coefficient $N(u_0)>0$. Equation \eqref{qgs:eq:qnormid} proves the result.
\end{proof}

\begin{theorem}[\TC\ Gauge-invariant leading curvature and positive action]\label{qgs:thm:wilson}
Over a Hahn field, assume every face holonomy satisfies $\st q_f=1$. For each nonflat face $q_f\ne1$, put
\[
 r_f=v(\Ima q_f)=v(q_f-1),\qquad
 u_f=\st\bigl(t^{-r_f}\Ima q_f\bigr).
\]
Let $\lambda_f>0$ be arbitrary Hahn weights, not necessarily finite, and define the finite-face action
\[
 \mathcal S=\sum_f\lambda_f(1-\Rea q_f).
\]
If at least one face is nonflat, then
\begin{equation}\label{qgs:eq:WilsonVal}
 v(\mathcal S)=\min_{f:q_f\ne1}\bigl(v(\lambda_f)+2r_f\bigr)=:\gamma,
\end{equation}
and
\begin{equation}\label{qgs:eq:WilsonLC}
 \lc(\mathcal S)=
 \frac12\sum_{\substack{f:q_f\ne1\\v(\lambda_f)+2r_f=\gamma}}
 \lc(\lambda_f)N(u_f)>0.
\end{equation}
Each $r_f$ is gauge invariant. The initial vector $u_f$ transforms by the ordinary rotation induced by the residue of the gauge transformation at the base vertex; its squared norm is therefore invariant. In particular $\mathcal S=0$ if and only if every face in the sum is flat.
\end{theorem}
\begin{proof}
Every unit quaternion has finite coordinates, so a gauge transformation $g$ has an ordinary unit-quaternion residue $g_0$. Conjugation preserves real part and maps the imaginary part to $g(\Ima q_f)g^{-1}$, preserving its quaternion norm. \Cref{qgs:lem:squares} therefore implies invariance of $r_f$, and after division by $t^{r_f}$ its residue is $g_0u_fg_0^{-1}$. This is the ordinary orthogonal action on imaginary quaternions. Since $v(1-s)=2r_f>r_f$ by \cref{qgs:lem:smallq}, $v(q_f-1)=v(\Ima q_f)$.

\Cref{qgs:lem:smallq} gives the valuation and coefficient of each nonzero Wilson weight. Multiplication by $\lambda_f$ adds valuations and multiplies leading coefficients. All those coefficients are positive. Since there are finitely many faces, the coefficient at the least resulting exponent is exactly the positive sum in \eqref{qgs:eq:WilsonLC}; it cannot cancel. The flat case is immediate.
\end{proof}

\status{\TC\ Source 07 (Theorem 10.2 there). Source 04 proves the valuation formula \eqref{qgs:eq:WilsonVal} and the zero criterion over any real closed field, with $\kappa_P=v(q_P-1)$ and $w_P=\norm{q_P-1}^2/2$ (its Proposition 9.1); source 07 adds the leading coefficient and gauge covariance of the initial curvature. \TA\ The hypothesis of positive weights is substantive: with sign-changing or complex weights, leading contributions can cancel. This is a theorem about a finite lattice action, not about a quantum path integral. It contains no claim of a Yang--Mills mass gap, confinement, existence of a continuum measure, or a thermodynamic limit.}

\TA\ An interesting regime is $v(\lambda_f)=-2r_f$. Then a loop that is exactly flat in the ordinary shadow contributes an appreciable action after coupling rescaling. That is a finite algebraic analogue of retaining a leading curvature term while lattice holonomies approach the identity. It is a candidate tool for checking scale assignments, not a derivation of a continuum limit. A positive action weight of negative valuation may make an infinitesimal defect appreciable, but that algebraic weighting is not automatically a bounded measurement effect or an experimentally accessible coupling. The usefulness of the ordered-group formulation is that different faces or couplings may live at different ranks: it gives an exact leading-action certificate and a test for whether an alleged truncation has deleted the first nonzero curvature.

\TC\ A homogeneous two-direction face with links $Q(u_1)$ and $Q(u_2)$ has exactly the commutator of \cref{qgs:hol:cor:comm}, so \eqref{qgs:hol:eq:comm-w} is a closed expression for the noncommuting part of its curvature. If $u_1=aA$ and $u_2=bB$ for small dimensionless step sizes and fixed imaginary directions $A,B$, its leading defect is $2a^2b^2\norm{A\times B}^2$. \TA\ This isolates the constant-connection commutator contribution; it does \emph{not} include the derivative terms of a general inhomogeneous continuum field strength.

\subsection{A two-scale commutator plaquette}
For a purely imaginary ordinary quaternion $a$ and an infinitesimal $\epsilon$, the formal exponential $e^{\epsilon a}=\sum_{n\ge0}\epsilon^n a^n/n!$ is strongly summable and has norm one. If $a\ne0$, it equals $\cos(\epsilon\abs a)+(a/\abs a)\sin(\epsilon\abs a)$, where sine and cosine are their positive-order formal series.

\begin{proposition}[\TC\ Non-Abelian curvature at a mixed scale]\label{qgs:prop:plaquette}
Let $a,b$ be purely imaginary ordinary quaternions with $[a,b]\ne0$, and let $\epsilon=t^\alpha$, $\zeta=t^\beta$ with $\alpha,\beta>0$. Define
\[
 Q=e^{\epsilon a}e^{\zeta b}e^{-\epsilon a}e^{-\zeta b}.
\]
Then
\begin{align}
 Q&=1+\epsilon\zeta[a,b]+R,
 &v(R)&>\alpha+\beta,\label{qgs:eq:commloop}\\
 v(1-\Rea Q)&=2(\alpha+\beta),
 &\lc(1-\Rea Q)&=\tfrac12 N([a,b]).\label{qgs:eq:commenergy}
\end{align}
Identifying imaginary quaternions with ordinary three-vectors, the leading Wilson weight is $2\epsilon^2\zeta^2\abs{a\times b}^{2}$. If $a$ and $b$ commute, $Q=1$ exactly.
\end{proposition}
\begin{proof}
Multiply the four formal series. Terms depending on only one of the two parameters cancel because the expression is one when either parameter is zero. The coefficient of $\epsilon\zeta$ is $ab-ba$. Every remaining nonconstant monomial has positive powers of both parameters and total degree at least three, so its valuation exceeds $\alpha+\beta$. Strong evaluation justifies the multiplication at arbitrary rank. The bracket of imaginary quaternions is imaginary, so the initial imaginary part is $[a,b]$ at order $\alpha+\beta$. \Cref{qgs:lem:smallq} yields \eqref{qgs:eq:commenergy}. Finally, $[a,b]=2a\times b$ in vector notation. If the two generators commute, all exponentials commute and the product cancels exactly.
\end{proof}

\status{\TC\ Source 07. \TI\ The expansion \eqref{qgs:eq:commloop} is the group commutator of \texttt{squat:prop:bch} in the surquaternion report \cite{repoQuaternion}, whose \texttt{squat:eq:groupcomm} gives the stronger remainder bound $v(R)\ge\alpha+\beta+\min\{\alpha,\beta\}$; source 07 did not cite it. The exact identities of \cref{qgs:ops:thm:commutator,qgs:hol:cor:comm} (sources 04 and 08) sharpen the leading-term statement to an exact norm identity.}

\TE\ For $a=\qi$ and $b=\qj$ there is the exact identity
\begin{equation}\label{qgs:eq:exactplaquette}
 1-\Rea\left(e^{\epsilon\qi}e^{\zeta\qj}
 e^{-\epsilon\qi}e^{-\zeta\qj}\right)
 =2\sin^2\epsilon\,\sin^2\zeta,
\end{equation}
which source 07's verification checks as a polynomial identity in sine and cosine coordinates subject to the two unit-circle relations. \merge\ More generally, since $\Ima e^{\epsilon a}=(a/\abs a)\sin(\epsilon\abs a)$, \cref{qgs:ops:thm:commutator} gives, for all nonzero ordinary imaginary $a,b$,
\[
 1-\Rea Q=\frac{2\sin^2(\epsilon\abs a)\sin^2(\zeta\abs b)}{\abs a^2\abs b^2}\,\abs{a\times b}^2,
\]
of which \eqref{qgs:eq:exactplaquette} and \eqref{qgs:eq:commenergy} are the special case and the leading term.

\TA\ In the rank-two hierarchy $\zeta\ll\epsilon^N$, this mixed curvature lies below every pure algebraic order in $\epsilon$. Yet its initial bracket and its action coefficient remain exact. This is the non-Abelian counterpart of the information retained in \cref{qgs:thm:alljets}: a chosen one-scale truncation can be perfectly accurate at every displayed algebraic order and still omit a mathematically meaningful sector.

\section{Infinitesimal holonomy as a quantum resource}\label{qgs:sec:discrimination}
\status{\TC\ Source 04, over an arbitrary real closed $F\supseteq\R$ with $K=F[i]$ and its natural valuation; the value group need not embed in $\R$. All dimensions, numbers of Kraus operators, circuit depths and query counts are ordinary finite integers. \TI\ State separation is due to Chefles and Barnett \cite{CheflesBarnett}, with later work on general priors \cite{Bagan}; the unitary-discrimination geometry is Ac\'in's \cite{Acin}. These classical formulas are credited, not claimed as new; source 04 gives algebraic proofs over $F$, an adaptive finite-query bound in that setting, and the combined arbitrary-rank interpretation for gauge curvature.}

\TA\ Source 04 asks exactly what a hierarchy of weak couplings, near alignment and measurement efficiencies permits a finite protocol to distinguish. The holonomy may represent a plaquette in an algebraic lattice-gauge model or the net unitary of a sequence of noncommuting controls. The proposal is to put the \emph{holonomy, the success probability, and the calibration error in the same ordered scale group}, rather than recording them in unrelated asymptotic estimates; the resulting resource inequalities retain their meaning when one scale is smaller than every power of another. Non-Abelian holonomy in quantum control is classical \cite{ZanardiRasetti}. The physical uses proposed below are a scale calculus for gauge probes, noncommuting control and symbolic certification; no actual infinitesimal coupling or heralding device is established, and candidate originality means originality of the synthesis to source 04's investigation, not certified historical priority.

\subsection{The main result in operational form}\label{qgs:hol:sec:main}
Fix nonzero infinitesimal imaginary quaternions $u_1,u_2\in\HF$ with $u_1\times u_2\neq0$. The two hypotheses supplied to the experimenter are the known gates
\[
 I_2\qquad\text{and}\qquad U=\chi(c),\qquad c=Q(u_1)Q(u_2)Q(u_1)^{-1}Q(u_2)^{-1}.
\]
The experimenter knows the two possibilities, not which one was supplied. Arbitrary finite-dimensional ancillas and hypothesis-independent quantum controls over $K$ are allowed. Write $\kappa=v(u_1\times u_2)>0$ and let $T_m$ be the Chebyshev polynomial. For every fixed $m\ge1$ we prove in \cref{qgs:hol:thm:query,qgs:hol:thm:frontier,qgs:hol:thm:valuation}:
\begin{enumerate}
\item The optimal error-free conclusive probability with equal priors, over protocols with at most $m$ queries, finite ancillas and adaptive control, is
\[
 p_{\mathrm{ua}}^{(m)}=w_m=1-T_m(1-w),\qquad
 w=\frac{2\norm{u_1\times u_2}^2}{(1+\norm{u_1}^2)(1+\norm{u_2}^2)}.
\]
It has valuation $2\kappa$ and leading term $m^2w$.
\item For a common heralding probability $p\in F$, $0<p\leq1$, the maximum distinguishability of the conditional outputs is
\begin{equation}\label{qgs:hol:eq:intro-frontier}
 \TD_{m,\max}(p)=
 \begin{cases}
 1,&p\leq w_m,\\
 \sqrt{2w_m/p-(w_m/p)^2},&p>w_m,
 \end{cases}
 \qquad
 v(\TD_{m,\max}(p))=\max\{0,\kappa-v(p)/2\}.
\end{equation}
\item No ordinary finite number of queries permits deterministic perfect discrimination. An appreciable heralding probability cannot produce appreciably different conditional outputs.
\end{enumerate}
Error-free conclusive identification is \emph{unambiguous discrimination}: the protocol may announce ``inconclusive'', but may not announce the wrong gate.

\subsection{Chebyshev algebra and the optimal query overlap}
We first treat an arbitrary unit quaternion $q$ with $q\sim1$, $q\neq1$, not necessarily a commutator. Put
\[
 U=\chi(q),\quad s=\Rea q=1-w,\quad
 r=\norm{\Ima q}=\sqrt{1-s^2}.
\]
Then $0<w\in\mm$, $s>0$, and the eigenvalues of $U$ are $s\pm ir$. Explicitly, $A=(U-sI)/(ir)$ is a Hermitian involution of trace zero, because $(U-sI)^2=-r^2I$ and $(U-sI)^*=-(U-sI)$. Applying \cref{qgs:prop:spectral} to $A$ supplies an orthonormal eigenbasis with eigenvalues $+1,-1$, and hence the asserted eigenbasis of $U$.

The Chebyshev polynomials $T_m$ and $\mathsf U_m$ are defined by
\[
 T_0=1,\ T_1=X,\ T_{m+1}=2XT_m-T_{m-1},\qquad
 \mathsf U_0=1,\ \mathsf U_1=2X,\
 \mathsf U_{m+1}=2X\mathsf U_m-\mathsf U_{m-1}.
\]
\TE\ Finite multiplication of $(s+ir)^m$ gives
\begin{equation}\label{qgs:hol:eq:cheby}
 (s+ir)^m=T_m(s)+ir\,\mathsf U_{m-1}(s),\qquad
 T_m(s)^2+(1-s^2)\mathsf U_{m-1}(s)^2=1.
\end{equation}
Because $s\sim1$, for every ordinary $m\geq1$ we have $T_m(s)\sim1$ and $\mathsf U_{m-1}(s)\sim m>0$. Thus their signs are known without trigonometry. In particular, $\sqrt{1-T_m(s)^2}=r\,\mathsf U_{m-1}(s)$.

\begin{theorem}[\TC\ Optimal finite-query overlap]\label{qgs:hol:thm:query}
For the two known gates $I_2$ and $U$, any protocol using at most $m\geq1$ queries, finite ancillas, and hypothesis-independent adaptive quantum control admits common-circuit purifications of its two final states whose overlap is at least
\[
 c_m=T_m(s)>0.
\]
The bound is attained by a pure qubit probe, $m$ consecutive applications of the unknown gate, and no ancilla. Consequently the maximum output trace distance is
\begin{equation}\label{qgs:hol:eq:Dm}
 \TD_m=\sqrt{1-T_m(s)^2},
\end{equation}
and the optimal minimum-error success probability is $(1+\TD_m)/2$.
\end{theorem}
\begin{proof}
Purify the initial state and dilate all controls, deferring their finite classical records as in \cref{qgs:sec:instruments}. Consider the usual sequence of $m+1$ final hybrid pure states $\Psi_0,\ldots,\Psi_m$: the first $j$ query slots use $U$, the remaining slots use $I$. Adjacent hybrids differ in one slot. Everything before that slot is the same in the two hybrids, and everything after it is the same isometry. Their inner product is therefore an expectation of $U\otimes I$ in a unit vector.

Since $(U+U^*)/2=sI$, its real part is $s$, and the modulus of this expectation is at least $s$. If a branch uses fewer queries, a coherent controlled skip has Hermitian part at least $sI$, giving the same lower bound. Thus $\abs{\Psi_{j-1}^*\Psi_j}\geq s$. Inductively apply \cref{qgs:hol:lem:triangle}. If $\abs{\Psi_0^*\Psi_j}\geq T_j(s)$, then
\begin{align*}
 \abs{\Psi_0^*\Psi_{j+1}}
 &\geq T_j(s)s-\sqrt{1-T_j(s)^2}\,r
 =T_j(s)s-(1-s^2)\mathsf U_{j-1}(s)
 =T_{j+1}(s).
\end{align*}
The final equality is the real-part multiplication identity in \eqref{qgs:hol:eq:cheby}. Its right side is positive in the present infinitesimal regime. This establishes the upper bound on distinguishability for all such protocols. Discarding the purifying environment cannot improve trace distance by \cref{qgs:lem:tracenorm}.

For attainability, choose orthonormal eigenvectors $e_+,e_-$ of $U$ and prepare $\psi=(e_++e_-)/\sqrt2$. The final two possible states are $\psi$ and $U^m\psi$, with overlap
\[
 \psi^*U^m\psi=\frac{(s+ir)^m+(s-ir)^m}{2}=T_m(s).
\]
Now apply \cref{qgs:hol:lem:puredistance} and \eqref{qgs:hol:eq:helstrom}.
\end{proof}

\TA\ The algebraic proof is related to the ordinary spectral-arc geometry of gate discrimination \cite{Acin}. Its useful feature here is uniformity over real closed fields and arbitrary infinitesimal scales; no angle or Archimedean threshold is concealed in the argument. Multi-query and sequential discrimination of unitaries has further antecedents beyond Ac\'in, for example Duan, Feng and Ying's perfect discrimination by sequential schemes without entanglement \cite{DuanFengYing}; source 04's comparison was targeted and did not survey them.

\subsection{Failure of eventual finite perfect discrimination}
\begin{corollary}[\TC\ No finite perfect gate test]\label{qgs:hol:cor:nofinite}
If $q\sim1$ but $q\neq1$, then $I_2$ and $\chi(q)$ are distinct quantum channels, yet no ordinary finite-query protocol distinguishes them perfectly with certainty. In fact every such protocol has infinitesimal trace-distance advantage over random guessing.
\end{corollary}
\begin{proof}
The only scalar matrices in $\SU(2,K)$ are $\pm I$. Near $I$, the condition $q\neq1$ excludes both scalar-channel coincidences. For each fixed $m$, \cref{qgs:hol:thm:query} gives strictly positive overlap $T_m(s)\sim1$ and infinitesimal $\TD_m$. Perfect certain discrimination would require orthogonal outputs.
\end{proof}

\TC\ For ordinary complex unitary channels that differ beyond global phase, sufficiently many uses allow perfect discrimination \cite{Acin}. The existential statement ``there is a finite number of uses'' has changed. There is no contradiction with real-closed-field algebra: quantification over all ordinary natural query counts is not a single fixed-size polynomial statement. Nor is it legitimate to insert an infinite omnific integer for $m$ in a finite circuit or in a Chebyshev polynomial indexed by an ordinary integer; that would require an additional model of infinite computation and measurement.

\TA\ Source 04 adds: ``The absence of an Archimedean query threshold is a genuine structural difference from ordinary complex quantum theory.'' This is true of the model \emph{with literal infinitesimal gate differences} only, and must not be read as a physical difference. Literal infinitesimals are not values of real parameters (\cref{qgs:hol:prop:noeval}); every ordinary specialization of the same finite formulas is an ordinary pair of gates, with Ac\'in's finite threshold, which merely grows without bound along a real family whose gate difference tends to zero; and the optimal Chebyshev branch used here is only the small-angle branch of that real family (\cref{qgs:hol:sec:real}).

\TC\ The obstruction is broader than the qubit example. If two finite surcomplex channels have identical reduced Kraus-instrument behaviour, every finite circuit built from them reduces to the same ordinary output distribution (\cref{qgs:prop:shadow}). A deterministic zero-versus-one event is then impossible. The qubit theorem supplies the exact quantitative refinement.

\begin{proposition}[\TC\ Finite queries improve coefficients, not valuations]\label{qgs:hol:prop:queryscale}
Set $w_m=1-T_m(1-w)$. For fixed $m\geq1$,
\begin{equation}\label{qgs:hol:eq:wmexp}
 w_m=m^2w-\frac{m^2(m^2-1)}6w^2+w^3R_m(w),\qquad R_m\in\Q[w],
\end{equation}
where zero coefficients are interpreted in the obvious way for small $m$. Hence
\[
 w_m\sim m^2w,\qquad
 \TD_m\sim m\sqrt{2w},\qquad
 v(w_m)=v(w),\quad v(\TD_m)=\tfrac12v(w).
\]
\end{proposition}
\begin{proof}
Differentiate the Chebyshev recurrence at $1$ to obtain $T_m'(1)=m^2$ and $T_m''(1)=m^2(m^2-1)/3$ by induction. The finite Taylor identity for the polynomial gives \eqref{qgs:hol:eq:wmexp}. Since $w$ is infinitesimal, the leading coefficient is a positive real number. Finally $\TD_m^2=2w_m-w_m^2$.
\end{proof}

\TA\ The term ``finite'' here does not mean a conveniently small numerical value of $m$. It means an ordinary integer fixed independently of a literal infinitesimal. Any such integer has valuation zero.

\subsection{Exact postselection and unambiguous identification}
A success map $\Phi$ has common success probability $p$ for a pair of states when $\tr\Phi(\rho)=\tr\Phi(\sigma)=p$. Keeping these probabilities equal prevents an undeclared change of prior odds when success is announced. The separate task of optimizing unambiguous identification with equal prior probabilities will not require this equality assumption for the upper bound.

\begin{theorem}[\TC\ Sharp equal-success separation; classical formula over $F$]\label{qgs:hol:thm:separation}
Let $x,y$ be unit vectors with $0\leq c=\abs{x^*y}<1$. Among all finite success instruments of common probability $0<p\leq1$, the largest conditional trace distance is
\begin{equation}\label{qgs:hol:eq:sep}
 \TD_{\max}(c,p)=\sqrt{1-\left(\max\{0,(c+p-1)/p\}\right)^2}.
\end{equation}
This bound allows mixed successful outputs and is attained by a single diagonal filter on the span of $x,y$.
\end{theorem}
\begin{proof}
Dilate the success instrument, retaining all its environment registers. Its successful vectors $x_s,y_s$ have norm $\sqrt p$; its failure vectors have norm $\sqrt{1-p}$. Preservation of inner products gives $x^*y=x_s^*y_s+x_f^*y_f$. Therefore $\abs{x_s^*y_s}\geq c-\abs{x_f^*y_f}\geq c-(1-p)$. The normalized successful pure states have overlap at least $d=\max\{0,(c+p-1)/p\}$. Their trace distance is at most $\sqrt{1-d^2}$ by \cref{qgs:hol:lem:puredistance}. Discarding an environment cannot improve it.

For equality, adjust the initial phases and basis so the two states are
\[
 \psi_\pm=\sqrt{\frac{1+c}{2}}\ket0
 \ \pm\sqrt{\frac{1-c}{2}}\ket1,
\]
define desired outputs $\phi_\pm=\sqrt{(1+d)/2}\ket0\pm\sqrt{(1-d)/2}\ket1$, and the filter
\begin{equation}\label{qgs:hol:eq:filterK}
 M_{\mathrm{s}}=\diag\left(
 \sqrt{\frac{p(1+d)}{1+c}},\,
 \sqrt{\frac{p(1-d)}{1-c}}\right).
\end{equation}
It sends $\psi_\pm$ to $\sqrt p\,\phi_\pm$. If $p\leq1-c$, then $d=0$ and both diagonal squares are at most one. If $p>1-c$, the second diagonal square equals one, while the first equals $(2p-1+c)/(1+c)\leq1$. Thus $M_{\mathrm{s}}^*M_{\mathrm{s}}\leq I$. The complementary filter $(I-M_{\mathrm{s}}^*M_{\mathrm{s}})^{1/2}$ completes a legitimate instrument. The outputs have the required distance, proving attainability.
\end{proof}

\TA\ The real-complex ancestor of this result is precisely the state-separation bound of Chefles and Barnett \cite{CheflesBarnett}. The construction is included because it establishes attainability over $F$, including when $p$ is infinitesimal, and because it prevents the operational interpretation from relying on a formal division without a valid physical map.

\begin{lemma}[\TC\ Equal-prior unambiguous bound]\label{qgs:hol:lem:ua}
For two pure states of overlap $c$, the average success probability of any error-free conclusive discrimination instrument with equal prior probabilities is at most $1-c$. Equality is attained with equal success probabilities.
\end{lemma}
\begin{proof}
Purify the full instrument and include the outcome label. Error-free conclusive branches have orthogonal labels under the two hypotheses, so their contribution to the inner product vanishes. If failure probabilities are $q_0,q_1$, then $c\leq\sqrt{q_0q_1}\leq(q_0+q_1)/2$. The mean success probability is at most $1-c$. When $c=1$, always declaring failure attains this zero bound. For $c<1$, equality follows from \cref{qgs:hol:thm:separation} with $p=1-c$, whose successful outputs are orthogonal, followed by their exact projective identification.
\end{proof}

\TI\ \merge\ Over $\C$ this is the Ivanovic--Dieks--Peres bound \cite{Ivanovic,Dieks,Peres}, which source 04 used without naming.

\begin{theorem}[\TC\ Query and postselection frontier]\label{qgs:hol:thm:frontier}
For $I_2$ versus $U=\chi(q)$ with $q\sim1$, $q\neq1$, and at most $m$ queries, let $w=1-\Rea q$ and $w_m=1-T_m(1-w)$. Then
\begin{align}
 p_{\mathrm{ua}}^{(m)}&=w_m,\label{qgs:hol:eq:ua-m}\\
 \TD_{m,\max}(p)&=
 \begin{cases}
 1,&0<p\leq w_m,\\
 \sqrt{2w_m/p-(w_m/p)^2},&w_m<p\leq1.
 \end{cases}\label{qgs:hol:eq:frontier}
\end{align}
The first optimum is over all error-free conclusive instruments with equal prior probabilities, whether or not their two success probabilities coincide. The second is over common-success-$p$ instruments. Both bounds allow finite ancillas and adaptive control and are attained by the repeated-query probe of \cref{qgs:hol:thm:query} followed by \eqref{qgs:hol:eq:filterK}.
\end{theorem}
\begin{proof}
Purified outputs before the final instrument have overlap at least $c_m=T_m(1-w)$. Apply \cref{qgs:hol:lem:ua}. For the conditional distance, \eqref{qgs:hol:eq:sep} is nonincreasing in $c$, so $c=c_m$ gives a universal upper bound. Substitute $c_m=1-w_m$. The repeated-query probe attains that overlap, and the explicit filter attains the state-separation bound.
\end{proof}

\TA\ It is important that discarded histories remain accounted for. An adaptive scheme that first postselects an unlikely branch is still a finite instrument when its failure histories are retained. Its total successful probability is included in $p$ or $p_{\mathrm{ua}}^{(m)}$. It cannot evade \cref{qgs:hol:thm:frontier} by declaring intermediate failures to be outside the experiment. The frontier is attained, so it is the exact price that \cref{qgs:thm:cost} guarantees must be paid for this task.

\subsection{The arbitrary-rank gauge-resource law}
\begin{theorem}[\TC\ Valuative curvature--success tradeoff]\label{qgs:hol:thm:valuation}
For the commutator of \cref{qgs:hol:cor:comm}, suppose $u_1,u_2$ are infinitesimal and $u_1\times u_2\neq0$. Let
\[
 \kappa=v(u_1\times u_2),\qquad \beta=v(p),\qquad 0<p\leq1.
\]
For every ordinary finite $m\geq1$,
\begin{align}
 v\bigl(p_{\mathrm{ua}}^{(m)}\bigr)&=2\kappa,\label{qgs:hol:eq:valua}\\
 v\bigl(\TD_{m,\max}(p)\bigr)&=
 \max\{0,\kappa-\beta/2\}.\label{qgs:hol:eq:valD}
\end{align}
These equalities hold in an arbitrary divisible ordered value group. If $\beta<2\kappa$, the stronger leading-term statement is
\begin{equation}\label{qgs:hol:eq:leadingD}
 \TD_{m,\max}(p)\sim \frac{2m\norm{u_1\times u_2}}{\sqrt p}.
\end{equation}
\end{theorem}
\begin{proof}
By \cref{qgs:hol:cor:comm,qgs:hol:prop:queryscale}, $w_m\sim2m^2\norm{u_1\times u_2}^2$, with valuation $2\kappa$. This proves \eqref{qgs:hol:eq:valua}. If $\beta<2\kappa$, then $p>w_m$ and $r_0=w_m/p$ is a positive infinitesimal. The exact formula gives $\TD^2=2r_0-r_0^2\sim2r_0$, proving \eqref{qgs:hol:eq:valD} and \eqref{qgs:hol:eq:leadingD} in this case.

If $\beta>2\kappa$, then $p<w_m$ and $\TD=1$. If $\beta=2\kappa$, either $p\leq w_m$ and again $\TD=1$, or $0<w_m/p<1$ is an appreciable unit. In the latter case $2w_m/p-(w_m/p)^2$ has positive real standard part. Its square root therefore has valuation zero. These are all cases in a totally ordered group.
\end{proof}

\begin{corollary}[\TC\ Appreciable conditional signal requires a rare event]\label{qgs:hol:cor:rare}
For this known-pair discrimination task, an appreciable conditional trace distance is attainable at probability $p$ if and only if $v(p)\geq2\kappa$. In particular, $p$ must be infinitesimal. With appreciable $p$, every finite-query conditional signal is infinitesimal of valuation at least $\kappa$.
\end{corollary}
\begin{proof}
An attained maximum is appreciable precisely when its valuation is zero. Apply \eqref{qgs:hol:eq:valD}. The remaining assertion follows because $v(p)=0$ for appreciable $p$.
\end{proof}

\TC\ For any achieved nonzero conditional distance with valuation $\delta$, the theorem gives the necessary inequality
\begin{equation}\label{qgs:hol:eq:resourceineq}
 \beta+2\delta\geq2\kappa.
\end{equation}
The lower boundary is attained. \TA\ This is not a statement that \emph{every} output distance has that valuation; deliberately poor probes may have much smaller distance or no signal at all.

\begin{figure}[ht]
\centering
\begin{tikzpicture}[x=1.35cm,y=1.25cm]
\draw[->] (0,0)--(4.6,0) node[right] {$\beta/\kappa$};
\draw[->] (0,0)--(0,2.25) node[above] {$\delta/\kappa$};
\draw[thick] (0,1)--(2,0)--(4.15,0);
\draw (0,1)--(-0.06,1);
\node[left] at (0,1) {$1$};\node[below] at (2,0) {$2$};
\node at (2.7,1.55) {allowed by the valuation bound};
\node[rotate=-27] at (0.94,0.8) {attained frontier};
\node at (3.0,-0.45) {rarer success};
\end{tikzpicture}
\caption{A rank-one visualization of \eqref{qgs:hol:eq:resourceineq} (source 04). Smaller $\delta$ means a larger signal; larger $\beta$ means a rarer successful event. The theorem itself is ordered-group valued and does not require numerical division by $\kappa$.}
\label{qgs:hol:fig:frontier}
\end{figure}

\TA\ The three quantities $w(c)$, $\TD_m$ and $p_{\mathrm{ua}}^{(m)}$ are not interchangeable. Their valuations for a commutator are respectively $2\kappa$, $\kappa$, and $2\kappa$. The first is a regular scalar gauge invariant, the second is an optimal deterministic coherent distinguishability, and the third is an error-free successful event rate. The signal can be linearly visible to a coherent probe while quadratically suppressed in a Wilson defect and in conclusive-event probability.

\TC\ At the crossover, write $p=\lambda w_m$ for an appreciable positive real constant $\lambda$, with $p\leq1$. Then $\TD_{m,\max}(p)=1$ for $0<\lambda\le1$ and $\sqrt{2/\lambda-1/\lambda^2}$ for $\lambda>1$. The numerical crossover survives even though the event rate itself is infinitesimal. A standard-part-only description of the \emph{unnormalized} experiment records a zero-probability branch and therefore loses this information. Repeating a fixed successful protocol an ordinary finite number of times does not make it appreciable (\cref{qgs:sec:copies}).

\subsection{Higher-rank examples and precision certificates}
\begin{example}[\TC\ Orthogonal weak controls]\label{qgs:hol:ex:orthogonal}
Take $u_1=\epsilon\qi$, $u_2=\zeta\qj$, with positive infinitesimal $\epsilon,\zeta$. Then
\[
 w=\frac{2\epsilon^2\zeta^2}{(1+\epsilon^2)(1+\zeta^2)},\qquad
 \kappa=v(\epsilon)+v(\zeta).
\]
For a common success probability of valuation $\beta<2\kappa$, $\TD_{m,\max}(p)\sim2m\epsilon\zeta/\sqrt p$.
\end{example}

\begin{example}[\TE\ A suppressed transverse component controls the whole answer]\label{qgs:hol:ex:aligned}
Take $u_1=\epsilon\qi$ and $u_2=\epsilon\qi+\epsilon^3\qj$. Both control norms have valuation $v(\epsilon)$, but
\[
 u_1\times u_2=\epsilon^4\qk,
 \qquad
 w=\frac{2\epsilon^8}{(1+\epsilon^2)(1+\epsilon^2+\epsilon^6)}.
\]
Thus the Wilson defect and unambiguous probability are of valuation $8v(\epsilon)$, not the generic $4v(\epsilon)$ suggested by multiplying the two leading control norms. The optimal deterministic distance is $\sim2m\epsilon^4$. Replacing $u_2$ by its leading term $\epsilon\qi$ changes the exact holonomy to the identity and deletes the signal altogether.
\end{example}

\TA\ This is not merely an additional term in a Baker--Campbell--Hausdorff expansion. It is a change from a nonzero resource to no resource. The exact cross product in \eqref{qgs:hol:eq:comm-w} detects it before any truncation is made.

\begin{example}[\TC\ Lexicographic rank two]\label{qgs:hol:ex:ranktwo}
In the rank-two workspace \eqref{qgs:eq:ranktwo} put $u_1=\epsilon\qi$ and $u_2=\epsilon\qi+\zeta\qj$. Both $\epsilon,\zeta$ are infinitesimal, but $\zeta<\epsilon^n$ for every ordinary positive integer $n$. We have
\[
 u_1\times u_2=\epsilon\zeta\qk,\quad \kappa=(1,1),\quad
 w=\frac{2\epsilon^2\zeta^2}{(1+\epsilon^2)(1+\epsilon^2+\zeta^2)}\sim2\epsilon^2\zeta^2.
\]
For $p=\epsilon^\varrho$ with fixed positive rational $\varrho$,
\[
 \beta=(0,\varrho)<(2,2)=2\kappa,
 \qquad
 v(\TD_{m,\max}(p))=(1,1-\varrho/2)>0.
\]
No such power-law heralding probability makes the signal appreciable, however large the fixed exponent $\varrho$. At the much rarer probability $p=w_m\sim2m^2\epsilon^2\zeta^2$, perfect conditional discrimination is attainable. One concrete surreal realization is $\epsilon=\omega^{-1}$ and $\zeta=\omega^{-\omega}$.
\end{example}

\TA\ This example illustrates why an arbitrary-rank value group is useful: the failure of every finite-order strategy is an order comparison, not a vague assertion that an exponential is ``very small.''

\begin{theorem}[\TC\ Certificate preserving curvature and operational scales]\label{qgs:hol:thm:precision}
Let $u_1,u_2$ be infinitesimal imaginary quaternions with $u_1\times u_2\neq0$, and write
\[
 a=v(u_1),\qquad b=v(u_2),\qquad \kappa=v(u_1\times u_2)\geq a+b.
\]
Suppose approximations $u_1'=u_1+e$, $u_2'=u_2+f$ satisfy
\begin{equation}\label{qgs:hol:eq:precisionhyp}
 v(e)>\kappa-b,\qquad v(f)>\kappa-a.
\end{equation}
Then $v(u_1'\times u_2'-u_1\times u_2)>\kappa$, and the two Wilson defects obey $w'/w\in1+\mm$. Their optimal $m$-query unambiguous probabilities have the same leading term, and their conditional optimal distances at each fixed $p>0$ have the same valuation.
\end{theorem}
\begin{proof}
The cross-product difference is $e\times u_2+u_1\times f+e\times f$. The first two terms have valuation greater than $\kappa$. The last has valuation greater than $2\kappa-a-b\geq\kappa$. Thus the cross products have the same leading vector. By positivity and \eqref{qgs:hol:eq:normval}, their squared norms have quotient in $1+\mm$. The hypotheses also keep $u_1',u_2'$ infinitesimal, so both denominators in \eqref{qgs:hol:eq:comm-w} lie in $1+\mm$. This proves $w'/w\in1+\mm$. The leading query result follows from $w_m\sim m^2w$. Finally \cref{qgs:hol:thm:valuation} depends only on $\kappa$ and the fixed valuation of $p$.
\end{proof}

\TA\ The excess $\kappa-a-b$ is an alignment-cancellation depth. Merely knowing that $e$ is small relative to $u_1$ and $f$ small relative to $u_2$ is not enough when that depth is positive. The strict inequality in \eqref{qgs:hol:eq:precisionhyp} is a real boundary: in \cref{qgs:hol:ex:aligned}, taking $e=0$, $f=-\epsilon^3\qj$ gives $v(f)=\kappa-a$ and annihilates the curvature. Thus a uniform theorem cannot replace the strict inequality by a weak one without extra information. The preceding theorem concerns errors in holonomy coefficients; errors that enter when a rare measurement event is normalized are priced by \cref{qgs:thm:cost,qgs:hol:cor:calibration}.

\subsection{Applications to gauge probes and quantum control}\label{qgs:hol:sec:physics}
\TA\ \emph{Holonomy sensing with a quantum probe.} A spin-$\tfrac12$ probe acted on by $U=\chi(q)$ is the simplest finite representation in which non-Abelian loop data have an operational meaning. For a known reference loop and a known candidate perturbation, \cref{qgs:hol:thm:frontier} specifies the best possible distinguishability under an explicit query budget. It yields three concrete uses. First, a symbolic gauge calculation can report not just ``the loop is nontrivial'' but the valuation of the optimum deterministic signal and the optimum conclusive event rate. Second, \cref{qgs:hol:thm:precision} tells a truncation algorithm how many transverse coefficient scales must be retained to preserve that prediction. Third, \cref{qgs:thm:cost} and \cref{qgs:hol:cor:calibration} provide a calibration requirement before interpreting postselection as an amplifier. A change of gauge frame conjugates $U$ while transforming the preparation and measurement accordingly; the optimized scalar values in the theorems are unchanged. The attainability construction assumes the two gates are known and that their eigenbasis can be used in the preparation. It is not a universal estimator for an unknown orientation, and it does not solve a gauge-reference alignment problem. Those are distinct physical tasks.

\emph{Noncommuting control and error isolation.} A four-step control sequence $aba^{-1}b^{-1}$ is a finite way of isolating a noncommuting component. The surquaternion report already proves the filtered Lie-bracket structure of such operations \cite{repoQuaternion}; the contribution here is a measurement resource assigned to the resulting group element. The square-root chart \eqref{qgs:hol:eq:Q} keeps the group operations exactly unitary without choosing a global exponential. Non-Abelian holonomies can implement quantum gates in suitable Hamiltonian families \cite{ZanardiRasetti}. We do not prove that an arbitrary pair $Q(u_1),Q(u_2)$ arises as an adiabatic holonomy of a given apparatus. Instead, once a finite unitary model is supplied, the present results certify what a probe can resolve. This is particularly relevant to nearly commuting controls: cancellation in $u_1\times u_2$ makes a naive norm-only estimate of the residual gate insufficient.

\emph{Postselection and metrological accounting.} Conditional amplification is not a free gain in information. The exact tradeoff here is a finite binary-discrimination statement, not a Fisher-information theorem for a continuous unknown parameter. It says what happens when the full successful probability is included. The wider literature on weak-value metrology also emphasizes resource accounting \cite{FerrieCombes}, but practical conclusions depend on detector saturation, technical noise, and the adopted cost model; critiques explicitly contest the scope of blanket negative claims \cite{Vaidman}. The theorem does not adjudicate that experimental debate. Within the present model, all preparations are counted, failures are retained, the hypotheses have equal prior probabilities, and allowed instruments are finite Kraus maps. Under those hypotheses the optimum is an identity, not an asymptotic heuristic. A real apparatus with a different bottleneck may still have a good engineering reason to postselect; a proof of such an advantage would need to include that bottleneck in the model.

\subsection{Real specializations}\label{qgs:hol:sec:real}
\begin{proposition}[\TC\ No faithful real evaluation of the full scalar field]\label{qgs:hol:prop:noeval}
If $F\supseteq\R$ contains a positive infinitesimal, there is no ordered-field homomorphism $F\to\R$ fixing $\R$.
\end{proposition}
\begin{proof}
Such a homomorphism is injective. The positive infinitesimal would map to a real $r>0$ satisfying $r<1/n$ for every ordinary positive integer $n$, which is impossible.
\end{proof}

\TA\ Standard part is instead a homomorphism only on $\OO$, with a nontrivial kernel. This explains why a literal surreal prediction and a family of ordinary real predictions must not be identified without explanation. Nevertheless, the finite formulas proved here are exact ordinary identities as well. One can specialize the finitely many chart parameters to positive real values and use the resulting ordinary quantum instruments; a whole Hahn field need not be evaluated to do this. The infinitesimal notation then serves as a rigorous organization of an asymptotic family, not as a measured new kind of number. For the optimal $m$-query formula one must stay on its small-angle branch: in ordinary real notation a sufficient condition is
\[
 0<w<1-\cos\!\left(\frac{\pi}{2m}\right).
\]
This condition holds eventually in the real family below for every fixed $m$. We do not claim that the same Chebyshev branch remains optimal for arbitrary real parameter values or for a query count growing without control.

\TC\ For the rank-two example, use the ordinary real functions $\epsilon=h$, $\zeta=e^{-1/h}$, $h>0$, and the same finite expressions for $u_1,u_2,w$. No surreal exponential is being used: this is an ordinary real comparison family. As $h\downarrow0$,
\[
 w(h)\sim2h^2e^{-2/h},\qquad
 w_m(h)\sim2m^2h^2e^{-2/h}.
\]
For $p(h)=h^\varrho$ with fixed $\varrho>0$, the exact formula yields
\[
 \TD_{m,\max}(p(h))\sim2m h^{1-\varrho/2}e^{-1/h}.
\]
The signal remains smaller than every fixed power of $h$ after any such power-law postselection. At $p=w_m$, conditional perfect discrimination is possible, but the event rate is correspondingly exponentially suppressed. These statements follow from the displayed finite rational identities and the elementary real estimate $e^{-1/h}=o(h^N)$ for every fixed $N$; they do not require convergence of an arbitrary formal transseries. \TA\ This is a concrete positive application: surreal or Hahn scales organize nonperturbative resource costs while finite ordinary formulas remain available for numerical or laboratory modelling. Whether the chosen real family describes a particular physical system is a separate question.

\section{A nearly dark port that heralds a finite unitary}\label{qgs:sec:dark}
\status{\TC\ Source 08, over an arbitrary real closed $F$ with $K=F[i]$, except the Hahn precision statement. \TA\ The controlled implementation and its relative path phase are hypotheses; a channel specified only up to an uncontrolled global phase is not sufficient experimental data.}

\subsection{The controlled-loop experiment}
Assume a calibrated controlled-$U$ gate is available, where $U=\chi(c)\in\SU(2,K)$. Prepare an ordinary path qubit in $\ket+=(\ket0+\ket1)/\sqrt2$, apply $U$ to the spin only on path $1$, and measure the path in the $\ket\pm$ basis. The spin Kraus operators are
\begin{equation}\label{qgs:ops:eq:ports}
 M_+=\frac{I+U}{2},\qquad M_-=\frac{I-U}{2}.
\end{equation}
A direct expansion gives $M_+^*M_++M_-^*M_-=I$ for every unitary $U$. These two operations therefore form a finite instrument in the sense of \cref{qgs:sec:instruments}.

\begin{theorem}[\TC\ Exact heralded unitary from a loop]\label{qgs:ops:thm:dark}
For a unit quaternion $c$ and $U=\chi(c)$, put
\begin{equation}\label{qgs:ops:eq:dark-prob}
 p_- =\frac{1-\Rea c}{2}.
\end{equation}
Then
\begin{equation}\label{qgs:ops:eq:dark-effect}
 M_-^*M_-=p_-I.
\end{equation}
If $c\ne1$, then $p_->0$ and
\begin{equation}\label{qgs:ops:eq:dark-unitary}
 V=\frac{I-U}{2\sqrt{p_-}}
   =\chi\!\left(\frac{1-c}{2\sqrt{p_-}}\right)\in\SU(2,K).
\end{equation}
For every density matrix, success has the same probability $p_-$ and its conditional spin state is exactly $V\rho V^*$. If $c=q_1q_2q_1^{-1}q_2^{-1}$ with $q_i=s_i+u_i$, then
\begin{equation}\label{qgs:ops:eq:dark-cross}
 p_-=\abs{u_1\times u_2}^2.
\end{equation}
For a nontrivial infinitesimal commutator, $v(p_-)=2v(c-1)=2\kappa$.
\end{theorem}
\begin{proof}
The quaternion representation gives $U+U^*=2(\Rea c)I$. Expanding $M_-^*M_-$ yields $(2I-U-U^*)/4=p_-I$. If $c\ne1$, the positive norm identity $N(1-c)=2(1-\Rea c)$ shows $p_->0$. The quaternion in \eqref{qgs:ops:eq:dark-unitary} has norm one, since $N(1-c)=4p_-$. Thus $V$ is special unitary, not just a scalar-normalized matrix of norm one on a particular input. Its Kraus expression is $M_-=\sqrt{p_-}V$, proving the probability and conditional-state claims. The commutator formulas follow from \cref{qgs:ops:thm:commutator}.
\end{proof}

\TA\ The state-independent effect \eqref{qgs:ops:eq:dark-effect} is stronger than a generic postselection statement: the port heralds a unitary operation on \emph{every} spin input and on inputs entangled with a finite ancilla. It is a finite $\SU(2)$ fact. Its surreal content is the exact separation of its disappearing success scale from its surviving conditional action.

\begin{remark}[\merge\ Consistency with the optimal identification rate]\label{qgs:mrg:rem:darkrate}
In the affine chart $q_i=Q(u_i)$ of \cref{qgs:hol:cor:comm}, the dark-port probability is $p_-=w/2$, while the optimal one-query unambiguous identification probability of \cref{qgs:hol:thm:frontier} is $w_1=w$. There is no conflict. The dark port is a heralded-gate device: with the identity it never clicks, so used as an unambiguous discriminator between $I$ and $U$ with equal priors it succeeds with probability $w/4<w_1$. It is not claimed to be an optimal discriminator; its point is the exact unitary it heralds.
\end{remark}

\begin{corollary}[\TE\ A nontrivial conditional shadow of an identity-shadow loop]\label{qgs:ops:cor:dark-limit}
For the perpendicular-axis example \eqref{qgs:ops:eq:leading-comm}, with $a=t^\alpha,b=t^\beta$ and $\alpha,\beta>0$,
\begin{equation}\label{qgs:ops:eq:dark-example}
 \st U=I,\qquad
 p_-=\frac{t^{2\alpha+2\beta}}
              {(1+t^{2\alpha})(1+t^{2\beta})},\qquad
 \st V=-\chi(\qk).
\end{equation}
Thus the conditional spin operation has a finite nonidentity shadow, while its outcome probability is infinitesimal of valuation $2(\alpha+\beta)$ and leading coefficient one.
\end{corollary}
\begin{proof}
The probability is \eqref{qgs:ops:eq:axis-trace} divided by two. Its positive square root has leading term $t^{\alpha+\beta}$. Equation \eqref{qgs:ops:eq:leading-comm} gives leading numerator $-2t^{\alpha+\beta}\chi(\qk)$ in \eqref{qgs:ops:eq:dark-unitary}. Normalize and take the residue.
\end{proof}

\TA\ In the chosen representation $\chi(\qk)=\bigl(\begin{smallmatrix}0&i\\i&0\end{smallmatrix}\bigr)$. Consequently the successful output of a spin prepared in $\ket0$ has ordinary shadow $\ket1$ up to phase. This is an order-one conditional change produced by a loop whose unconditional unitary shadow is the identity. It is not a deterministic spin flip at an infinitesimal control cost: almost all trials exit the bright port, and the rate of successful heralds carries the squared commutator scale. Under a common gauge conjugation $c\mapsto gcg^{-1}$, the probability is unchanged and $V\mapsto\chi(g)V\chi(g)^{-1}$. A calibrated input or spin analyzer transforms accordingly. The scalar event probability is gauge invariant; an axis-resolved spin readout also uses that reference data.

\subsection{Calibration error and incoherent backgrounds}
The heralded gate divides an amplitude by a small norm. Its operational interpretation must therefore include a precision and noise test.

\begin{proposition}[\TC\ Amplitude error contract for the dark port]\label{qgs:ops:prop:gate-precision}
Over a Hahn field, suppose $c-1\ne0$ is infinitesimal with $\kappa=v(c-1)$, and replace the unitary $U$ by another unitary $U'$. If $v(U'-U)>\kappa$, then for every fixed limited input density, the dark-port conditional output has the same ordinary shadow under $U'$ as under $U$. The leading probability coefficient and valuation are also unchanged.
\end{proposition}
\begin{proof}
The Kraus change is $M'_--M_-=-(U'-U)/2$. For the original input factors $R=M_-\rho^{1/2}$, \cref{qgs:ops:thm:dark} gives $\tr RR^*=p_-$, so \cref{qgs:lem:squares} gives $v(R)=\kappa$. The change in $R$ has valuation strictly larger than $\kappa$ because $\rho^{1/2}$ is limited. Apply \cref{qgs:ops:cor:amplitude-precision}.
\end{proof}

\TA\ The strict inequality is a sufficient calibration condition, not a complete noise model. A parasitic amplitude at the same order can change the normalized rotation axis. A stronger parasitic amplitude can determine the entire conditional limit.

For an explicitly incoherent contamination model, let the desired unnormalized output be $p\rho_{\rm d}$, where $\rho_{\rm d}$ is a density matrix, and let false heralds add $b\rho_{\rm f}$ with $\rho_{\rm f}$ another density matrix. Assume $p,b>0$ and $p+b\le1$. This defines the normalized accepted state
\begin{equation}\label{qgs:ops:eq:noise-mix}
 \rho_{\rm acc}=\frac{p\rho_{\rm d}+b\rho_{\rm f}}{p+b}.
\end{equation}
The model describes addition of alternative accepted events, not coherent addition of an unknown field amplitude.

\begin{proposition}[\TC\ Exact background threshold]\label{qgs:ops:prop:background}
Over a Hahn field, let $\mu=v(p)$ and $\lambda=v(b)$, with positive leading coefficients $c_p,c_b$. Then
\begin{equation}\label{qgs:ops:eq:background-cases}
 \st\rho_{\rm acc}=
 \begin{cases}
  \st\rho_{\rm d},&\lambda>\mu,\\[2pt]
  \dfrac{c_p\st\rho_{\rm d}+c_b\st\rho_{\rm f}}{c_p+c_b},&\lambda=\mu,\\[7pt]
  \st\rho_{\rm f},&\lambda<\mu.
 \end{cases}
\end{equation}
For the commutator gate, the probability threshold is $\mu=2v(c-1)$.
\end{proposition}
\begin{proof}
Divide numerator and denominator of \eqref{qgs:ops:eq:noise-mix} by the monomial of the smaller valuation. Positive leading coefficients ensure that the denominator has nonzero residue. Reduction gives the three cases, with no cancellation. Use \cref{qgs:ops:thm:dark} for the final identification.
\end{proof}

\TA\ Thus there are two distinct experimental smallness requirements: a coherent amplitude error must be below the commutator amplitude scale, while an incoherent false-herald probability must be below its squared scale to leave the ideal conditional shadow unchanged. Apparatus whose ordinary noise rate stays nonzero as the control scale tends to zero will not exhibit the ideal rare-event limit. This is a derived limitation of the declared contamination model, not a verdict that all postselection is useless in realistic metrology.

\subsection{How this connects the mechanisms}
\TA\ The matrix $M_-\rho^{1/2}$ is the Gram factor of the accepted quantum branch. Normalizing it discards a vanishing scalar but retains a finite rotation. The matrix $BD^{-1/2}$ is the Gram factor controlling the soft-mode defect of \cref{qgs:sec:schur}. Normalizing the coupling by a vanishing stiffness retains a finite relaxed-energy correction. Both constructions require positive Gram algebra, and both are invisible after the relevant leading factor has been replaced by zero. Their costs are different: rare samples and calibration in the first, unrestricted displacements and a narrow spectral regime in the second. The analogy is exact at the algebraic level but is not an identification of the physical resources.

\section{Hahn-valued gauge fields on an ordinary manifold}\label{qgs:sec:gauge}
\status{\TC\ Source 03 (Sections 9--12 there). The coefficient field is a Hahn field $\R((t^\Gamma))$ of arbitrary rank. \TI\ Chern--Weil integrality of the ordinary charge, the elliptic deformation complex of anti-self-dual connections and its Hodge theory are imported \cite{AHS,Itoh}; Kuranishi's construction is classical. \TA\ The proposed contribution is the arbitrary-support evaluation, uniqueness without a cofinality assumption, and the valuation-resolved obstruction and action statements.}

\subsection{The coefficient space is part of the hypothesis}
Let $M$ be a closed, oriented, smooth Riemannian four-manifold, and let $P\to M$ be a fixed ordinary principal $\SU(2)$ bundle. Its adjoint bundle is an ordinary real rank-three bundle with fiber the imaginary quaternions. For an ordinary vector bundle $E$ write
\begin{equation}
 \Omega^k(M,E)((t^\Gamma))
 =\left\{\sum_{\gamma\in\Sigma}a_\gamma t^\gamma:
 \Sigma\text{ well ordered},\ a_\gamma\in\Omega^k(M,E)\right\}.
\end{equation}
A positive-support element has $\Sigma\subset\Gamma_{>0}$. The support here is \emph{global}: $a_\gamma$ is a whole smooth section, and it is zero only if it vanishes identically. A statement about these series is stronger than a pointwise assertion that every value happens to be infinitesimal.

Fix an ordinary smooth connection $A_0$. A Hahn connection in its positive-support neighbourhood is written
\begin{equation}\label{qgs:obs:eq:connection}
 A=A_0+a,\qquad
 a\in\Omega^1(M,\ad P)((t^\Gamma)),\quad v(a)>0.
\end{equation}
\TA\ This notation describes coefficient extension of the connection forms on the fixed ordinary bundle. It does not construct a new manifold with surreal coordinate charts or a principal bundle whose base is the whole surreal class.

Use the two-dimensional complex representation $\chi$ to define matrix products and traces, with skew-Hermitian traceless matrices for the Lie algebra. Then
\begin{equation}\label{qgs:obs:eq:curvature}
 F_A=F_{A_0}+D_{A_0}a+a\wedge a,
 \qquad D_{A_0}a=da+A_0\wedge a+a\wedge A_0.
\end{equation}
Every coefficient is a finite sum, and support control is preserved by wedge products, brackets, the ordinary Hodge star, and ordinary covariant derivatives. Integration over $M$ is defined coefficientwise. A real-linear operator on smooth sections, including an ordinary Green operator when available, cannot create new exponents; it may annihilate some coefficients and thereby increase valuation. The finite-rank formulas $F_A=dA+A\wedge A$ and $d_AF_A=0$ over coherent coefficients are also developed in the vector-and-tensor-fields report \cite{repoVTF}.

\subsection{A valuation on curvature and action}
For a nonzero smooth-coefficient series $b$ let $v(b)$ be its least exponent with a nonzero global coefficient. Define its squared $L^2$ norm by
\begin{equation}\label{qgs:obs:eq:L2}
 \norm b_{L^2}^2=-\int_M\tr_2(b\wedge *b).
\end{equation}
For Lie-algebra-valued forms this is the usual positive inner product extended coefficientwise. In particular, if $b=t^\gamma b_\gamma+o_v(t^\gamma)$, then
\begin{equation}\label{qgs:obs:eq:L2lead}
 \norm b_{L^2}^2
 =t^{2\gamma}\norm{b_\gamma}_{L^2(\R)}^2+o_v(t^{2\gamma}),
 \qquad \norm{b_\gamma}_{L^2(\R)}^2>0.
\end{equation}
The positivity follows because a nonzero smooth section is nonzero on an ordinary open set of positive volume. Thus no cancellation can erase the leading squared norm. This is the gauge-theoretic counterpart of the leading-Gram positivity in \cref{qgs:thm:rare}.

Define a signed topological charge and the Yang--Mills action by
\begin{equation}\label{qgs:obs:eq:chargeaction}
 k(A)=-\frac1{8\pi^2}\int_M\tr_2(F_A\wedge F_A),
 \qquad \SYM(A)=\frac1{g^2}\norm{F_A}_{L^2}^2,
\end{equation}
where $g>0$ is an ordinary real coupling. This sign convention makes self-dual nonzero curvature have positive $k$ and anti-self-dual nonzero curvature have negative $k$. It is the opposite of some conventions for the second Chern number; only the displayed definition is used below. \TI\ For an ordinary connection, $k(A_0)$ is an integer by ordinary Chern--Weil theory. Its integrality is an imported topological fact, not derived from the Hahn arithmetic \cite{AHS,Itoh}.

\subsection{Topological rigidity and a sharp action-defect law}
\begin{theorem}[\TC\ Exact charge rigidity on a fixed bundle]\label{qgs:obs:thm:chern}
For the connection \eqref{qgs:obs:eq:connection},
\begin{equation}\label{qgs:obs:eq:chern-rigid}
 k(A)=k(A_0)\in\Z
\end{equation}
as an exact Hahn equality, not merely after standard part. More precisely,
\begin{align}\label{qgs:obs:eq:transgression}
 \tr_2(F_A\wedge F_A-F_{A_0}\wedge F_{A_0})
 =d\tr_2\left(2a\wedge F_{A_0}
 +a\wedge D_{A_0}a+\frac23a\wedge a\wedge a\right).
\end{align}
\end{theorem}
\begin{proof}
Introduce an auxiliary ordinary polynomial parameter $x$ and set $A_x=A_0+xa$. Its curvature is $F_x=F_{A_0}+xD_{A_0}a+x^2a\wedge a$. The Bianchi identity $D_{A_x}F_x=0$ follows by the usual algebraic expansion of $F_x=dA_x+A_x\wedge A_x$, valid coefficientwise. Using the graded cyclic trace and $\partial_xF_x=D_{A_x}a$ gives
\[
 \partial_x\tr_2(F_x\wedge F_x)
 =2\tr_2(D_{A_x}a\wedge F_x)
 =2d\tr_2(a\wedge F_x).
\]
Integrate this polynomial identity from $x=0$ to $x=1$. The integrals of $1,x,x^2$ produce the coefficients $2,1,2/3$ in \eqref{qgs:obs:eq:transgression}. The difference $a$ of two connections is a global adjoint-valued form, so the traced transgression is a global three-form. All its Hahn coefficients are finite sums of ordinary smooth expressions. Since $M$ has no boundary, ordinary Stokes' theorem makes the integral of each exact coefficient vanish. This proves \eqref{qgs:obs:eq:chern-rigid}.
\end{proof}

The transgression argument works for any support-controlled connection difference for which the displayed products exist, not only for positive-support differences. Positivity is retained because it defines a deformation with fixed ordinary residue and makes the asymptotic interpretation transparent.

\TA\ The theorem excludes fractional or infinitesimal changes of this charge within the specified fixed-bundle model. It does not prohibit ordinary instanton sectors, a sum over bundles in a path integral, tunnelling in a larger dynamical setting, boundary flux, or a singular limit that leaves the space of smooth global connections. Those processes change a hypothesis that the theorem keeps fixed.

\begin{theorem}[\TC\ Exact action gap and its valuation]\label{qgs:obs:thm:action}
Let $F_A^\pm=(F_A\pm *F_A)/2$, and let $k=k(A_0)$. Then
\begin{align}\label{qgs:obs:eq:duality-identities}
 8\pi^2k&=\norm{F_A^+}_{L^2}^2-\norm{F_A^-}_{L^2}^2,\\
 \SYM(A)&=g^{-2}\bigl(\norm{F_A^+}_{L^2}^2+
                         \norm{F_A^-}_{L^2}^2\bigr).\nonumber
\end{align}
If $k\ge0$, set $F_{\rm bad}=F_A^-$; if $k\le0$, one may instead set $F_{\rm bad}=F_A^+$. At $k=0$ either choice works. Then
\begin{equation}\label{qgs:obs:eq:action-gap}
 \SYM(A)-\frac{8\pi^2\abs k}{g^2}=\frac2{g^2}\norm{F_{\rm bad}}_{L^2}^2.
\end{equation}
When $F_{\rm bad}\ne0$,
\begin{equation}\label{qgs:obs:eq:action-value}
 v\left(\SYM(A)-\frac{8\pi^2\abs k}{g^2}\right)
 =2v(F_{\rm bad}).
\end{equation}
If $F_{\rm bad}=t^\gamma f_\gamma+o_v(t^\gamma)$, the leading coefficient of the action gap is $2\norm{f_\gamma}_{L^2(\R)}^2/g^2>0$.
\end{theorem}
\begin{proof}
The Hodge star on ordinary two-forms has eigenvalues $\pm1$, and the two eigenspaces are orthogonal. Extend these finite linear identities coefficientwise. The norm square of $F_A$ is the sum of the two norm squares, whereas $-\int\tr_2(F_A\wedge F_A)$ is their difference. \Cref{qgs:obs:thm:chern} identifies the latter with $8\pi^2k$. These are \eqref{qgs:obs:eq:duality-identities}; subtraction gives \eqref{qgs:obs:eq:action-gap}. Finally apply \eqref{qgs:obs:eq:L2lead}. The coupling and the factor two are ordinary nonzero reals, so they do not change the valuation.
\end{proof}

\TA\ The usual topological action bound is therefore not merely inherited as an inequality. It acquires a sharp order-of-smallness invariant. Near an ordinary (anti-)self-dual background, the first nonzero curvature component of the wrong duality determines both the first nonzero action correction and its strictly positive coefficient. Its abelian analogue is the exact doubling of the Maxwell energy scale, \texttt{vtf:mink:thm:energyscale} of the vector-and-tensor-fields report \cite{repoVTF}.

\TC\ The statements are gauge invariant under support-controlled unitary gauge transformations. Curvature transforms by conjugation and its traced quadratic expressions are unchanged. At the leading coefficient, the ordinary residue of a unitary gauge transformation acts by a pointwise isometry of the adjoint bundle. Thus it cannot turn a nonzero leading curvature section into zero or change its $L^2$ norm.

\begin{example}[\TE\ An infinitesimal non-Abelian field with zero charge]\label{qgs:obs:ex:torus}
On the unit-volume flat four-torus with the trivial bundle, let
\begin{equation}
 A=t^\alpha \qi\,dx^1+t^\beta \qj\,dx^2,
 \qquad \alpha,\beta>0.
\end{equation}
The coefficients are constant, but they do not commute. Since $[\qi,\qj]=2\qk$,
\begin{equation}
 F_A=2t^{\alpha+\beta}\qk\,dx^1\wedge dx^2.
\end{equation}
This curvature is nonzero, so the connection is not gauge equivalent to a flat connection. Nevertheless $F_A\wedge F_A=0$, because the decomposable form $dx^1\wedge dx^2$ wedges with itself to zero. Hence $k(A)=0$ exactly. Since $\chi(\qk)^2=-I_2$,
\begin{equation}\label{qgs:obs:eq:torus-action}
 \SYM(A)=\frac8{g^2}t^{2(\alpha+\beta)}.
\end{equation}
This is genuine non-Abelian curvature with a positive infinitesimal action, not a fractional instanton number. Its connection and curvature have zero ordinary residue. The example works for arbitrary relative ranks of $\alpha$ and $\beta$; it is the continuum analogue of the plaquette of \cref{qgs:prop:plaquette}.
\end{example}

\subsection{A support-controlled Hahn Kuranishi construction}\label{qgs:obs:sec:kuranishi}
We now address an actual nonlinear field equation, not only a topological identity. Let $M$ be connected and let $A_0$ be an ordinary irreducible anti-self-dual $\SU(2)$ connection, so $F_{A_0}^+=0$. A perturbation $a$ in background Coulomb gauge satisfies
\begin{equation}\label{qgs:obs:eq:ASD}
 d_{A_0}^*a=0,
 \qquad d_{A_0}^+a+(a\wedge a)^+=0.
\end{equation}
Here $d_{A_0}^+=\operatorname{pr}_+D_{A_0}$ on one-forms. Its linearization is the familiar elliptic deformation complex. \TI\ The identification of infinitesimal deformations and obstruction spaces is classical \cite{AHS,Itoh}: the first cohomology is represented by the harmonic kernel of the Coulomb linearization, while the second records obstructions. We require the following precise ordinary splitting.

\begin{assumption}[\TI\ Hodge--Green splitting]\label{qgs:obs:ass:hodge}
Let
\begin{equation}
 \Lop=d_{A_0}^*\oplus d_{A_0}^+:
 \Omega^1(M,\ad P)\longrightarrow
 \Omega^0(M,\ad P)\oplus\Omega^{2,+}(M,\ad P).
\end{equation}
Put $\mathcal H^1=\ker\Lop$ and $\mathcal H^2=\ker(d_{A_0}^+)^*$, both finite dimensional over $\R$. Let $\pi_1,\pi_2$ be the ordinary $L^2$ orthogonal projections onto these spaces and define $\Pi(f,b)=(0,\pi_2b)$. Assume an ordinary real-linear operator $\Green$ taking smooth pairs to smooth one-forms satisfies
\begin{equation}\label{qgs:obs:eq:Q}
 \Lop\Green=I-\Pi,\qquad \Green\Lop=I-\pi_1,\qquad \pi_1\Green=0.
\end{equation}
The irreducibility assumption includes $\ker d_{A_0}|_{\Omega^0}=0$.
\end{assumption}

\TA\ These are the standard smooth Hodge identities for the compact elliptic problem. One obtains $\Green$ from the Green operator for $\Lop\Lop^*$; irreducibility removes the degree-zero cokernel, and the remaining cokernel is $\mathcal H^2$. Elliptic regularity and finite dimensionality here are ordinary analytic inputs. We do not claim to prove them by changing the scalar field. Conversely, the theorem below only needs \eqref{qgs:obs:eq:Q} and those finite-dimensional spaces, so its dependencies are explicit and separable from its coefficient construction.

All these maps extend to Hahn series coefficientwise. In particular,
\begin{equation}\label{qgs:obs:eq:Q-value}
 v(\Green b)\ge v(b),\qquad v(\pi_jb)\ge v(b)
\end{equation}
when the corresponding arguments are defined. No norm bound uniform over all surreal values is being assumed: the inequalities follow simply because an ordinary linear map cannot introduce an exponent absent from its input.

\begin{theorem}[\TC\ Hahn Kuranishi map]\label{qgs:obs:thm:kuranishi}
Under \cref{qgs:obs:ass:hodge}, for every
\begin{equation}
 u\in\mathcal H^1((t^\Gamma)),\qquad v(u)>0,
\end{equation}
there is a unique positive-support solution of
\begin{equation}\label{qgs:obs:eq:K-fixed}
 a(u)=u-\Green\bigl(0,(a(u)\wedge a(u))^+\bigr).
\end{equation}
If $\Sigma$ is the union of the positive supports of the coordinates of $u$ in an ordinary harmonic basis, then
\begin{equation}\label{qgs:obs:eq:K-support}
 \supp a(u)\subseteq\Sigma^*\setminus\{0\},\qquad
 a(u)-u=\Oval{2v(u)}
\end{equation}
for $u\ne0$; for $u=0$ the solution is $a=0$. Define the finite-dimensional Hahn obstruction
\begin{equation}\label{qgs:obs:eq:kappa}
 \ob(u)=\pi_2\bigl(a(u)\wedge a(u)\bigr)^+
            \in\mathcal H^2((t^\Gamma)).
\end{equation}
Then
\begin{equation}\label{qgs:obs:eq:K-identities}
 \pi_1a(u)=u,\qquad d_{A_0}^*a(u)=0,\qquad
 F_{A_0+a(u)}^+=\ob(u).
\end{equation}
Consequently $A_0+a(u)$ is anti-self-dual exactly when $\ob(u)=0$. Every positive-support solution of the Coulomb ASD equations \eqref{qgs:obs:eq:ASD} arises uniquely in this way from its harmonic projection.
\end{theorem}
\begin{proof}
Choose an ordinary basis $e_1,\ldots,e_h$ of $\mathcal H^1$ and first let $u=\sum_jz_je_j$ with formal commuting variables $z_j$. Write $a=\sum_{n\ge1}a_n$, where $a_n$ is homogeneous of total degree $n$. Equation \eqref{qgs:obs:eq:K-fixed} forces
\begin{equation}\label{qgs:obs:eq:K-recursion}
 a_1=u,\qquad
 a_n=-\Green\left(0,\sum_{r=1}^{n-1}(a_r\wedge a_{n-r})^+\right)
 \quad(n\ge2).
\end{equation}
This recursively defines a unique formal series. At every fixed degree there are finitely many monomials in the $h$ variables and each coefficient is an ordinary smooth form, since $\Green$ takes smooth inputs to smooth outputs.

Now replace each $z_j$ by its given positive-support Hahn coordinate. \Cref{qgs:obs:lem:support} proves strong evaluation in the space of smooth forms, gives the support bound, and shows that the evaluated series satisfies \eqref{qgs:obs:eq:K-fixed}. The degree-one part is $u$; every remaining term has at least two positive-valuation input factors. This proves the valuation bound in \eqref{qgs:obs:eq:K-support}.

Suppose $a,b$ are two positive-support solutions for the same $u$, and put $c=a-b$. Their difference is $c=-\Green\bigl(0,(c\wedge a+b\wedge c)^+\bigr)$. If $c\ne0$, both summands on the right have valuation strictly above $v(c)$; \eqref{qgs:obs:eq:Q-value} cannot lower it. This contradicts equality. Thus the positive-support solution is unique, without invoking topological convergence of successive approximations.

Apply $\pi_1$ to \eqref{qgs:obs:eq:K-fixed} and use $\pi_1\Green=0$ to obtain $\pi_1a=u$. Apply $\Lop$ and use $\Lop u=0$ and $\Lop\Green=I-\Pi$. With $b=(a\wedge a)^+$ one obtains $\Lop a=-(I-\Pi)(0,b)=(0,-b+\pi_2b)$. The degree-zero-form component says $d_{A_0}^*a=0$; the self-dual two-form component gives $d_{A_0}^+a+(a\wedge a)^+=\pi_2b=\ob(u)$, proving \eqref{qgs:obs:eq:K-identities}.

Conversely, let $a$ be a positive-support Coulomb ASD solution. Then $\Lop a=(0,-(a\wedge a)^+)$. Applying $\Green$ and using $\Green\Lop=I-\pi_1$ gives $a=\pi_1a-\Green\bigl(0,(a\wedge a)^+\bigr)$. Set $u=\pi_1a$ and use the uniqueness just proved. The harmonic projection has positive support or is zero, so it is in the stated parameter domain.
\end{proof}

\begin{corollary}[\TC\ Unobstructed infinitesimal instanton families]\label{qgs:obs:cor:regular}
If $\mathcal H^2=0$, every positive-support harmonic parameter $u$ lifts to a unique Coulomb ASD connection $A_0+a(u)$. Its charge and action are exactly
\begin{equation}
 k(A_0+a(u))=k(A_0),\qquad
 \SYM(A_0+a(u))=\frac{8\pi^2\abs{k(A_0)}}{g^2}.
\end{equation}
The assertion holds for all ranks of the set-sized exponent group and all well-ordered positive supports of the finitely many harmonic coordinates.
\end{corollary}
\begin{proof}
The obstruction space is zero, so \eqref{qgs:obs:eq:K-identities} is the ASD equation. Use \cref{qgs:obs:thm:chern,qgs:obs:thm:action} for the two exact invariants.
\end{proof}

\TA\ This provides exact multiscale formal instanton deformations, not additional topological sectors. It also does not create an ordinary regular instanton background from nothing. Such a background and its ordinary elliptic data must be supplied; the theorem then extends its deformation problem to the specified Hahn parameters.

\begin{proposition}[\TC\ No infinitesimal gauge redundancy in the slice]\label{qgs:obs:prop:slice}
Assume $\ker d_{A_0}|_{\Omega^0}=0$. Let $A_0+a$ and $A_0+b$ be positive-support connections satisfying $d_{A_0}^*a=d_{A_0}^*b=0$. If a support-controlled $\SU(2)$ Hahn gauge transformation with residue identity relates them, then that transformation is identity and $a=b$.
\end{proposition}
\begin{proof}
Suppose $g\ne I$, $\st g=I$, and write $g=I+t^\delta\xi+o_v(t^\delta)$ with $\delta>0$ and a nonzero ordinary smooth coefficient $\xi$. The unitary and determinant-one equations at leading order make $\xi$ skew-Hermitian and traceless, hence a section of $\ad P$. The inverse has leading expansion $g^{-1}=I-t^\delta\xi+o_v(t^\delta)$. Using the convention $A^g=g^{-1}Ag+g^{-1}dg$, the leading coefficient of $A^g-A$ at $t^\delta$ is $d_{A_0}\xi$; every term involving $a$ has strictly higher valuation. Apply $d_{A_0}^*$ to the difference of the two Coulomb conditions. Its coefficient at $\delta$ is $d_{A_0}^*d_{A_0}\xi=0$. Ordinary integration by parts gives $\norm{d_{A_0}\xi}_{L^2}^2=0$, so $d_{A_0}\xi=0$. The assumed degree-zero kernel vanishes, forcing $\xi=0$, a contradiction.
\end{proof}

\TA\ We do not claim a global moduli-space classification or prove here that every arbitrary Hahn connection can be put in this gauge. The proposition has the narrower content stated: within the slice, a gauge transformation whose residue is identity cannot supply a hidden identification. Ordinary stabilizers and global gauge topology are separate issues.

\subsection{Obstructions cannot be repaired by adding smaller scales}
The first nonlinear term of the Kuranishi map is
\begin{equation}\label{qgs:obs:eq:kappa2}
 \ob_2(u)=\pi_2(u\wedge u)^+.
\end{equation}
The leading-valuation formulation gives a clean no-go theorem. Unlike an unspecified assertion that a formal series is obstructed, it identifies a coefficient that no smaller scale can cancel.

\begin{theorem}[\TC\ Persistence of a quadratic instanton obstruction]\label{qgs:obs:thm:obstruction}
Let $u_0\in\mathcal H^1$ be ordinary with $\ob_2(u_0)\ne0$, and let $\gamma>0$. There is no positive-support connection correction of the form
\begin{equation}\label{qgs:obs:eq:obstructed-a}
 a=t^\gamma u_0+o_v(t^\gamma)
\end{equation}
for which $A_0+a$ is anti-self-dual. The assertion allows arbitrary well-ordered supports for the higher terms and does not require them to be integer powers of $t^\gamma$. It also does not require $a$ to be in Coulomb gauge.
\end{theorem}
\begin{proof}
By definition of the harmonic obstruction space, $\pi_2d_{A_0}^+c=0$ for every ordinary smooth one-form $c$; this identity extends coefficientwise. Therefore $\pi_2F_{A_0+a}^+=\pi_2(a\wedge a)^+$. Insert \eqref{qgs:obs:eq:obstructed-a}. All terms containing a higher-order factor have valuation strictly above $2\gamma$, so
\begin{equation}\label{qgs:obs:eq:obstruction-leading}
 \pi_2F_{A_0+a}^+
 =t^{2\gamma}\ob_2(u_0)+o_v(t^{2\gamma}).
\end{equation}
Its leading coefficient is nonzero. Thus $F_{A_0+a}^+$ cannot vanish. A linear correction at exponent $2\gamma$, even if available, lies in the image of $d_{A_0}^+$ and is annihilated by $\pi_2$; it cannot cancel the obstruction coefficient.
\end{proof}

\TA\ The theorem concerns a nonzero quadratic obstruction. If that coefficient vanishes, higher obstruction terms can interact with appropriately chosen higher corrections. We make no blanket claim that the first nonzero value of every higher homogeneous term is immutable under all choices of tails. The exact equation $\ob(u)=0$ is the correct test in that case.

\begin{theorem}[\TC\ Sharp action floor at an obstructed tangent]\label{qgs:obs:thm:floor}
Under the hypotheses of \cref{qgs:obs:thm:obstruction}, set $c=\norm{\ob_2(u_0)}_{L^2(\R)}^2>0$ and $k_0=k(A_0)\le0$. For every correction \eqref{qgs:obs:eq:obstructed-a}, let
\[
 \mathcal E(a)=\SYM(A_0+a)-\frac{8\pi^2\abs{k_0}}{g^2}.
\]
Then $\mathcal E(a)>0$ and
\begin{equation}\label{qgs:obs:eq:floor-value}
 v(\mathcal E(a))\le4\gamma.
\end{equation}
If equality holds, its leading coefficient is at least $2c/g^2$. For the corrected connection $a(t^\gamma u_0)$ supplied by \cref{qgs:obs:thm:kuranishi}, the bound is attained in the stronger form
\begin{equation}\label{qgs:obs:eq:floor-attained}
 \mathcal E\bigl(a(t^\gamma u_0)\bigr)
 =\frac{2c}{g^2}t^{4\gamma}+\Oval{5\gamma}.
\end{equation}
Thus additional smaller scales cannot produce an action defect of strictly higher valuation while retaining this leading tangent.
\end{theorem}
\begin{proof}
Since $A_0$ is ASD and charge is fixed, \cref{qgs:obs:thm:action} gives $\mathcal E(a)=2\norm{F_{A_0+a}^+}_{L^2}^2/g^2$. The ordinary orthogonal projection $\pi_2$ extends coefficientwise and the Pythagorean identity remains exact:
\[
 \norm F_{L^2}^2=\norm{\pi_2F}_{L^2}^2+\norm{(I-\pi_2)F}_{L^2}^2
 \ge\norm{\pi_2F}_{L^2}^2.
\]
Apply this to $F=F_{A_0+a}^+$ and use \eqref{qgs:obs:eq:obstruction-leading} and \eqref{qgs:obs:eq:L2lead}. The right side has valuation $4\gamma$ and leading coefficient $c$. The order inequality proves \eqref{qgs:obs:eq:floor-value}, positivity, and the asserted leading coefficient lower bound when the valuations agree.

For the constructed correction with $u=t^\gamma u_0$, the homogeneous recursion has $a=u+\Oval{2\gamma}$, and hence $F_{A_0+a(u)}^+=\ob(u)=t^{2\gamma}\ob_2(u_0)+\Oval{3\gamma}$. The action defect is exactly twice this squared norm divided by $g^2$. Its leading term is $2ct^{4\gamma}/g^2$; all cross and higher terms have valuation at least $5\gamma$. This is \eqref{qgs:obs:eq:floor-attained}.
\end{proof}

\TA\ The inequality direction in \eqref{qgs:obs:eq:floor-value} is important. A larger valuation would mean a \emph{smaller} action defect. The theorem rules that out. Other choices of higher corrections can introduce a defect earlier than $4\gamma$; the Kuranishi-corrected representative avoids those earlier components and realizes the least possible leading coefficient at the best possible valuation.

\TA\ The positive result is an exact multiscale construction of deformations whenever the obstruction map permits them. The negative result is that hierarchies of new scalar sizes do not remove a nonzero cohomological obstruction whose leading coefficient has been identified, and the action-floor theorem gives a sharp quantitative form of the same obstruction. The classical Kuranishi mechanism predates this report. What has been proved here, rather than assumed by analogy with a one-parameter formal series, is its evaluation on arbitrary positive well-ordered Hahn supports, its uniqueness without a cofinality assumption, and the resulting valuation-resolved obstruction and action statements. Whether this exact package is new in the broader literature on deformation functors and valued coefficient rings remains a priority question. It is a concrete candidate for further mathematical development, not an established physical discovery.

\subsection{A concrete two-scale torus obstruction}\label{qgs:obs:sec:torus-obstruction}
The obstruction theorem was stated around an irreducible background to share the Kuranishi setup. Its projection argument itself does not need irreducibility. On a flat torus it gives an explicit example, with no unconstructed instanton background and no unspecified obstruction coefficient.

\begin{theorem}[\TC\ Noncommuting torus directions cannot be repaired]\label{qgs:obs:thm:torus-obstruction}
On the trivial $\SU(2)$ bundle over the unit-volume flat oriented four-torus, let $0<\alpha\le\beta$ and consider
\begin{equation}\label{qgs:obs:eq:torus-tail}
 A=t^\alpha \qi\,dx^1+t^\beta \qj\,dx^2+h,
 \qquad v(h)>\beta.
\end{equation}
Then $A$ is not anti-self-dual. Its charge is zero, and its positive action satisfies $v(\SYM(A))\le2(\alpha+\beta)$. If equality holds, its leading coefficient is at least $8/g^2$. For $h=0$ the bound is attained exactly, as in \eqref{qgs:obs:eq:torus-action}. All assertions allow arbitrary well-ordered supports for $h$ and arbitrary rank of the exponent group.
\end{theorem}
\begin{proof}
At the flat background $A_0=0$, the harmonic self-dual two-forms have constant coefficients. Let $\pi_{\rm harm}$ be their ordinary $L^2$ projection, tensored with the imaginary quaternions. It annihilates $d^+h$, coefficientwise. The two square terms from the displayed leading one-forms vanish, and their cross term is $2t^{\alpha+\beta}\qk\,dx^1\wedge dx^2$. Every product involving $h$ has valuation strictly above $\alpha+\beta$, since $v(h)>\beta\ge\alpha$. With the orientation $dx^1\wedge dx^2\wedge dx^3\wedge dx^4$ this gives
\begin{equation}
 \pi_{\rm harm}F_A^+
 =t^{\alpha+\beta}\qk\,(dx^1\wedge dx^2+dx^3\wedge dx^4)
       +o_v(t^{\alpha+\beta}).
\end{equation}
Its nonzero leading coefficient has squared norm $4$: the two coordinate two-forms are orthonormal and $-\tr_2\chi(\qk)^2=2$. Thus $F_A^+\ne0$, proving non-ASD. Charge zero follows from \cref{qgs:obs:thm:chern} on the trivial bundle. The action is $2\norm{F_A^+}_{L^2}^2/g^2$; orthogonal projection cannot increase the squared norm, so the displayed harmonic component gives the valuation and leading coefficient bounds. With $h=0$, direct computation gives \eqref{qgs:obs:eq:torus-action}, attaining them.
\end{proof}

\TA\ A change at the already prescribed $t^\beta$ scale could alter the result; the theorem forbids corrections strictly smaller than that scale. This is exactly the kind of hypothesis a multiscale calculation must retain. Merely allowing more infinitesimals supplies no cancellation mechanism for the identified non-Abelian harmonic flux.

\subsection{What the gauge application does and does not establish}\label{qgs:sec:vtfrel}
\TA\ The gauge results suggest keeping the ordinary elliptic geometry and letting only the finite harmonic deformation coordinates carry hierarchical scales. The expensive analytic input is then the ordinary Hodge splitting; once it is available, the recursive nonlinear corrections are built from $\Green$, wedge products, and projections. The output distinguishes three outcomes: an unobstructed harmonic parameter gives an exact supported ASD lift; a nonzero leading obstruction gives a certificate of nonexistence for all allowed smaller-scale tails; a parameter with vanishing initial obstruction but unknown higher terms requires more computation and is not automatically solved. The action-floor theorem additionally reports the best leading scale of the unavoidable Yang--Mills defect for an obstructed tangent. This use is more specific than saying that surreal numbers might regularize a field configuration: it preserves exact charge, positivity, and a nonlinear equation while identifying the obstruction to satisfying that equation, and can explain why a perturbative ansatz fails without mistaking the failure for numerical instability. A full application to a particular ordinary instanton moduli problem would need an explicit background, computable harmonic bases, enough control of $\Green$ to calculate the desired coefficients, and an independent analysis of any limit of ordinary physical configurations. The present report proves the extension mechanism and gives a fully explicit flat-torus obstruction; it does not claim to have carried out a new classification of instantons.

\TA\ \emph{Relation to the vector-and-tensor-fields report.} Among that report's open directions (\texttt{vtf:subsec:open}) is ``support certificates for specific gauge-fixed nonlinear systems'', recorded there as partly answered for a regular metric, stationary flow on the three-torus and a cubic wave equation, with ``gauge fixing, causal propagation and the nonlinear gravitational equations, and any global Einstein, Yang--Mills or Navier--Stokes evolution statement'' left open. \Cref{qgs:obs:thm:kuranishi} is a support certificate for one further specific gauge-fixed nonlinear system: the Coulomb-gauge anti-self-duality equations on a closed Riemannian four-manifold, around an ordinary irreducible instanton, under \cref{qgs:obs:ass:hodge}. It supplies the gauge fixing for this elliptic, Riemannian, static system and conditional existence and uniqueness where that report's non-claim records that its variational and gauge equations come without existence theorems. It does \emph{not} touch the evolution part: nothing here is Lorentzian, hyperbolic or dynamical, and no Yang--Mills evolution statement is proved. That item therefore stays open, re-scoped by one more specific instance.

\section{Algebraic obstructions that scalar extension does not remove}\label{qgs:sec:nogo}

\subsection{The Bell bound survives}
\begin{theorem}[\TC\ The finite algebraic Tsirelson bound]\label{qgs:thm:bell}
Let $A_0,A_1,B_0,B_1$ be Hermitian involutions over $K$, or over $\HF$, with $[A_j,B_k]=0$ for all $j,k$, over an arbitrary real closed $F$. Put
\[
 \mathcal B=A_0(B_0+B_1)+A_1(B_0-B_1).
\]
Then $-2\sqrt2\,I\le\mathcal B\le2\sqrt2\,I$ exactly, not merely after standard part; for every state in the corresponding finite theory, $\abs{\langle\mathcal B\rangle}\le2\sqrt2$.
\end{theorem}
\begin{proof}
Set $R_0=A_0-(B_0+B_1)/\sqrt2$ and $R_1=A_1-(B_0-B_1)/\sqrt2$. Using the cross-commutation relations and involution identities gives the exact sum of squares
\begin{equation}\label{qgs:eq:BellSOS}
 2\sqrt2 I-\mathcal B=\frac1{\sqrt2}(R_0^2+R_1^2)\ge0.
\end{equation}
Changing both $A_j$ to $-A_j$ gives the other bound. Taking a positive normalized state expectation preserves these inequalities. For quaternionic matrices the same central real coefficients and algebraic identity apply, and positivity may equivalently be checked through $\chi$.
\end{proof}

\status{\TC\ Sources 04 (Proposition 3.4, over $K$) and 07 (Theorem 11.1, also over $\HF$). \TI\ The bound is classical \cite{Cirelson}; source 04 includes it only to delimit the proposal. \TA\ The proof shows why non-Archimedean coefficients cannot evade it while the same positivity and compatibility assumptions are retained: the ability to encode infinitesimal scales is not, under the stated measurement model, an ability to exceed the usual quantum correlation bound. In a finite classical hidden-variable model with nonnegative $F$-valued weights summing to one, deterministic local $\pm1$ outcomes still give the CHSH bound two by the same elementary convexity argument. Setting-dependent postselection is not covered by that classical assertion: changing the sampled ensemble changes the hypotheses, not the scalar arithmetic.}

\subsection{Finite-dimensional canonical commutation remains impossible}
\begin{proposition}[\TC\ A complex finite-matrix obstruction]\label{qgs:prop:CCR}
For $n>0$ and $0\ne\hbar\in F$, no matrices $Q,P\in M_n(K)$ satisfy $[Q,P]=i\hbar I_n$.
\end{proposition}
\begin{proof}
Over the commutative field $K$, matrix trace is cyclic, so the trace of the left side is zero. The trace of the right side is $ni\hbar\ne0$ in characteristic zero.
\end{proof}

\TA\ This remains true even when $\hbar$ is infinitesimal. The proposition concerns the usual complex central right-hand side; it is not a claim that quaternion-valued traces have the same cyclic property. A continuum oscillator or quantum field cannot be replaced by a fixed finite matrix model satisfying exact canonical commutation relations merely by extending the scalars. (Source 07.)

\subsection{Where the present results stop}\label{qgs:sec:stop}
\TA\ The finite protocol theorems do not construct an infinite-dimensional Hilbert space, a spectral measure, a countably additive probability theory, or a quantum field. The finite lattice theorem does not justify interchanging standard part with a continuum limit or with an infinite-volume limit. The implicit-series theorem does not supply a Borel summation prescription, and the finite-phase exponential does not define arbitrary infinitely rapid oscillations.

There is also a topological issue specific to the full surreal field: its fine order topology is not the ordinary real topology enlarged by a harmless set of tiny distances. The collection documents strong restrictions on small continuous paths in that topology \cite{repo,repoPhysics}. Our use of a set-sized Hahn workspace and ordinary readouts does not assert those restrictions disappear; it avoids assuming the full fine topology is the appropriate topology for physical time.

Finally, the present results do not change the earlier negative analysis of gravitational singularities. Replacing a divergent real expression by an unlimited surreal expression does not provide a regular endpoint, a continuation law, or a measurement rule \cite{repoPhysics}. The positive applications here concern controlled scale information in explicitly stated finite models, not a mechanism of singularity resolution (\cref{qgs:sec:boundaries}).

\section{Real families and beyond-all-orders scale realizations}\label{qgs:sec:realization}
\status{\TC\ Source 08, with the real families of sources 04 and 07 cross-referenced.}

\subsection{Power-scale examples are genuine ordinary asymptotics}
\TA\ An abstract Hahn series need not define a function of a real parameter. In particular, arbitrary coefficients can grow too fast for a formal series to converge at any positive real argument. The algebraic statements above do not rely on such convergence. To apply their valuations as predictions about a family of real systems, a separate realization must be specified.

\begin{proposition}[\TC\ Convergent specialization for the power-scale examples]\label{qgs:ops:prop:puiseux}
Consider a finite expression built from ordinary real constants, monomials $t^q$ with $q\in\Q$, finite arithmetic operations with nonzero denominators, and positive square roots of positive nonzero expressions. Its Hahn value has a convergent Puiseux expansion near $t=0$ on the positive real side. Substitution $t=s>0$ is defined for all sufficiently small $s$, preserves every algebraic identity among such expressions, and, for a nonzero value with leading term $at^q$, gives
\begin{equation}\label{qgs:ops:eq:real-asymptotic}
 f(s)=a s^q(1+o(1)).
\end{equation}
For a limited value, its standard part is its ordinary limit. The statement applies coordinatewise to the explicitly displayed rational-power examples in this report.
\end{proposition}
\begin{proof}
Start with convergent Laurent series in $s^{1/N}$ for a common denominator of the finitely many initial exponents. They are closed under finite sums and products. A nonzero such series factors as $as^q(1+h(s))$ with $h(s)\to0$ and convergent $h$; its reciprocal has the convergent geometric expansion after shrinking the interval. For a positive expression, $a>0$, and its positive square root is $\sqrt a\,s^{q/2}(1+h(s))^{1/2}$, with the convergent binomial expansion. Only finitely many square-root steps are performed, so a new finite common denominator suffices. Induction proves convergence and agreement with the Hahn operations. A nonzero convergent series has a first nonzero coefficient and hence \eqref{qgs:ops:eq:real-asymptotic}; if it is identically zero, uniqueness of convergent power-series coefficients gives the zero Hahn series. Positivity and denominator nonvanishing therefore hold on a sufficiently small positive interval. The limit assertion follows from the leading exponent and constant term.
\end{proof}

\TC\ For example, the Schur construction \eqref{qgs:ops:eq:realize} with $\varepsilon=s^q$ has \emph{exactly} the same relaxed matrix for every $s>0$, not merely the same limit; its minimizing displacement diverges like $s^{-q}N^*x$. The Bell success probability tends to zero like $2s^{2a}$ when $a\in\Q_{>0}$. The perpendicular-axis dark-port gate has probability asymptotic to $s^{2\alpha+2\beta}$ and conditional unitary limit $-\chi(\qk)$. These are ordinary real families whose asymptotic structure is summarized by the Hahn values.

\subsection{A rank-two hierarchy that no single power scale captures}
To see a use beyond ordinary Puiseux series, take the rank-two workspace \eqref{qgs:eq:ranktwo}, so $0<\zeta<\epsilon^N$ for every ordinary positive integer $N$. A useful ordinary family of monomials is
\begin{equation}\label{qgs:ops:eq:exp-realization}
 t^{(a,b)}\longmapsto e^{-a/s}s^b,
 \qquad s\downarrow0,
\end{equation}
under which $\epsilon\mapsto s$ and $\zeta\mapsto e^{-1/s}$. \TA\ This is a choice of asymptotic realization, not an identity between the surreal $\omega$-map and the real exponential.

\begin{proposition}[\TC\ Finite rank-two monomial realization]\label{qgs:ops:prop:rank-two}
Map a finite sum $P=\sum c_{a,b}t^{(a,b)}$ by \eqref{qgs:ops:eq:exp-realization}, combining equal exponent pairs first. For nonzero $P$ with least exponent $(a_0,b_0)$ and leading coefficient $c_0$, its realized function satisfies
\begin{equation}\label{qgs:ops:eq:rank-two-leading}
 P(s)\sim c_0 e^{-a_0/s}s^{b_0}.
\end{equation}
The map extends injectively to fractions of finite sums and preserves leading ratios and eventual signs. Positive square roots of positive such fractions have the corresponding positive leading square root. No evaluation of an arbitrary infinite Hahn sum is asserted.
\end{proposition}
\begin{proof}
After division by the least monomial, every other term has ratio $e^{-(a-a_0)/s}s^{b-b_0}$. If $a>a_0$, this tends to zero because an exponential with negative reciprocal argument dominates every fixed power; if $a=a_0$, then $b>b_0$ and the power tends to zero. There are only finitely many terms, so their sum tends to zero. Nonzero finite sums are therefore eventually nonzero with the claimed sign. The rules for products and fractions follow, proving injectivity of the fraction-field realization. Taking the ordinary positive square root of a positive asymptotic relation proves the last assertion.
\end{proof}

\TC\ For a concrete loop, choose $a=\epsilon$ and $b=\zeta$ in \eqref{qgs:ops:eq:axis-rotations} with perpendicular axes. Its exact dark-port probability is
\begin{equation}\label{qgs:ops:eq:exp-gate}
 p_- =\frac{\epsilon^2\zeta^2}{(1+\epsilon^2)(1+\zeta^2)}.
\end{equation}
Under \eqref{qgs:ops:eq:exp-realization}, it is asymptotic to $s^2e^{-2/s}$, while the conditional unitary still tends to $-\chi(\qk)$. Likewise, putting $D=\zeta^2I$ and $B=\zeta N$ gives the same finite Schur defect $NN^*$ as in \cref{qgs:sec:schur}, although both coefficients vanish beyond all algebraic powers of $s$; the static energy window now requires $\abs z\ll e^{-2/s}$ for the sufficient criterion of \cref{qgs:ops:thm:window}. The same realization is used in \cref{qgs:thm:alljets} ($\epsilon=h$, $\zeta=e^{-1/h}$), in \cref{qgs:cor:realrealization} ($\delta_1=h$, $\delta_2=e^{-1/h}$) and in \cref{qgs:hol:sec:real}.

\TA\ The associated value group makes the ordering of these scales exact before a real parameterization is chosen. A different physical parameterization could realize the same finite hierarchy. Conversely, a general arbitrary-rank Hahn workspace need not have a realization by one real variable compatible with all its infinite sums. None of the main theorems depends on that much stronger claim.

\subsection{Why this is not an empirical argument for surreal scalars}
\TA\ All explicit examples above admit ordinary real asymptotic models. The finite algebra does not select between a literal non-Archimedean theory and an ordinary parameter-dependent family. Its demonstrated benefit is structural: exact scale comparison, positivity-sensitive normalization, and explicit tests for whether reduction loses essential information. It does not show that nature contains a measurable infinitesimal, or that $\No$ rather than a much smaller ordered field is physically necessary. Every construction here lives in a set-sized workspace; none needs the full proper class of surreal numbers. For a single power-law experiment, a Puiseux field can be sufficient; a rank-two field is useful for simultaneous power and beyond-all-orders scales. The larger surreal setting provides a common language in which such workspaces can be situated, not a license to ignore their analytical or experimental interfaces.

\TA\ Most theorems in this report hold in any real-closed valued field with the stated residue and strong-evaluation properties; their physical mechanisms are not exclusive to $\No$. The surreal field is attractive because a wide range of ordered exponent hierarchies can be represented in one normal-form language, while set-sized subfields isolate the resources needed by a calculation. That universality is a representational advantage. To turn it into a uniquely surreal physical claim one would need a principle selecting a particular embedding, a particular class of allowed coefficients, or a particular readout that could not be replaced by a smaller valued-field model. No such principle is derived here. The theorem statements are deliberately invariant under ordered exponent embeddings wherever possible, which makes their mathematical content more robust but their dependence on the full surreal class less essential. (Sources 07 and 08.)

\section{Relation to the physics report}\label{qgs:sec:physrel}
\TA\ The collection's physics report, \texttt{docs/physics/surreal-scalars-and-spacetime} \cite{repoPhysics}, defines its value as a firewall and reaches a negative verdict: surreal scalars make a singular regime exactly representable, and they do not make it regular. This report is consistent with that stance. It claims representational and bookkeeping gains in finite models, never new physics, and every source states the same conclusion independently (\cref{qgs:app:nonclaims}). Its relation to specific statements of that report is as follows; the labels cited exist there at the placement commit.
\begin{itemize}
\item \emph{Confirmed and sharpened.} The status note of \texttt{phys:thm:quantumshadow} says that an outcome of infinitesimal probability has zero ordinary shadow, so reduction cannot determine the state obtained by dividing by that infinitesimal. \Cref{qgs:thm:alljets} shows that not even all one-scale power data determine it, and \cref{qgs:ops:prop:rare-embedding} that identical complete shadows can conceal any ordinary channel. The replacement datum is supplied by \cref{qgs:thm:rare}: the leading Gram matrix, which is leading data rather than reduced data, so there is no contradiction.
\item \emph{Now covered.} The status note of \texttt{phys:ex:rare} says that postselection on an infinitesimal-probability event lies outside \texttt{phys:thm:quantumshadow}. It remains outside that theorem, and is now priced by \cref{qgs:thm:cost}.
\item \emph{Answered for postselection.} The status note of \texttt{phys:ex:amplify} demands that a proposed observable effect state how a duration, coupling, postselection or amplification is physically realized, and identify the physical gain and its cost. For postselection the cost is now a theorem: the success probability, $\max(p,q)\,\TD_{\rm out}\le\TD_{\rm in}$ with equality attained (\cref{qgs:thm:cost}); the exact attained frontier for a small holonomy (\cref{qgs:hol:thm:frontier,qgs:hol:thm:valuation}); and calibration and background thresholds (\cref{qgs:hol:cor:calibration,qgs:ops:prop:gate-precision,qgs:ops:prop:background}). For duration, \cref{qgs:obs:thm:dynamics} generalizes \texttt{phys:ex:amplify} to any exactly isolated effective block but, like it, supplies no physical realization. The demand for coupling and amplification in general, and for what fixes the scale, stands.
\item \emph{Continued.} The report's register of non-claims (\texttt{phys:app:nonclaims}) excludes conditioning on standard-probability-zero events and infinite or hyperfinite experiments. This report continues the first exclusion: its rare-branch theorems compute a conditional state \emph{inside} the valued model and prove that an ordinary finite experiment cannot read it with positive ordinary probability; they do not reinstate conditioning on a standard-probability-zero event as an ordinary readout. The second exclusion stays: all resources here are ordinary finite.
\item \emph{Instances.} The finite-angle exponential of \cref{qgs:obs:lem:exponential} is an instance of the first two tractable classes in \texttt{phys:sub:dynamics}, and no global unitary evolution for an unlimited Hamiltonian is inferred.
\item \emph{Still open.} The report's realization bridge (its Stage D, \texttt{phys:sub:stageD}) remains open: \cref{qgs:cor:realrealization} and \cref{qgs:sec:realization} are finite-matrix realization bridges, not the Einstein--matter target named there.
\end{itemize}
This report is filed beside the physics report rather than merged into it because only its rare-event part touches that report, and what it continues there is a stated exclusion and an assessment-tier demand; the effective-Hamiltonian, discrimination, Schur-defect, entanglement and Yang--Mills material shares no named question with a report whose subject is gravitational singularities.

\section{Applications, benchmarks, and a decision table}\label{qgs:sec:applications}
\TA\ Everything in this section is tier \TA: these are proposals and tests, not results.

\subsection{Concrete uses}
\paragraph{Rare-event asymptotic tomography.}
Given symbolic Kraus matrices and a purification with known leading support, \cref{qgs:thm:rare} computes the state in a rare recorded branch. It can expose dark-state leakage and distinguish competing small amplitudes that disappear from the unconditional density matrix. A successful implementation must carry initial matrices and re-evaluate the order after cancellations, rather than simply add gate valuations. Its physical output must be accompanied by the branch weight and the resource bound of \cref{qgs:thm:cost}. A symbolic backend could attach a valuation and leading Gram certificate to each finite instrument branch, with the two precision contracts of \cref{qgs:ops:cor:amplitude-precision,qgs:ops:cor:distance-scale}. A useful benchmark would compare a computed leading output to an exact real family at several small parameters, keeping both the normalized fidelity and the acceptance rate; success would mean detecting branches whose apparent order-one improvement depends on omitted input or calibration data. This requires no arbitrary-Hahn equality algorithm: the test domain can be finite algebraic expressions with certified leading terms.

\paragraph{A certificate-producing asymptotic quantum calculator.}
The input would be a finite matrix or Kraus instrument whose entries are represented in an effective sublanguage of a chosen Hahn field. Each operation would produce both its mathematical value and a certificate: positive support for a formal evaluation, an invertible residue gap for a graph lift, or a positive leading Gram trace for a rare-event normalization. For a branch, the useful output is not just ``probability zero to this order'' but either exact impossibility or a record $(v(p),\lc(p),\st(X/p))$ with enough leading amplitude data to verify it. For a Hamiltonian block, the output is a leading splitting scale, a Hermitian effective residue, and a statement of which rescaled time variable makes that residue govern the dynamics; a finite-depth spectral tree then records which degeneracies are exact and which are merely unresolved at an earlier scale. A practical implementation should begin with finite lexicographic rational exponent groups and rational or algebraic ordinary coefficients. Unrestricted surreal notation, arbitrary support comparison, and equality testing must not be advertised as solved by these theorems. The advantage over a single formal parameter is clearest when a model needs scales such as $\zeta\ll\epsilon^N$ for every ordinary $N$; nevertheless multivariate formal series, transseries, and other valued fields already provide related capabilities, and the surreal viewpoint offers a common ambient interpretation, not an exclusive claim to the underlying computational idea.

\paragraph{Certified multiscale effective Hamiltonians.}
\Cref{qgs:thm:elimination} supplies an exact target for eliminating several auxiliary sectors, including sectors separated beyond all powers. The first correction is Hermitian and includes the normalization anticommutator, which is lost by treating the nonorthonormal graph generator $A+BY$ as though it were already a physical Hermitian Hamiltonian. The real-parameter counterpart supplies uniform finite-scale error bounds under explicit boundedness hypotheses. The three-level model of \cref{qgs:prop:virtual} distinguishes direct restriction from elimination through virtual paths; a generalized scale-ordering algorithm should reproduce the symmetric ground-state shadow before inspecting a smaller direct bias. More complicated models would require repeated block elimination (\cref{qgs:obs:thm:hierarchy}), not an unproved universal rule that each next coefficient is simply compressed to the preceding ground eigenspace.

\paragraph{Minimal soft-mode explanations of relaxed actions.}
Given ordinary $H_0$ and a candidate relaxed matrix $E$, compute $\Xi=S_0-E$. The classification gives a complete finite test: $0\le\Xi\le S_0$ and $\rank\Xi\le r$. When it passes, the factor $N$ in \eqref{qgs:ops:eq:realize} constructs a witness. When it fails, this particular positive soft-mode mechanism is excluded. A physical follow-up must bound allowed displacements and establish the relevant energy regime. An extension to nonlinear positive potentials would be substantially more informative than blindly enlarging the coefficient field: the analogue of the defect could depend on nonlinear saturation rather than only on a Gram factor.

\paragraph{Gauge-invariant scale diagnostics and controlled loops.}
\Cref{qgs:thm:wilson} computes the exact leading face-energy scale from initial loop curvature and positive weights; because the coefficient is a positive sum of squared norms, it detects the leading action without cancellations between faces, and offers a consistency check for anisotropic or nested lattice scale assignments. It does not establish the convergence of a family with an increasing number of faces; that would require uniform estimates absent from a fixed finite theorem. The finite two-rotation sequence of \cref{qgs:sec:dark} provides a tractable testbed for scale-sensitive gauge or spin control: its exact dark-port probability, conditional unitary, calibration threshold and background threshold are all in closed form. An implementation study would have to model control errors, available inverse gates, phase calibration, and whether preserving successful spin outputs has value relative to deterministic alternatives. A theorem of conditional unitarity alone is not a metrological advantage or an efficient gate synthesis result. \Cref{qgs:sec:discrimination} gives the matching optimal-discrimination and precision statements.

\paragraph{Multiscale instanton deformation and obstruction diagnosis.}
See \cref{qgs:sec:vtfrel}.

\paragraph{Renormalization and transseries: a direction, not a theorem here.}
Surreal normal forms can provide an ambient language for transseries-like scale hierarchies, and compatible differential structures have been constructed \cite{BM}; ordered Laurent coefficient fields already play a role in formal quantum state constructions \cite{BW}. It is therefore reasonable to investigate a Hahn-valued effective-action calculation with explicit support and renormalization data. However, none of the present theorems constructs a continuum quantum field theory, a path-integral measure, a renormalized interacting Hamiltonian, or a regulator-independent observable. Defining a divergent expression as an unlimited scalar does not perform those tasks. A credible next calculation would specify a concrete regularized ordinary theory, prove its coefficient identities in a controlled workspace, and show that an explicitly prescribed finite readout is independent of the allowed regulator changes. The physics report already explains why an arbitrary finite-part extraction is not a substitute for this work \cite{repoPhysics}.

\subsection{Formalization targets and effective computation}\label{qgs:sec:formal}
The collection's documented standard-part, strong-summation, and finite-algebra layers are plausible prerequisites for formalizing these statements. The initial-state theorem primarily needs finite sums of squared moduli, positivity, and coefficient extraction. The postselection theorem needs a finite trace norm and a failure-flag construction. The elimination theorems need matrix-valued formal implicit evaluation, invariant graphs, and positive square roots near the identity. The discrimination results need finite dilations, the overlap triangle and Chebyshev identities. The gauge theorems need a quaternion $*$-algebra, finite loop products, and the leading-coefficient rule for positive sums. This is a dependency assessment (source 07), not a claim that these modules have been added or checked. An actual Lean implementation would need to reconcile all interfaces with the repository's current definitions and verify its imports and axioms. Similarly, arbitrary Hahn series are not automatically effective data structures. A computational application needs a computable exponent group, explicit support descriptions, and certified ways to determine initial terms after cancellation. The all-power-data obstruction shows why a universal finite truncation rule cannot substitute for those requirements.

\subsection{Benchmarks that can fail}
A research proposal becomes more useful when its failure conditions are specified. \Cref{qgs:tab:benchmarks} lists tests that separate a working implementation from a relabelling of formal series (source 07, with additions from sources 04 and 08).

\begin{longtable}{@{}L{0.25\textwidth}L{0.40\textwidth}L{0.27\textwidth}@{}}
\caption{Concrete mathematical tests, not forecasts of experimental advantage.}\label{qgs:tab:benchmarks}\\
\toprule
Benchmark & Required result & Failure to detect\\
\midrule
\endfirsthead
\toprule
Benchmark & Required result & Failure to detect\\
\midrule
\endhead
\bottomrule
\endfoot
Rare rank-two branch
& Compute $p$, its leading coefficient, and the conditional projector even when $\st p=0$.
& Returning an undefined $0/0$ or replacing the branch by zero.\\[3pt]
Indistinguishable power data
& Distinguish the two conditional outputs of \cref{qgs:thm:alljets} by retaining $\zeta$-data.
& Claiming all $\epsilon$-power data suffice.\\[3pt]
Noncommuting high block
& Verify the Riccati residual and the anticommutator correction when $D\Delta\ne\Delta D$.
& Silently diagonalizing scales independently of $D$.\\[3pt]
Virtual ground selection
& Recover $\st(\lambda_{\min}/\epsilon^2)=-2$ and $\st P_{\min}=\proj{+}$.
& Selecting $\ket0$ from the direct bias.\\[3pt]
Soft-mode relaxation
& Reproduce every defect $0\le\Xi\le S_0$ of rank at most $\dim\ker D_0$, and no other.
& Discarding $B$ because $B_0=0$ on the soft subspace.\\[3pt]
Two-scale plaquette
& Recover mixed order $2(\alpha+\beta)$ and bracket norm coefficient.
& Treating non-Abelian products as commuting scale factors.\\[3pt]
Nearly aligned controls
& Recover $v(w)=8v(\epsilon)$ in \cref{qgs:hol:ex:aligned}.
& Multiplying leading control norms, or truncating $u_2$ to $\epsilon\qi$.\\[3pt]
Quaternionic composite
& Respect the $36$ versus $28$ obstruction for two full two-level systems.
& Declaring a single-system doubling map to be a tensor-product theory.\\
\end{longtable}

\subsection{A unified decision table}
The most useful output of the theory is a criterion for when an ordinary approximation is safe. \Cref{qgs:ops:tab:decision} (source 08, extended) summarizes the proved tests; its conditions should be checked \emph{before} discarding small coefficients.

\begin{longtable}{@{}L{0.22\textwidth}L{0.33\textwidth}L{0.37\textwidth}@{}}
\caption{Proved algebraic conditions, not universal claims about apparatus.}\label{qgs:ops:tab:decision}\\
\toprule
Operation & Safe reduction regime & Data or cost at the singular boundary\\
\midrule
\endfirsthead
\toprule
Operation & Safe reduction regime & Data or cost at the singular boundary\\
\midrule
\endhead
\bottomrule
\endfoot
Finite quantum circuit
 & Limited Kraus data; unnormalized finite branches
 & Ordinary shadow theorem; no new finite-probability output solely from scalar extension\\[3pt]
Conditioning a branch
 & Success probability with nonzero residue
 & Leading Gram matrices; error smaller than success scale; reciprocal real trial budget\\[3pt]
Eliminating unlimited modes
 & Critical scaling, $D_0>0$
 & Exact graph; anticommutator correction; uniform remainder is not a deep-scale certificate\\[3pt]
Eliminating quadratic modes
 & $D_0>0$, or all soft whitened couplings have zero residue
 & Positive defect of rank at most $\dim\ker D_0$; unrestricted displacements may diverge\\[3pt]
Using a static spectral approximation
 & $\abs z/\lambda_{\min}(D)$ infinitesimal
 & Energy-dependent Feshbach pencil; transition scales cannot be dropped\\[3pt]
Reading a near-identity loop
 & Retain the order required by the observable
 & Regular invariant scalars start at second order; reference or normalization changes the question\\[3pt]
Discriminating a near-identity gate
 & None at finite query count
 & Optimal rate $w_m$ of valuation $2\kappa$; appreciable signal only at rare success\\[3pt]
Heralding the loop gate
 & Controlled calibrated unitary and sufficiently low accepted noise
 & Success scale $2v(c-1)$; coherent precision at scale $v(c-1)$\\
\end{longtable}

\section{Boundaries and counterexamples that must remain visible}\label{qgs:sec:boundaries}
\TA\ (Source 03, with additions.)

\paragraph{Infinitesimals are not error bars without a model.}
The statement $v(R)\ge3\alpha$ says where the first possible residual coefficient lies. It does not by itself imply a numerical estimate $\norm{R(\varepsilon)}\le C\varepsilon^3$ for a family of ordinary functions: there may be no chosen realization $t^\alpha\mapsto\varepsilon$, no convergence theorem, and no uniform bound on the ordinary coefficients. The exact three-level example has such an ordinary realization because its closed algebraic formula is explicit; \cref{qgs:ops:prop:puiseux} gives one for every finite rational-power expression. A general arbitrary-support Hahn series does not acquire one automatically. Similarly, $v(p)>0$ does not say how many ordinary experimental trials are needed until a success occurs (\cref{qgs:sec:copies}).

\paragraph{Pointwise smallness does not control derivatives.}
The gauge theorems require global smooth-coefficient supports. Without such control, standard part need not commute with differentiation. The familiar ordinary family $f_\varepsilon(x)=\varepsilon\sin(x/\varepsilon)$ converges uniformly to zero while $f_\varepsilon'(0)=1$. This is one of the derivative/reduction warnings already discussed in the physics report \cite{repoPhysics}. It illustrates why one cannot replace our coefficient space by arbitrary pointwise infinitesimal perturbations and expect the curvature or elliptic arguments to survive unchanged. In our coefficient space, differentiation behaves well for an algebraic reason: it acts on each ordinary coefficient without changing its exponent. This is an explicit admissibility condition, not a conclusion that every surreal-valued field automatically has the required regularity.

\paragraph{Unlimited phases and proper-class spacetime.}
The dynamical theorems work only after the relevant generator has become finite by rescaling and scalar-phase removal. They say nothing about the uniqueness or physical adequacy of a global phase map at unlimited imaginary arguments. The collection analyzes distinct global exponential proposals, their periods, and their incompatibilities with particular derivations \cite{repoPhysics,repoQuaternion}; we have deliberately avoided using an unjustified global phase rule. Nor do we replace an ordinary spacetime manifold by the proper class $\No^4$. Size, topology, differential equations, causal structure, and an observable interpretation would all need to be specified for that proposal. The no-path phenomena of the full fine topology and the lack of ordinary completeness arguments are not repaired merely by a field embedding.

\paragraph{Gravity and singularity claims.}
The endpoint equation $r^6\mathcal K=48m^2$, $m\ne0$, for the Schwarzschild Kretschmann scalar $\mathcal K$ has no field-valued solution at $r=0$ in characteristic zero. Evaluating $\mathcal K$ at a positive infinitesimal radius does not change that algebraic fact. The physics report derives and analyzes this obstruction in detail \cite{repoPhysics}; speculative surreal-physics proposals, such as \cite{Nieto2016}, should not be read as supplying the missing regular endpoint or the missing dynamics. The positive results here are consistent with that negative lesson. An infinitesimal can encode the scale of an effect; it cannot force an inconsistent endpoint equation, a nonzero harmonic obstruction, or a forbidden topological charge to vanish. What can change is the question: a rescaled quantum experiment may have a nontrivial residue even when an unrescaled one does not, and an unobstructed ordinary gauge background may have a rich hierarchy of coefficient-valued deformations. Sources 04, 07 and 08 defer to the same obstruction; none claims singularity resolution or new spacetime dynamics.

\section{Open questions}\label{qgs:sec:open}
\TA\ These are research questions raised by the sources, not claims to identify or solve previously named open problems; none of them is claimed resolved.
\begin{enumerate}
\item \emph{Hierarchical elliptic operators} (03). How broadly does the support-controlled Kuranishi construction extend when both the background geometry and the elliptic operator have hierarchical perturbations? An ordinary inverse with nonzero residue is easy to treat by another formal lift, but an operator whose small eigenvalues occur at several ranks needs a combined spectral/deformation analysis. The finite matrix theorems are a model for this issue, not its infinite-dimensional solution; the collection's infinite-dimensional spectral report treats compact self-adjoint operators under noncommuting Hahn perturbations, which is relevant input.
\item \emph{Faithful realizations} (03). Which classes of ordinary asymptotic families admit faithful realizations of the Hahn certificates, with uniform error bounds and admissible rescalings? \Cref{qgs:ops:prop:puiseux,qgs:ops:prop:rank-two,qgs:cor:realrealization} answer this for finite rational-power and rank-two monomial expressions and for the elimination theorem only.
\item \emph{Comparison with deformation theory} (03). Compare the exact statements here with general formal deformation theory and valued algebraic geometry, to identify which support/valuation formulations are already implicit or explicit there.
\item \emph{Infinite resources} (03). Extending the finite operational theorem to a specified infinite-resource or infinite-dimensional setting requires a choice of tensor products, completeness, states, and probability measures. The proofs isolate precisely where finiteness was used; they do not license deleting that hypothesis.
\item \emph{$\SU(d)$ holonomies} (04). An exact arbitrary-rank resource law for general $\SU(d)$ holonomies. The two-eigenvalue geometry reduces everything to one scalar $w$; in higher dimension several eigenvalue separations and numerical-range faces may enter. A useful result would identify a finite collection of gauge-invariant valuations controlling optimal postselection, together with attainment certificates.
\item \emph{Noise-constrained discrimination} (04). Prescribe valuation bounds on control errors or restrict available effects, and ask for the achievable region. \Cref{qgs:thm:cost,qgs:hol:thm:precision} provide lower-level ingredients but not a complete solution.
\item \emph{Growing query budgets} (04). A many-plaquette probe with a dimension or query budget that grows along an ordinary real asymptotic family; this is different from using an infinite surreal integer as a circuit size, and needs uniform estimates for a family of finite ordinary circuits.
\item \emph{Relational observables} (04). A systematic invariant-theoretic treatment of pure gauge scalar observables versus relational matter-probe observables for networks with several loops and reference systems.
\item \emph{Schrieffer--Wolff extensions} (04). Extending the rigorous ordinary Schrieffer--Wolff theory must track spectral-gap scales, locality, and time of validity; this is not accomplished by the finite holonomy theorem, and \cref{qgs:sec:elimination,qgs:sec:block} treat finite matrices only.
\item \emph{Indefinite pencils} (08). A classification analogous to \eqref{qgs:ops:eq:classification-set} for nonpositive or indefinite pencils, where the Gram positivity mechanism fails and divergent cancellations are possible.
\item \emph{Higher-dimensional heralded unitaries} (08). For any unitary $U$ over $K$, the scalar-effect condition $(I-U)^*(I-U)\propto I$ is equivalent to $U+U^*=aI$ for a real scalar $a$; multiplying by $U$ gives $U^2-aU+I=0$, so the spectrum lies in the at-most-two-element set of conjugate unit-circle roots with common real part $a/2$, and conversely such a spectrum gives $U+U^*=aI$ by unitary diagonalization (a unitary and its adjoint preserve the same eigenspace complements). The $\SU(2)$ case is special because the condition holds for every element. This elementary criterion is not presented as an unsolved classification problem; the open question is physical realization by a given gate set.
\item \emph{Coherent recombination} (08). A finite multi-branch invariance theory that includes coherent recombination of several rare pathways and gives a minimal set of leading amplitude certificates.
\item \emph{Evolution} (this merge). The evolution part of the vector-and-tensor-fields report's open item on gauge-fixed nonlinear systems (\cref{qgs:sec:vtfrel}).
\end{enumerate}

\section{Conclusions}\label{qgs:sec:conclusion}
\TA\ The strongest positive conclusion is an exact method, not a new law of nature. A nonzero rare quantum branch has a well-defined leading normalized state; an independently scaled high-energy sector has an exact dressed finite Hamiltonian under explicit positivity assumptions, and an isolated residue cluster an exact support-certified block whose dynamics on the inverse splitting scale is an ordinary quantum evolution; the possible shadows of a soft-mode Schur complement are exactly classified; the resolution of a small non-Abelian holonomy by finite quantum protocols has an exact attained frontier; and a near-flat non-Abelian loop or connection has a gauge-invariant initial curvature whose squared norm controls its positive action.

The strongest negative conclusions are equally informative. All perturbative power data in one chosen scale can fail to determine a rare conditioned state, and identical shadows can conceal any ordinary channel. Any appreciable amplification of initially infinitesimal distinguishability by the same postselection operation has a correspondingly small success weight, and no ordinary finite number of queries discriminates an infinitesimally different gate perfectly. A nonzero leading quadratic obstruction and a fixed-bundle charge cannot be changed by smaller scales. A naive quaternionic composite cannot accommodate independent local preparation and effect statistics at the prescribed dimension. Positivity and compatibility still enforce the ordinary Bell bound, and characteristic-zero trace algebra still forbids exact finite-dimensional complex canonical commutation relations.

These results support a concrete research direction: use surreal or Hahn coefficients to preserve scale information, and use explicit reduction and resource theorems to decide which parts can influence an ordinary model --- keep ordinary spacetime and physical postulates explicit, choose a set-sized coefficient workspace, prove support and positivity, and state which normalization precedes reduction. Every recovery comes with a hypothesis that must not be suppressed: a rare outcome, an unrestricted soft displacement, an energy scale, a reference, or a singular normalization. The results do not support the inference that replacing scalars, by itself, repairs singularities, constructs quantum fields, or reveals empirically new degrees of freedom, and their potential novelty remains to be assessed at the level of the precise theorems rather than their terminology.

\appendix
\section{Exact verification and reproducibility}\label{qgs:app:verification}
\TE\ Four programs accompany the report, one from each source, kept byte-identical under the file prefixes of \cref{qgs:tab:sources}. Each performs exact symbolic or rational calculations with SymPy, uses no floating-point acceptance tolerance, and raises an exception if a check fails. None is an implementation of the surreal field, a Hahn support construction, or a Lean proof.

\begin{longtable}{@{}L{0.30\textwidth}L{0.12\textwidth}L{0.50\textwidth}@{}}
\caption{The four verification programs.}\label{qgs:tab:verification}\\
\toprule
Program (in \texttt{code/}) & Checks & What is checked\\
\midrule
\endfirsthead
\toprule
Program (in \texttt{code/}) & Checks & What is checked\\
\midrule
\endhead
\bottomrule
\endfoot
\nolinkurl{07-scales-in-physics-verify_examples.py} & 43 named & Quaternion multiplication, conjugation, norm and complex representation; the unit-circle form \eqref{qgs:eq:exactplaquette} of the commutator plaquette; the Bell sum of squares and a saturating state; the scalar effective eigenvalue expansion; the three-level characteristic polynomial and symmetric eigenvalues of \cref{qgs:prop:virtual}; the sharp diagonal postselection example; the rare pure branch of \cref{qgs:ex:rare3}; the Riccati coefficient recursion, metric self-adjointness of the graph generator and the anticommutator correction for a noncommuting $2\times2$ example; the quaternionic dimension gap.\\[3pt]
\nolinkurl{03-observable-limits-verification.py} & 40 & The quaternionic representation and $[\qi,\qj]=2\qk$; the tagged example and the conditional qubit (\cref{qgs:obs:ex:qubit}); leading Gram products; the three-level characteristic equation, the expansion \eqref{qgs:obs:eq:lambda} through degree eight and the transition probability \eqref{qgs:obs:eq:rabi}; a noncommuting Riccati recursion, graph isometry, unitary normalization and effective block through degree six; a nonconstant noncommutative Chern--Weil transgression; the torus curvature, zero charge density, action coefficient $8$ and harmonic self-dual norm $4$; a finite-dimensional toy Kuranishi recursion through degree eleven (\cref{qgs:app:recursions}).\\[3pt]
\nolinkurl{04-gauge-holonomy-verify.py} & 148 & Generic commutator identities in six real coordinates; multiplicativity of the complex representation; Chebyshev norm identities and first and second derivatives at $1$ for $m=1,\ldots,16$; exact filter normalization, contraction and overlaps for sample $(w,p)$; the nearly aligned example \cref{qgs:hol:ex:aligned} and leading query coefficients for $m\le8$; conditioning sharpness examples; the quaternionic dimension identity; simulation unitarity entries.\\[3pt]
\nolinkurl{08-quantum-operations-verify.py} & 1,861 & Leading Gram states and Kraus mixing, including \cref{qgs:ops:ex:mixed} (54); rational diagonal postselection inequalities with unequal success probabilities and exact saturation (756); the two-qubit filter and partial-transpose characteristic polynomial (25); soft-mode realizations, rank-restricted defects and full block congruences, with one ordinary and several soft modes at different powers (960); the Feshbach determinant, transition regimes and scalar positivity threshold (22); quaternion multiplication, loop trace, port completeness and the state-independent dark effect (44). Deterministic seed 20260923.\\
\end{longtable}

\paragraph{Recorded runs.} The delivered records are \nolinkurl{data/07-scales-in-physics-verification_results.txt}, \nolinkurl{data/03-observable-limits-verification_results.txt}, \nolinkurl{data/04-gauge-holonomy-verification.json} with \texttt{.txt}, and \nolinkurl{data/08-quantum-operations-verification.json}; each source ran Python 3.13.5 and SymPy 1.14.0. For this report all four were rerun on copies with Python 3.14.4 and SymPy 1.14.0: 43, 40, 148 and 1,861 checks passed, exit status zero, and the outputs agree with the delivered records except for the recorded Python version (and line endings of the files that source 04's and 08's programs rewrite).

\paragraph{What the checks do not establish.} A successful run verifies finite identities and examples only. It does not establish real closure of a Hahn field, the Neumann support lemma, an arbitrary-rank implicit theorem, a universally quantified spectral, operational, discrimination or classification theorem, the optimality of adaptive protocols, the elliptic Hodge theorem, Chern--Weil integrality, the analytic invariant-germ theorem, physical realizability, or novelty. Source 03's toy recursion is labelled as such and is not a computation of a Green operator on a four-manifold; checking a transgression identity on polynomial local forms does not replace the general proof or Stokes' theorem. No numerical approximation is used as evidence for an exact identity, and no computer-algebra output is represented as a Lean proof. The README records how to rerun each program without overwriting delivered files.

\section{Recursions and formulas in implementable form}\label{qgs:app:recursions}

\subsection{Matrix graph coefficients (source 03)}
For a single bookkeeping variable $z$, suppose $P_{11}=zP_{11}^{(1)}$, $P_{12}=zP_{12}^{(1)}$, $P_{22}=zP_{22}^{(1)}$, and write $X=\sum_{n\ge1}z^nX_n$. The recurrence used by source 03's check is
\begin{align}
 X_1&=-\Gap^{-1}P_{12}^{(1)*},\\
 X_n&=-\Gap^{-1}P_{22}^{(1)}X_{n-1}+\Gap^{-1}X_{n-1}P_{11}^{(1)}
       +\Gap^{-1}\sum_{\substack{r+s=n-1\\r,s\ge1}}X_rP_{12}^{(1)}X_s
       \quad(n\ge2).
\end{align}
For general multivariable inputs, replace coefficient indices by multi-indices and use total degree. This is the construction of \cref{qgs:obs:thm:block}. The bookkeeping variable is not presumed to be a cofinal monomial of the final exponent group; it is used to generate a formal expression before Hahn evaluation. To compute the orthonormal block to degree $N$, truncate each multiplication to that degree and use $(I+X^*X)^{-1/2}=\sum_{j\ge0}\binom{-1/2}{j}(X^*X)^j$. Since $X^*X$ begins in degree two, at most $\lfloor N/2\rfloor$ nonconstant powers are relevant. Hermiticity, unitarity, and the Riccati residual can then be checked as exact polynomial coefficient identities.

\subsection{A finite-dimensional obstruction toy model (source 03)}
The toy version of the nonlinear splitting is not an $\SU(2)$ background. Let $V=\R^2$ with coordinates $(x,y)$, let $W=\R^2$, and set
\[
 \Lop(x,y)=(y,0),\quad \Green(s,t)=(0,s),\quad
 \pi_1(x,y)=(x,0),\quad \Pi(s,t)=(0,t).
\]
These satisfy the same linear splitting identities as \eqref{qgs:obs:eq:Q}. Take the quadratic nonlinearity $\mathcal N(x,y)=(x^2+xy,\ x^2+y^2)$. \TE\ The fixed-point equation $a=(u,0)-\Green\mathcal N(a)$ has the exact solution
\begin{equation}
 x=u,\qquad y=-\frac{u^2}{1+u},\qquad
 \Lop a+\mathcal N(a)=\left(0,u^2+\frac{u^4}{(1+u)^2}\right).
\end{equation}
The degree-one tangent is $u$, the correction begins in degree two, and the residual lies entirely in the obstruction coordinate. Its nonzero quadratic coefficient cannot be cancelled by later terms. The example tests the recurrence, factor order, and projection mechanism while making no claim to test elliptic regularity or a geometric moduli problem.

\subsection{Exact formulas for holonomy discrimination (source 04)}
For a unit quaternion $q=s+x\qi+y\qj+z\qk$,
\[
 U=\chi(q)=\begin{pmatrix}s+ix&y+iz\\-y+iz&s-ix\end{pmatrix},
 \qquad w=1-s,
 \qquad \TD_1=\sqrt{2w-w^2}.
\]
\TE\ For ordinary $m=1,2,3,4$, the exact conclusive-event polynomials are
\begin{align*}
 w_1&=w, &
 w_2&=4w-2w^2,\\
 w_3&=9w-12w^2+4w^3, &
 w_4&=16w-40w^2+32w^3-8w^4.
\end{align*}
In the partial-separation regime $p>w_m$, the successful filter in the symmetric two-state basis has the form $M_{\mathrm s}=\diag\bigl(\sqrt{(2p-w_m)/(2-w_m)},\ 1\bigr)$; in the orthogonal-output regime $p\le w_m$ it is $M_{\mathrm s}=\diag\bigl(\sqrt{p/(2-w_m)},\sqrt{p/w_m}\bigr)$. Every displayed diagonal entry is finite and at most one. The normalization of the conditional state divides by $\sqrt p$, but the \emph{physical filter itself} contains no unbounded matrix entry. \TA\ This distinction is essential: the large conditional amplification is paid for by success probability, not implemented by a forbidden unbounded quantum operation.

For a commutator it is unnecessary to compute $c$ first when only $w$ and the optimal scalar resource are needed. One may calculate the three coordinates of $u_1\times u_2$, its squared norm, and the two scalar denominators in \eqref{qgs:hol:eq:comm-w}, then apply the Chebyshev recurrence and the piecewise formula \eqref{qgs:hol:eq:frontier}. For exact leading-scale certification, retain enough coefficient data to establish \eqref{qgs:hol:eq:precisionhyp}. A leading-zero test alone is not a certificate that the full cross product vanishes.

\section{Proof dependencies}\label{qgs:app:dependencies}
\TA\ Real-closed field foundations give finite Hermitian diagonalization and positive square roots (\cref{qgs:prop:spectral}); positive Gram algebra then gives limited normalized matrices and their residues (\cref{qgs:lem:squares,qgs:lem:finitequantum}). From there the principal chains are:
\begin{itemize}
\item Gram valuation $\to$ leading Kraus state (\cref{qgs:thm:rare}) $\to$ amplitude precision (\cref{qgs:ops:cor:amplitude-precision}) and dark-port precision (\cref{qgs:ops:prop:gate-precision}).
\item Positive trace-preserving contraction and the failure flag $\to$ the cost theorem (\cref{qgs:thm:cost}) $\to$ valuation budget, calibration cost, and the classical improvement of \texttt{fsp:thm:stability}.
\item Formal implicit evaluation and the support lemma $\to$ invariant graphs (\cref{qgs:thm:elimination,qgs:obs:thm:block}) $\to$ dressed and inverse-scale dynamics, hierarchy, virtual selection.
\item Positive square completion $\to$ bounded whitening $\to$ defect formula and realization (\cref{qgs:ops:thm:defect,qgs:ops:thm:classification}) $\to$ displacement obstruction and energy window.
\item Finite positive linear algebra $\to$ overlap triangle and instrument dilation $\to$ query overlap bound (\cref{qgs:hol:thm:query}); with the explicit separation filter $\to$ exact frontier (\cref{qgs:hol:thm:frontier}); with the exact Wilson defect $\to$ the valuation law (\cref{qgs:hol:thm:valuation}). None of these steps requires a surreal derivative, a Hahn-summed exponential, a compact unitary group, or a Banach-space theorem.
\item Hamilton multiplication $\to$ commutator norm (\cref{qgs:ops:thm:commutator}) $\to$ heralded unitary (\cref{qgs:ops:thm:dark}) and Wilson action (\cref{qgs:thm:wilson}).
\item Coefficientwise transgression and Stokes $\to$ charge rigidity $\to$ action gap; Hodge--Green splitting and the support lemma $\to$ Kuranishi map $\to$ obstruction and action floor.
\end{itemize}
The class-function result uses only conjugacy of equal-norm imaginary vectors and even polynomial restriction; the coefficient certificate uses bilinearity of the cross product; the quaternionic composite obstruction uses only finite dimension, positivity spans, and linear independence or bilinearity. The analytic invariant-germ theorem, the implicit-function constructions, the finite-angle exponential and the Kuranishi map are the places where formal Hahn evaluation is used. The real-specialization propositions impose further function-theoretic hypotheses and are not used to prove the algebraic classification.

\section{Source and novelty audit}\label{qgs:app:audit}

\subsection{What each source inspected}
\TA\ All four literature searches and repository inspections were performed on 23 September 2026. Primary sources were preferred, and each source's bibliography is merged into the one below.
\begin{itemize}
\item \emph{Source 07} (pin \texttt{c0a36ee}) inspected the repository README, the physics report's README, directory metadata, and selected code-index searches. A local clone was unavailable; no build was performed; the Lean source was not audited; the results are not added to the verified inventory.
\item \emph{Source 03} (pin \texttt{934810a}) inspected the root README, the opening 190 lines of the formalization ledger, and the opening portions of the physics and surquaternion guides, as recorded in its delivered audit file \nolinkurl{03-observable-limits-source_audit.md}. It read printed pages 25--26 of Itoh's paper, including Proposition 2.4 and the beginning of Section 3, in rendered images; it identified Atiyah--Hitchin--Singer through Itoh's reference list without a full fresh reading; it did not use Bordemann--Waldmann for a new GNS theorem; exploratory searches returned irrelevant results and a recent quaternionic preprint, none of which was used as a dependency or for a priority claim.
\item \emph{Source 04} (pin \texttt{934810a}) inspected the physics, surquaternion and spectral guides and the public tree through a connected repository reader, since a direct clone was unavailable. It checked Chefles--Barnett (including its success-probability bound) and Ac\'in (including the finite-use spectral-angle threshold) directly, and consulted primary sources for quaternionic simulation and composition, gauge theory, holonomic control, and postselection. Its incomplete code search for ``postselection'' was not used as an absence proof (and was a false negative, \cref{qgs:sec:stale}).
\item \emph{Source 08} (pin \texttt{c0a36ee}) inspected the README (opening), the documentation catalogue and reading routes, the physics, finite-probability and surquaternion guides, and lines 1--260 of the spectral article, as recorded in its delivered audit file \nolinkurl{08-quantum-operations-RESEARCH_AUDIT.md}. It inspected Bordemann--Waldmann, Vidal, Dusson--Sigal--Stamm and Gallier in relevant text, Watrous's author manuscript for Corollary 3.40, and only publisher records or abstracts for Anderson--Trapp, Wilson, Torres--Salazar-Serrano and Barnum--Graydon--Wilce; Gonshor and Berarducci--Mantova were consulted as foundational references without a full re-audit. Its Schur-term code search returned no matches and was not used as an absence proof (a false negative, \cref{qgs:sec:stale}). It did not upgrade a guide's report of a successful build into independently checked formal verification, nor assume that pending manuscript claims are proved because they are in the repository.
\end{itemize}

\begin{longtable}{@{}L{0.22\textwidth}L{0.36\textwidth}L{0.34\textwidth}@{}}
\caption{Source 08's targeted comparison with the collection, not an assertion of exhaustive novelty.}\label{qgs:ops:tab:overlap}\\
\toprule
Existing report & Already available & Focus developed here\\
\midrule
\endfirsthead
\toprule
Existing report & Already available & Focus developed here\\
\midrule
\endhead
\bottomrule
\endfoot
Surreal scalars and spacetime \cite{repoPhysics}
 & Finite quantum shadows; their conditioning boundary; limits of scalar replacement
 & Leading Kraus/Gram states; quantified postselection cost\\[3pt]
Finite surreal probability \cite{repoProbability}
 & Classical rare-event conditional shadows and precision bounds
 & Mixed quantum branches and entanglement filtering\\[3pt]
Surcomplex spectral theory \cite{repoSpectral}
 & Finite Hermitian algebra, singular values, scales and Schur-type effective matrices
 & Complete fixed-shadow Schur-defect classification and bounded-displacement obstruction\\[3pt]
Surquaternions \cite{repoQuaternion}
 & Hamilton algebra, $\SU(2)$ realization, filtered commutators and local analysis
 & Exact loop trace loss, heralded unitary recovery, and noise-scale tests\\
\end{longtable}

\subsection{What is, and is not, presented as a contribution}
\TA\ Established background: the existence of surreal normal forms and real closure; finite spectral theory; ordinary quantum channels, trace-distance contraction, Helstrom, pure-state concentration; Schrieffer--Wolff reduction, the generalized Schur formula, the Feshbach kernel correspondence; state separation, unambiguous discrimination and unitary discrimination geometry; quaternionic complex simulation and the Hamilton representation of $\SU(2)$; Wilson loops; the Tsirelson bound; Kuranishi's construction and Chern--Weil theory. The collection's finite quantum shadow and its conditioning boundary, the classical conditional skeleton, the classical conditioning stability bound, the exact effective Schur matrix and the leading group commutator are prior collection results.

Candidate contributions of the sources, none certified as first in the literature: the rare-branch initial-form package with its all-one-scale-power obstruction and universal rare embedding (03, 07, 08); the sharp symmetric cost theorem as an explicit valuation budget (04, 07, 08); the independently scaled invariant-graph and remainder theorem (07); the exact positive-support Hahn block diagonalization with support certificates, inverse-scale dynamics, hierarchy and Gibbs reduction (03); the complete fixed-shadow Schur-defect classification with its displacement consequence (08); the exact arbitrary-rank gauge, query and postselection synthesis with its alignment certificate (04); the combined leading-Gram, commutator and conditioned-unitary application with coherent and incoherent precision thresholds (08); the positive initial-curvature formula (07); and the arbitrary-support charge, action, Kuranishi, obstruction and action-floor package (03).

None of these labels rules out an earlier equivalent theorem under different terminology. In particular, the elimination theorems are closely related to matrix Riccati equations and Schur-complement methods, and the rare-event formula uses familiar positivity and leading-term ideas; the Kuranishi constructions over formal coefficient rings and recursive effective-Hamiltonian methods are not new, and a publication-level priority claim would require a more exhaustive comparison with valued linear algebra, nonstandard probability, singular perturbation, deformation functors and valued coefficient rings than these targeted reviews supply. The standard ASD deformation references beyond those cited (for example the treatments of Freed--Uhlenbeck and Donaldson--Kronheimer) were not consulted by source 03, which disclaims invention. The report nevertheless makes falsifiable mathematical assertions: each theorem has stated hypotheses and a proof, and its examples have exact outputs. Independent review can therefore strengthen, correct, or reject a specific claim without having to interpret an undefined claim of a ``surreal breakthrough.'' No major named open problem in physics is claimed to be solved.

\section{Register of non-claims}\label{qgs:app:nonclaims}
\TA\ Every limitation, disclaimer and priority caveat stated by any source is kept. Most are stated at their point of use above; this register lists them all, by source, with that location. Counts: 38 (source 03), 25 (source 04), 31 (source 07), 29 (source 08).

\paragraph{Source 03 (38).}
(1) Scalar replacement alone gives no new dynamics, singularity resolution or observable change (\cref{qgs:sec:intro}). (2) The mechanisms are classical; the contribution is the valuation-resolved arbitrary-rank formulation; priority is not established; nothing is Lean-verified; no experimental departure is claimed (\cref{qgs:app:audit}). (3) The repository audit was targeted, not every theorem or manuscript (\cref{qgs:app:audit}). (4) Formal status is not inferred from a report's presence; no new Lean check; the program tests selected identities (\cref{qgs:app:verification}). (5) The literature search shows antecedents, not absence; formal Kuranishi constructions and recursive effective Hamiltonians are not new (\cref{qgs:app:audit}). (6) No proper-class sums; the surreal reading only after a fixed embedding (\cref{qgs:sec:surreal}). (7) Valuation statements are not numerical error bounds, and $v(p)>0$ says nothing about trial counts (\cref{qgs:sec:fields,qgs:sec:boundaries}). (8) No valuation-topology or Picard convergence is assumed (\cref{qgs:obs:ex:noconv}). (9) Derivatives are coefficientwise, not Berarducci--Mantova; strong summation is not fine-topology continuity (\cref{qgs:obs:sec:conventions}). (10) Reading $\st p$ as an observed probability is a modelling choice (\cref{qgs:prop:shadow}). (11) No surreal number of copies; infinite or hyperfinite products need a separate theory (\cref{qgs:sec:copies}). (12) The constant two in its postselection bound is not claimed optimal (\cref{qgs:obs:cor:postselection}). (13) No free distinguishability for finite experiments (\cref{qgs:sec:copies}). (14) The rare-event engine gives no finite-resource information from zero-readout events, no realizable apparatus with literal surreal probabilities, and needs an explicit ordinary family (\cref{qgs:obs:sec:rareapp}). (15) Not a new Schrieffer--Wolff principle; no many-body locality or size-uniform bound (\cref{qgs:sec:block}). (16) No phase for arbitrary infinite imaginary arguments; inter-block coherence with unlimited relative phase is not solved (\cref{qgs:obs:thm:dynamics}). (17) The insensitivity corollary is an independence statement, not existence of a high-frequency evolution (\cref{qgs:obs:cor:insensitive}). (18) The three-level sharper remainder is not a general improvement (\cref{qgs:obs:sec:threelevel}). (19) Finite depth is not finite computational cost; the scalar-block test may be undecidable for a notation; the final support is not the original monoid (\cref{qgs:obs:sec:hierarchy}). (20) The Gibbs exponential may leave the Hahn workspace; no phase transition or volume limit (\cref{qgs:obs:sec:gibbs}). (21) No equivalence of quaternionic and complex quantum mechanics in every formulation; locality and tensor structure are not preserved; Moretti--Oppio is distinct; not every quaternionic proposal is redundant (\cref{qgs:sec:quaternions}). (22) Coefficient extension on a fixed ordinary bundle, no surreal charts or class-based bundle; Chern--Weil integrality is imported (\cref{qgs:sec:gauge}). (23) Charge rigidity does not exclude instanton sectors, bundle sums, tunnelling, boundary flux or singular limits (\cref{qgs:obs:thm:chern}). (24) Elliptic regularity and Hodge splitting are assumed, not derived; no uniform norm bound (\cref{qgs:obs:ass:hodge}). (25) No new topological sectors; no background created (\cref{qgs:obs:cor:regular}). (26) No global moduli classification; not every connection is gauge-fixable; stabilizers are separate (\cref{qgs:obs:prop:slice}). (27) No claim about higher obstructions when the quadratic one vanishes (\cref{qgs:obs:thm:obstruction}). (28) Kuranishi is classical; priority is open; not a physical discovery (\cref{qgs:obs:thm:floor}). (29) Effective sublanguage only; equality and support comparison are not solved; transseries already provide related capabilities; no exclusive claim (\cref{qgs:sec:applications}). (30) No new instanton classification (\cref{qgs:sec:vtfrel}). (31) No continuum QFT, path-integral measure, renormalized Hamiltonian or regulator-independent observable (\cref{qgs:sec:applications}). (32) Pointwise smallness does not control derivatives (\cref{qgs:sec:boundaries}). (33) No global phase adequacy; no $\No^4$ spacetime (\cref{qgs:sec:boundaries}). (34) No regular Schwarzschild endpoint; Nieto's proposal is speculative (\cref{qgs:sec:boundaries}). (35) The toy recursion is not a Green operator; the checks do not prove real closedness, the support lemma, the Hodge theorem or integrality; no Lean build (\cref{qgs:app:verification}). (36) Its open questions are not claimed resolved (\cref{qgs:sec:open}). (37) Third-party papers and fonts are not included; the repository was not modified by the source. (38) Audit: Bordemann--Waldmann is not used for a new GNS theorem; Atiyah--Hitchin--Singer was not freshly re-read; irrelevant search hits and a quaternionic preprint were not used (\cref{qgs:app:audit}).

\paragraph{Source 04 (25).}
(1) Finite dimensions, ancillas, depths and queries; no proper-class Hilbert space, global surreal exponential or infinite-query experiment (\cref{qgs:sec:discrimination}). (2) Chefles--Barnett, Ac\'in, Graydon, Wilson, Zanardi--Rasetti and Cirel'son are credited, not claimed new (\cref{qgs:sec:discrimination,qgs:sec:quaternions,qgs:sec:holonomy,qgs:sec:nogo}). (3) A candidate synthesis; historical priority is not certified; independent review is appropriate (\cref{qgs:app:audit}). (4) No empirical departure, continuum Yang--Mills, ultraviolet completion, quantum gravity or proof-assistant verification (\cref{qgs:sec:intro,qgs:app:verification}). (5) Targeted repository comparison; an incomplete code search is not evidence of absence (\cref{qgs:sec:stale}). (6) The standard-part reduction is background, not a new impossibility theorem (\cref{qgs:prop:shadow}). (7) CHSH is included only to delimit (\cref{qgs:thm:bell}). (8) The Wilson identity is not a new quaternion theorem (\cref{qgs:hol:cor:comm}). (9) The regularity restriction is essential; the class-function theorem covers only polynomial functions of one $\SU(2)$ element (\cref{qgs:hol:thm:classfunctions}). (10) An infinite omnific $m$ is illegitimate in a finite circuit (\cref{qgs:hol:cor:nofinite}). (11) ``Finite'' means an ordinary integer fixed independently (\cref{qgs:hol:prop:queryscale}). (12) Not every output distance has the frontier valuation (\cref{qgs:hol:thm:valuation} and the remark after it). (13) No countably additive surreal probability; ``wait $1/p$ trials'' is not a finite operation (\cref{qgs:sec:copies}). (14) The strict precision inequality cannot be weakened uniformly (\cref{qgs:hol:thm:precision}). (15) The calibration bound is worst-case sufficient; structured noise may be milder; amplifiers need a noise model (\cref{qgs:hol:cor:calibration}). (16) The lattice action is a finite function, not a path integral; commutator part only; no thermodynamic limit, measure, confinement or mass gap; negative-valuation weights are not measurement effects (\cref{qgs:thm:wilson}). (17) Known gates and eigenbasis are assumed; no universal estimator; no reference-alignment solution (\cref{qgs:hol:sec:physics}). (18) Arbitrary pairs are not shown to be adiabatic holonomies (\cref{qgs:hol:sec:physics}). (19) Binary discrimination, not Fisher information; the weak-value debate is not adjudicated (\cref{qgs:hol:sec:physics}). (20) The simulation does not identify all doubled states nor settle composition (\cref{qgs:thm:quatsimulation}). (21) The composite no-go does not exclude Jordan-algebraic composites (\cref{qgs:hol:thm:composite}). (22) The Chebyshev branch holds only in the small-angle regime; no claim for arbitrary real parameters or growing query counts; physical relevance of the real family is separate (\cref{qgs:hol:sec:real}). (23) $\SU(d)$, noise-constrained, many-plaquette and invariant-theoretic problems are open; the Schrieffer--Wolff extension is not accomplished; no singularity or spacetime claim (\cref{qgs:sec:open,qgs:sec:boundaries}). (24) The checks do not verify adaptive optimality, the real-closed characterization, surreal constructions, physical realizability or novelty; running the program overwrites its outputs; rebuilding changes PDF metadata and checksums (\cref{qgs:app:verification}; README). (25) Authorship: ``Mathematical development and exposition: ChatGPT'' (\cref{qgs:tab:sources}). In addition, its sentence on the absence of an Archimedean query threshold is kept with the caveat of \cref{qgs:hol:cor:nofinite}.

\paragraph{Source 07 (31).}
(1) No continuum QFT, new experimental prediction or singularity resolution (title page). (2) Unrefereed and AI-assisted; the checks verify selected examples; no new Lean (title page, \cref{qgs:app:verification}). (3) ``Proposed'' is not a priority certification (\cref{qgs:sec:ledger}). (4) New physics would need derivation, readout, regime and difference; none follows (\cref{qgs:sec:intro}). (5) Scales are dimensionless (\cref{qgs:sec:intro}). (6) Nieto motivates only; Benci is different; no hyperreal identification or transfer (\cref{qgs:sec:sources}). (7) Not a discovery of virtual transitions, invariant graphs or Schur complements; quaternion doubling and composition are not novel; ``valuations in physics'' is too broad (\cref{qgs:sec:sources}). (8) Targeted review; no build; Lean not audited; not added to the verified inventory; no local clone (\cref{qgs:app:audit}). (9) $\zeta$ records an order hierarchy, not a surreal exponential; tunnelling laws need independent justification (\cref{qgs:sec:surreal}). (10) No cofinality argument (\cref{qgs:sec:strong}). (11) No computability or convergence at non-infinitesimal parameters (\cref{qgs:lem:implicit}). (12) The finite quantum algebra is established algebra re-proved (\cref{qgs:sec:quantum}). (13) No $\exp(i\Lambda)$ for unlimited $\Lambda$; no global oscillatory calculus (\cref{qgs:sec:expfin}). (14) The real-family analogy gives no embedding or summation prescription (\cref{qgs:thm:alljets}). (15) Finite idealized discrimination; not about detector saturation or noise (\cref{qgs:thm:cost}). (16) $p^{-1}$ is not a laboratory protocol; specialize a validated family (\cref{qgs:sec:copies}). (17) The uniform $O(\tau^2)$ is not a deep-scale certificate (\cref{qgs:rem:uniform}). (18) The real-parameter corollary does not show a continuum theory produces the scales (\cref{qgs:cor:realrealization}). (19) No phase for the high block; not generic initial states; time dependence needs an adiabatic theorem (\cref{qgs:sec:elimination}). (20) The virtual selection mechanism is well known (\cref{qgs:sec:virtual}). (21) No instanton action, prefactor or resummation; no uniquely surreal predictions (\cref{qgs:sec:virtual}). (22) Quaternion facts are classical; the simulation embeds a restricted family; the composite no-go does not exclude larger composites or other postulates (\cref{qgs:sec:quaternions}). (23) Positive weights are essential; finite lattice; no mass gap, confinement or continuum measure; not a continuum limit (\cref{qgs:thm:wilson}). (24) The Bell bound is classical; setting-dependent postselection is not covered; the CCR statement is not about quaternionic traces (\cref{qgs:sec:nogo}). (25) No infinite-dimensional Hilbert space, spectral measure, countably additive probability or QFT; no interchange with continuum or infinite-volume limits; no Borel prescription; fine-topology restrictions stand; the gravitational negative analysis is unchanged (\cref{qgs:sec:stop}). (26) The Wilson theorem gives no convergence with growing lattices; benchmarks are not forecasts (\cref{qgs:sec:applications}). (27) The mechanisms are not exclusive to $\No$; no principle selects an embedding (\cref{qgs:sec:realization}). (28) Formalization targets are a dependency assessment; arbitrary Hahn series are not effective (\cref{qgs:sec:formal}). (29) Verification scope (\cref{qgs:app:verification}). (30) No earlier equivalent is ruled out; no named open problem is solved (\cref{qgs:app:audit}). (31) Script variable names need not match the notation (\cref{qgs:sec:notation}).

\paragraph{Source 08 (29).}
(1) No new empirical law, ultraviolet regularization or solution of a named open problem (\cref{qgs:sec:intro,qgs:sec:open}). (2) An AI-assisted draft; candidate contributions are not first-in-literature; the 1,861 checks are not proof or formalization (\cref{qgs:app:verification,qgs:app:audit}). (3) No Hilbert completion, path integral, thermodynamic limit, continuum QFT or infinite observation is smuggled into a change of coefficient field (\cref{qgs:sec:instruments,qgs:sec:stop}). (4) Targeted inspection; build not rerun; the Schur code search is not evidence of absence (\cref{qgs:sec:stale,qgs:app:audit}). (5) Proofs are independent of the priority answer (\cref{qgs:app:audit}). (6) Theorem, model interpretation and proposal are separate levels; $\st$ is a readout rule, not a law (\cref{qgs:conv:tiers,qgs:obs:sec:reductions}). (7) Units are fixed; $t^{-\gamma}$ is not an experiment size (\cref{qgs:sec:intro}). (8) No canonical embedding; no class-sized space; the $\omega$-map, a derivation, a time derivative and evolution are not identified (\cref{qgs:sec:surreal}). (9) The Kraus form is the definition; no infinite operator-algebra theorem (\cref{qgs:sec:instruments}). (10) The mixed-output example is a Gram-output example (\cref{qgs:ops:ex:mixed}). (11) No uniform square-root Lipschitz bound (\cref{qgs:cor:initialstability}). (12) The rare embedding is not execution of a channel without implementing it (\cref{qgs:ops:prop:rare-embedding}). (13) Distinguishability, not detector performance (\cref{qgs:thm:cost}). (14) The precision contracts are sufficient, not necessary, and not interchangeable (\cref{qgs:cor:budget}). (15) $1/p$ is a trial count, not a gate-complexity bound; no infinite path measure (\cref{qgs:sec:copies}). (16) No unrestricted LOCC optimality; growing copy families are excluded; no new effect (\cref{qgs:ops:sec:entanglement}). (17) The Schur classification excludes only the positive soft-mode mechanism (\cref{qgs:ops:cor:rank-cost}). (18) No Gaussian or Hahn integration; no small-displacement extrapolation (\cref{qgs:ops:sec:quadratic-physics}). (19) Nothing is asserted near eigenvalues of $D$ (\cref{qgs:ops:thm:window}). (20) Positivity, truncation and a displacement regime are required; gauge zero modes, Minkowski signature, non-Hermitian and unbounded operators are excluded; an ultraviolet cutoff is not removable (\cref{qgs:ops:sec:quadratic-physics}). (21) Complex amplitudes are not replaced; no new tensor product; simultaneous conjugation needs based loops (\cref{qgs:sec:quaternions,qgs:sec:holonomy}). (22) The invariant-germ theorem covers only one analytic loop; it is not a ban on detecting small rotations (\cref{qgs:ops:thm:invariant}). (23) The controlled gate and path phase are hypotheses; not a deterministic flip; the strict inequality is only sufficient; the background result is not a verdict on postselection (\cref{qgs:sec:dark}). (24) Hahn series need not define functions; the rank-two realization is a choice; no infinite-sum evaluation; general workspaces need not be realizable (\cref{qgs:sec:realization}). (25) Not an empirical argument for surreal scalars; set-sized workspaces suffice (\cref{qgs:sec:realization}). (26) The decision table is not a set of universal apparatus claims; conditional unitarity is not a metrological advantage or gate synthesis; the open extensions are not named problems; the scalar-effect criterion is not an unsolved classification (\cref{qgs:sec:applications,qgs:sec:open}). (27) Limits: finite dimensions and protocols; real closed field; limited densities and blocks; positive-definite eliminated block before reduction; analytic one-loop invariance; no new measurement law, literal infinitesimal, continuum field, infinite-volume state or renormalization; scalar extension does not replace a model, observable, preparation and resource analysis (\cref{qgs:sec:schur,qgs:sec:dark,qgs:sec:realization}). (28) Verification is narrower than the proofs; no Lean (\cref{qgs:app:verification}). (29) Audit: classical background list; filtering is an elementary extension; no proof of absence; guide build reports are not upgraded; Anderson--Trapp and others from abstracts only (\cref{qgs:app:audit}).

\clearpage
\begin{thebibliography}{99}
\small\raggedright

\bibitem{repo}
V.~Reshetnikov and contributors, \emph{Surreal: Lean formalization and research reports}, repository \repo, README and documentation catalogue. The sources inspected it at commits \texttt{934810abdde4c6d74c08777b069445de43602934} (sources 03, 04) and \texttt{c0a36eebe5cf30aa72abfe4ef8ff9e46f63c3c28} (sources 07, 08). \url{https://github.com/VladimirReshetnikov/Surreal}.

\bibitem{repoLedger}
\emph{Formalization coverage and dependency ledger} and \emph{normal-form bridge}, \texttt{docs/FORMALIZATION.md} and \texttt{docs/NORMAL\_FORM\_BRIDGE.md}, in \cite{repo}.

\bibitem{repoPhysics}
\emph{Surreal Scalars in Theoretical Physics: Black-Hole Singularities, Asymptotic Structure, and the Limits of Replacing the Scalar Field}, merged research report, \texttt{docs/physics/surreal-scalars-and-spacetime/}, in \cite{repo}. Labels cited: \texttt{phys:thm:spectral}, \texttt{phys:thm:quantumshadow}, \texttt{phys:ex:rare}, \texttt{phys:ex:amplify}, \texttt{phys:sub:dynamics}, \texttt{phys:conv:tiers}, \texttt{phys:sub:stageD}, \texttt{phys:app:nonclaims}.

\bibitem{repoProbability}
\emph{Surreal Probability and Log-Odds: A Multiscale Theory of Belief, Information, and Infinite Sampling}, research report, \texttt{docs/surreal/finite-surreal-probability/}, in \cite{repo}. Labels cited: \texttt{fsp:thm:skeleton}, \texttt{fsp:ex:shadowloss}, \texttt{fsp:thm:stability}, \texttt{fsp:eq:precision}, \texttt{fsp:thm:multisoft}.

\bibitem{repoSpectral}
\emph{Surcomplex Spectral Theory}, research report, \texttt{docs/surcomplex/spectral-theory/}, in \cite{repo}. Labels cited: \texttt{prop:effective}, \texttt{ex:secondorder}, \texttt{thm:subspace}.

\bibitem{repoQuaternion}
\emph{Surquaternions: Algebra, Infinitesimal Geometry, and Support-Controlled Analysis}, research report, \texttt{docs/surquaternions/surquaternions/}, in \cite{repo}. Labels cited: \texttt{squat:prop:bch}, \texttt{squat:eq:groupcomm}, \texttt{squat:thm:graded}, \texttt{squat:thm:spectral}.

\bibitem{repoVTF}
\emph{Vector and Tensor Fields over the Surreal Numbers}, merged research report, \texttt{docs/surreal/vector-and-tensor-fields/}, in \cite{repo}. Labels cited: \texttt{vtf:subsec:open}, \texttt{vtf:mink:thm:energyscale}.

\bibitem{Gonshor}
H. Gonshor, \emph{An Introduction to the Theory of Surreal Numbers}, London Mathematical Society Lecture Note Series 110, Cambridge University Press, 1986.

\bibitem{Conway}
J. H. Conway, \emph{On Numbers and Games}, second edition, A K Peters, 2001. First edition, Academic Press, 1976.

\bibitem{vdDE}
L. van den Dries and P. Ehrlich, Fields of surreal numbers and exponentiation, \emph{Fundamenta Mathematicae} \textbf{167} (2001), 173--188. \url{https://doi.org/10.4064/fm167-2-3}.

\bibitem{BM}
A. Berarducci and V. Mantova, Surreal numbers, derivations and transseries, \emph{Journal of the European Mathematical Society} \textbf{20} (2018), 339--390. \url{https://doi.org/10.4171/JEMS/769}; arXiv:1503.00315.

\bibitem{ADH}
M. Aschenbrenner, L. van den Dries, and J. van der Hoeven, \emph{The surreal numbers as a universal $H$-field}, arXiv:1512.02267, 2015.

\bibitem{BournezGuilmant}
O. Bournez and Q. Guilmant, \emph{Surreal fields stable under exponential and logarithmic functions}, arXiv:2201.08199, 2022.

\bibitem{Nieto2016}
J. A. Nieto, \emph{Some Mathematical and Physical Remarks on Surreal Numbers}, arXiv:1611.09699, 2016.

\bibitem{Nieto2018}
J. A. Nieto, Duality, matroids, qubits, twistors, and surreal numbers, \emph{Frontiers in Physics} \textbf{6} (2018), article 106. \url{https://doi.org/10.3389/fphy.2018.00106}.

\bibitem{Benci}
V. Benci, \emph{Non-Archimedean Mathematics and the formalism of Quantum Mechanics}, arXiv:1809.04446, 2018.

\bibitem{BW}
M. Bordemann and S. Waldmann, Formal GNS construction and states in deformation quantization, \emph{Communications in Mathematical Physics} \textbf{195} (1998), 549--583. \url{https://doi.org/10.1007/s002200050402}; arXiv:q-alg/9607019.

\bibitem{BDL}
S. Bravyi, D. P. DiVincenzo, and D. Loss, Schrieffer--Wolff transformation for quantum many-body systems, \emph{Annals of Physics} \textbf{326} (2011), 2793--2826. \url{https://doi.org/10.1016/j.aop.2011.06.004}; arXiv:1105.0675.

\bibitem{DSS}
G. Dusson, I. M. Sigal, and B. Stamm, \emph{The Feshbach--Schur map and perturbation theory}, arXiv:2105.02058, 2021.

\bibitem{AT}
W. N. Anderson, Jr., and G. E. Trapp, Shorted operators. II, \emph{SIAM Journal on Applied Mathematics} \textbf{28} (1975), 60--71. \url{https://doi.org/10.1137/0128007}.

\bibitem{Gallier}
J. Gallier, \emph{The Schur Complement and Symmetric Positive Semidefinite (and Definite) Matrices}, lecture notes, University of Pennsylvania, 2019. \url{https://www.cis.upenn.edu/~jean/schur-comp.pdf}.

\bibitem{AHS}
M. F. Atiyah, N. J. Hitchin, and I. M. Singer, Self-duality in four-dimensional Riemannian geometry, \emph{Proceedings of the Royal Society of London, Series A} \textbf{362} (1978), 425--461.

\bibitem{Itoh}
M. Itoh, On the moduli space of anti-self-dual Yang--Mills connections on K\"ahler surfaces, \emph{Publications of the Research Institute for Mathematical Sciences} \textbf{19} (1983), 15--32. In particular Proposition 2.4 and the beginning of Section 3.

\bibitem{Watrous}
J. Watrous, \emph{The Theory of Quantum Information}, Cambridge University Press, 2018. In particular Chapter 3, Corollary 3.40. \url{https://cs.uwaterloo.ca/~watrous/TQI/}.

\bibitem{CheflesBarnett}
A. Chefles and S. M. Barnett, Quantum state separation, unambiguous discrimination and exact cloning, \emph{Journal of Physics A: Mathematical and General} \textbf{31} (1998), 10097--10103. \url{https://doi.org/10.1088/0305-4470/31/50/007}; arXiv:quant-ph/9808018.

\bibitem{Bagan}
E. Bagan, V. Yerokhin, A. Shehu, E. Feldman, and J. A. Bergou, \emph{A geometric approach to quantum state separation}, arXiv:1506.08241, 2015.

\bibitem{Ivanovic}
I. D. Ivanovic, How to differentiate between non-orthogonal states, \emph{Physics Letters A} \textbf{123} (1987), 257--259.

\bibitem{Dieks}
D. Dieks, Overlap and distinguishability of quantum states, \emph{Physics Letters A} \textbf{126} (1988), 303--306.

\bibitem{Peres}
A. Peres, How to differentiate between non-orthogonal states, \emph{Physics Letters A} \textbf{128} (1988), 19.

\bibitem{Acin}
A. Ac\'in, Statistical distinguishability between unitary operations, \emph{Physical Review Letters} \textbf{87} (2001), 177901. \url{https://doi.org/10.1103/PhysRevLett.87.177901}; arXiv:quant-ph/0102064.

\bibitem{DuanFengYing}
R. Duan, Y. Feng, and M. Ying, Entanglement is not necessary for perfect discrimination between unitary operations, \emph{Physical Review Letters} \textbf{98} (2007), 100503.

\bibitem{Vidal}
G. Vidal, Entanglement of pure states for a single copy, \emph{Physical Review Letters} \textbf{83} (1999), 1046--1049. \url{https://doi.org/10.1103/PhysRevLett.83.1046}; arXiv:quant-ph/9902033.

\bibitem{Graydon}
M. A. Graydon, \emph{Quaternionic quantum dynamics on complex Hilbert spaces}, arXiv:1103.3547, 2011.

\bibitem{BarnumWilce}
H. Barnum and A. Wilce, \emph{Local tomography and the Jordan structure of quantum theory}, arXiv:1202.4513, 2012.

\bibitem{BGW}
H. Barnum, M. A. Graydon, and A. Wilce, Composites and categories of Euclidean Jordan algebras, \emph{Quantum} \textbf{4} (2020), 359. \url{https://doi.org/10.22331/q-2020-11-08-359}; arXiv:1606.09331.

\bibitem{MorettiOppio}
V. Moretti and M. Oppio, \emph{Quantum theory in quaternionic Hilbert space: How Poincar\'e symmetry reduces the theory to the standard complex one}, arXiv:1709.09246, 2017. Its reduction depends on the hypotheses of that work and is not used as an unrestricted equivalence theorem here.

\bibitem{Wilson}
K. G. Wilson, Confinement of quarks, \emph{Physical Review D} \textbf{10} (1974), 2445--2459. \url{https://doi.org/10.1103/PhysRevD.10.2445}.

\bibitem{ZanardiRasetti}
P. Zanardi and M. Rasetti, Holonomic quantum computation, \emph{Physics Letters A} \textbf{264} (1999), 94--99. arXiv:quant-ph/9904011.

\bibitem{BanerjeeEtAl}
A. Banerjee, R. Jaiswal, M. Manjunath, and A. Narayan, \emph{A tropical geometric approach to exceptional points}, arXiv:2301.13485, 2023.

\bibitem{FerrieCombes}
C. Ferrie and J. Combes, Weak value amplification is suboptimal for estimation and detection, \emph{Physical Review Letters} \textbf{112} (2014), 040406. \url{https://doi.org/10.1103/PhysRevLett.112.040406}; arXiv:1307.4016.

\bibitem{JordanEtAl}
A. N. Jordan, J. Mart\'inez-Rinc\'on, and J. C. Howell, Technical advantages for weak-value amplification: when less is more, \emph{Physical Review X} \textbf{4} (2014), 011031. \url{https://doi.org/10.1103/PhysRevX.4.011031}.

\bibitem{Torres}
J. P. Torres and L. J. Salazar-Serrano, Weak value amplification: a view from quantum estimation theory that highlights what it is and what it isn't, \emph{Scientific Reports} \textbf{6} (2016), 19702. \url{https://doi.org/10.1038/srep19702}.

\bibitem{Vaidman}
L. Vaidman, \emph{Comment on ``Weak value amplification is suboptimal for estimation and detection''}, arXiv:1402.0199, 2014.

\bibitem{Cirelson}
B. S. Cirel'son, Quantum generalizations of Bell's inequality, \emph{Letters in Mathematical Physics} \textbf{4} (1980), 93--100. \url{https://doi.org/10.1007/BF00417500}.

\end{thebibliography}
\end{document}
